\documentclass[11pt,a4paper]{article}
\usepackage[T1]{fontenc}
\usepackage[utf8]{inputenc}
\usepackage{lmodern}
\usepackage[margin=25mm,headheight=15pt]{geometry}
\usepackage{amsmath,amssymb,amsthm,mathtools,mathrsfs,aliascnt}
\usepackage{microtype,booktabs,tabularx,array,enumitem}
\usepackage{xcolor,fancyhdr,titlesec,needspace}
\definecolor{Ink}{HTML}{243747}
\definecolor{Accent}{HTML}{566348}
\usepackage[colorlinks=true,linkcolor=Ink,citecolor=Accent,urlcolor=Accent]{hyperref}
\usepackage{xurl}
\usepackage[nameinlink,noabbrev]{cleveref}
\titleformat{\section}{\Large\bfseries\color{Ink}}{\thesection}{.65em}{}
\titleformat{\subsection}{\large\bfseries\color{Accent}}{\thesubsection}{.65em}{}
\pagestyle{fancy}\fancyhf{}
\fancyhead[L]{\small Infinite-dimensional Hahn spectral theory}
\fancyhead[R]{\small Spectra relative to named spaces and algebras}
\fancyfoot[C]{\thepage}
\setlength{\parskip}{.26em}
\setlength{\parindent}{1.1em}
\setlength{\emergencystretch}{3em}
\widowpenalty=10000\clubpenalty=10000
\setlist{itemsep=.15em,topsep=.3em,leftmargin=2em}
\setcounter{tocdepth}{2}
\numberwithin{equation}{section}
\newtheorem{theorem}{Theorem}[section]
\newaliascnt{proposition}{theorem}
\newtheorem{proposition}[proposition]{Proposition}
\aliascntresetthe{proposition}
\newaliascnt{lemma}{theorem}
\newtheorem{lemma}[lemma]{Lemma}
\aliascntresetthe{lemma}
\newaliascnt{corollary}{theorem}
\newtheorem{corollary}[corollary]{Corollary}
\aliascntresetthe{corollary}
\theoremstyle{definition}
\newaliascnt{definition}{theorem}
\newtheorem{definition}[definition]{Definition}
\aliascntresetthe{definition}
\newaliascnt{example}{theorem}
\newtheorem{example}[example]{Example}
\aliascntresetthe{example}
\theoremstyle{remark}
\newaliascnt{remark}{theorem}
\newtheorem{remark}[remark]{Remark}
\aliascntresetthe{remark}
\newaliascnt{warning}{theorem}
\newtheorem{warning}[warning]{Warning}
\aliascntresetthe{warning}
\crefname{theorem}{theorem}{theorems}
\Crefname{theorem}{Theorem}{Theorems}
\crefname{proposition}{proposition}{propositions}
\Crefname{proposition}{Proposition}{Propositions}
\crefname{lemma}{lemma}{lemmas}
\Crefname{lemma}{Lemma}{Lemmas}
\crefname{corollary}{corollary}{corollaries}
\Crefname{corollary}{Corollary}{Corollaries}
\crefname{definition}{definition}{definitions}
\Crefname{definition}{Definition}{Definitions}
\crefname{example}{example}{examples}
\Crefname{example}{Example}{Examples}
\crefname{remark}{remark}{remarks}
\Crefname{remark}{Remark}{Remarks}
\crefname{warning}{warning}{warnings}
\Crefname{warning}{Warning}{Warnings}
\newcommand{\C}{\mathbb C}
\newcommand{\R}{\mathbb R}
\newcommand{\Q}{\mathbb Q}
\newcommand{\N}{\mathbb N}
\newcommand{\Z}{\mathbb Z}
\newcommand{\No}{\mathbf{No}}
\newcommand{\K}{K}
\newcommand{\Fld}{F}
\newcommand{\mK}{\mathfrak m_K}
\newcommand{\mF}{\mathfrak m_F}
\newcommand{\Ocal}{\mathcal O_K}
\newcommand{\Hil}{\mathcal H_\Gamma}
\newcommand{\Hrf}{\mathcal H_I}
\newcommand{\B}{\mathcal B(H)}
\newcommand{\Ahh}{\mathcal A_\Gamma(H)}
\newcommand{\Arf}{\mathcal A_I}
\newcommand{\Dg}{\mathcal D_I}
\newcommand{\supp}{\operatorname{supp}}
\newcommand{\st}{\operatorname{st}}
\newcommand{\Ran}{\operatorname{Ran}}
\newcommand{\End}{\operatorname{End}}
\newcommand{\coker}{\operatorname{coker}}
\newcommand{\diag}{\operatorname{diag}}
\newcommand{\RCF}{\operatorname{RCF}}
\newcommand{\Cent}{\operatorname{Cent}}
\newcommand{\dist}{\operatorname{dist}}
\newcommand{\Id}{\mathbf 1}
\newcommand{\sH}{\mathop{\sum}\nolimits^{\mathrm H}}
\newcommand{\coeff}[2]{[t^{#1}]#2}
\newcommand{\bigo}{\mathcal O}
\newcommand{\eps}{\varepsilon}
\newcommand{\abs}[1]{\lvert#1\rvert}
\newcommand{\norm}[1]{\lVert#1\rVert}
\newcommand{\ip}[2]{\langle#1,#2\rangle}
\newcommand{\commitA}{608dd23c0539fce73723843949ee627641c27b30}
\newcommand{\commitB}{a3124af79f66b8b9c196d76b4cbc5ac3938907c4}
\newcommand{\commitC}{048b72cf7cbfc8ab246e4f73788c10460cb3f6e0}
% Part III (Drazin halos) and Part IV (Fredholm windows)
\newcommand{\Alg}{\mathfrak A}
\newcommand{\AlgG}{\mathfrak A_\Gamma}
\newcommand{\SigA}{\Sigma^{\Alg}_\Gamma}
\newcommand{\sigD}{\sigma_{\mathrm D,\Alg}}
\newcommand{\acc}{\operatorname{acc}}
\newcommand{\lc}{\operatorname{lc}}
\newcommand{\Tr}{\operatorname{tr}}
\newcommand{\Det}{\operatorname{Det}_\Gamma}
\newcommand{\TT}{\mathscr T}
\newcommand{\SH}{\mathcal S_1(H)}
\newcolumntype{Y}{>{\raggedright\arraybackslash}X}
\hypersetup{pdftitle={Infinite-Dimensional Hahn Spectral Theory},pdfauthor={Research manuscript for the Surreal project},pdfsubject={Inequivalent infinite-dimensional spectra over a Hahn scalar field}}
\begin{document}
\hypersetup{pageanchor=false}
\begin{titlepage}
\centering
\vspace*{12mm}
{\large\scshape A research manuscript for the Surreal project\par}
\vspace{10mm}
{\Huge\bfseries\color{Ink} Infinite-Dimensional\\[.15em] Hahn Spectral Theory\par}
\vspace{10mm}
{\Large Part I: infinitesimal thickening and range defects\\[.1em]
 in Hahn--Hilbert spaces\par}
\vspace{4mm}
{\Large Part II: exact synthesis and algebra-relative spectra\\[.1em]
 for row- and column-finite Hahn matrices\par}
\vspace{4mm}
{\Large Part III: Drazin halos and ramified spectral mapping\\[.1em]
 in arbitrary Hahn operator algebras\par}
\vspace{4mm}
{\Large Part IV: Fredholm windows for compact self-adjoint\\[.1em]
 operators under noncommuting Hahn perturbations\par}
\vspace{10mm}
{\large September 22, 2026\par}
\vfill
\begin{minipage}{.97\textwidth}\small
\textbf{Status.} Four candidate original theorem packages, each with
explicit proofs and its own documented literature comparison. None claims to
solve a named published open problem. Priority is not certified for any of
them. The general proofs have not been independently refereed or checked by
a proof assistant. The four accompanying programs check finite algebraic
identities, not the infinite-dimensional theorems. Classical Hilbert,
Drazin, Fredholm, Hahn--Neumann, Higman and surreal normal-form inputs are
identified rather than presented as new.
\end{minipage}
\newpage
\thispagestyle{empty}
\begin{minipage}{.97\textwidth}\small
\begin{abstract}
\noindent
This report assembles four infinite-dimensional spectral theories over one
Hahn scalar field $\K=\C((t^\Gamma))$. They are printed together because
they continue the same non-claim of the collection's finite-dimensional
spectral report, and they are printed as separate parts because the word
\emph{spectrum} does not mean the same thing in them: the vector spaces and
operator algebras are different. Parts I--III allow arbitrary nonzero
set-sized ordered value groups; Part IV's spectral classification assumes
$\Gamma$ divisible.

\emph{Part I} fixes an ordinary complex Hilbert space $H$ and a nonzero
set-sized ordered abelian group $\Gamma$, and studies
$\Hil=H((t^\Gamma))$ with its coefficientwise-convolution inner product.
An automatic-structure theorem identifies every everywhere-defined
adjointable operator, with no assumed continuity or support condition, with
a Hahn series of bounded operators on $H$. For an ordinary bounded normal
$T$, the set of $\lambda\in\K$ for which $T-\lambda I$ fails to be bijective
on $\Hil$ is exactly $\sigma_\C(T)\cup(\sigma_\C(T)'+\mK)$, where the prime
denotes ordinary accumulation points and $\mK$ the infinitesimals.
Accumulation points thicken into whole infinitesimal monads; isolated points
do not thicken; no new eigenvalues appear. An injective ordinary leading
operator retains its entire algebraic range defect under every
positive-order Hahn perturbation, so a coercive positive self-adjoint
operator can fail to be onto with a cokernel of dimension at least the
continuum.

\emph{Part II} uses the same generality for $\Gamma$ and studies the algebra $\Arf=\RCF_I(\C)((t^\Gamma))$ of Hahn
series of row- and column-finite matrices acting on
$\Hrf=\C^{(I)}((t^\Gamma))$. For $A=D+B$ with $D=\diag(d_i)$ of pairwise
distinct ordinary real entries and $B=B^*$ of globally well-ordered positive
support, there is a unique near-identity normalized Hahn-unitary
diagonalization, exact vectorwise spectral synthesis, a full commutant and
spectral-projection classification, and a coherent resolvent at every
noneigenvalue. The spectrum relative to $\Arf$ is exactly the labeled
eigenvalue set, even when the real residues accumulate. Weighted walks
control every coefficient, and a sharp order $2r+2$ finite-section threshold
holds for one-scale locally finite graphs.

\emph{Part III} takes any nonzero unital complex algebra $\Alg$ and proves
that the spectrum of a constant $T\in\Alg$ relative to $\Alg((t^\Gamma))$ is
$\sigma_\Alg(T)\cup(\sigma_{\mathrm D,\Alg}(T)+\mK)$, with
$\sigma_{\mathrm D,\Alg}$ the finite-index Drazin spectrum, without
normality, norm or divisibility. For $\Alg=\mathcal B(H)$ this is Part~I's
adjointable-algebra spectrum, and it recovers that half of Part~I's normal
formula; for $\Alg=\RCF_I(\C)$ it is Part~II's spectrum of constant
matrices. A ramification criterion decides exactly when polynomial spectral
mapping holds over a nondivisible $\Gamma$.

\emph{Part IV} returns to $H((t^\Gamma))$ for infinite-dimensional separable
$H$ and divisible $\Gamma$. For an
injective compact self-adjoint $T$ and a self-adjoint Hahn perturbation $E$
of positive order that need not commute with $T$, the spectrum of $T+E$ is
the whole monad $\mK$ together with finite Hermitian eigenvalue clusters over
the nonzero eigenvalues of $T$. These local pieces do not assemble: a
trace-class example has no Hahn-unitary diagonalizer with bounded
coefficients and no coefficientwise eigenvalue trace sum or determinant
product, although a coefficientwise Fredholm determinant with an exact
Fredholm alternative exists. The determinant construction and alternative
also hold for arbitrary complex Hilbert spaces and set-sized ordered
exponent groups, without divisibility.

\Cref{ihs:sec:interface} states the interface: what the spectra are, why the
accumulation statements of the parts are not in conflict, and which
hypothesis is doing the work in each.
\end{abstract}
\end{minipage}
\end{titlepage}
\hypersetup{pageanchor=true}
\pagenumbering{roman}
\tableofcontents
\clearpage
\pagenumbering{arabic}

\part*{Shared front matter}
\addcontentsline{toc}{part}{Shared front matter}

\section{Why this is one report, and why it is four parts}\label{ihs:sec:intro}

\subsection{The non-claim that the parts continue}\label{ihs:sec:nonclaim}
The collection already contains a substantial finite-dimensional spectral
report: Hermitian diagonalization, singular values, conditioning,
determinantal scales, spectral descent, and ramification of positive matrix
roots \cite{RepoSpectral}. Every dimension there is an ordinary finite
integer. That report closes by listing, in its own words, what it does
\emph{not} contain:
\begin{quote}
``There is no theorem here about infinite-dimensional Hilbert spaces over
$\No[i]$, bounded-operator spectra, compact or trace-class operators,
or spectral measures.''
\end{quote}
The first two manuscripts assembled here (Parts I and II) supply the first
two of those items and the compact half of the third. The two later ones add
to the second and third items: Part III computes the algebra-relative
spectrum of every constant bounded operator, normal or not, and Part IV
treats compact self-adjoint operators under noncommuting infinitesimal
perturbations and supplies a coefficientwise trace and Fredholm determinant.
A general trace-class theory and a classical Hilbert-space spectral-measure
theory are still not developed; Part II supplies an algebraic projection
calculus with vectorwise Hahn additivity.

\begin{itemize}
\item \emph{Infinite-dimensional Hilbert spaces over the surcomplex
 numbers.} \Cref{ihs:sec:construction,ihs:sec:adjoints} construct one
 explicit such space and determine its adjointable operator algebra
 exactly.
\item \emph{Bounded-operator spectra.} \Cref{ihs:hh:thm:spectrum} computes
 the spectrum of the scalar extension of every ordinary bounded normal
 operator. \Cref{ihs:dz:thm:halo} computes, for every constant element of
 any unital complex algebra, its spectrum relative to the Hahn extension of
 that algebra; for bounded operators this is the adjointable-algebra
 spectrum of Part~I, and the Volterra operator shows that isolated spectral
 points of nonnormal operators can acquire whole halos
 (\cref{ihs:dz:sec:volterra}).
\item \emph{Compact operators.} \Cref{ihs:hh:cor:compact} does the same for
 ordinary compact normal operators, and states explicitly that
 ``compact'' there describes the original operator on the ordinary Hilbert
 space and asserts nothing about compactness or compactoidness in the Hahn
 valuation topology. \Cref{ihs:fr:thm:main} adds an arbitrary
 noncommuting self-adjoint positive-order Hahn perturbation of an injective
 compact self-adjoint operator, with the same meaning of ``compact''.
\item \emph{Trace-class operators and spectral measures.}
 Parts I--III contain no trace-class theory. Part IV defines a
 coefficientwise trace on $\SH((t^\Gamma))$ and a Fredholm determinant on
 its nonnegative-valuation component, with an exact Fredholm alternative
 (\cref{ihs:fr:thm:fredholm}), and proves that ordinary coefficientwise
 eigenvalue sums and products need not represent them
 (\cref{ihs:fr:thm:nolidskii,ihs:fr:thm:noproduct}). This is not a general
 trace-class or nuclear-operator theory. \Cref{ihs:rf:thm:commutant} gives a
 complete Boolean algebra of spectral projections with vectorwise strong
 Hahn additivity. This does not assert topological countable additivity:
 \cref{ihs:rf:rem:notopsum} distinguishes the two and supplies a
 counterexample even for a rank-one value group.
\end{itemize}

The finite-dimensional report is already a long treatise on determinantal
finite-matrix theory. Absorbing these manuscripts into it would bury them.
This report is its sibling, not its continuation, and
\cref{ihs:sec:interface} records what the reports may and may not be read
as saying jointly. Parts III and IV arrived after Parts I and II and are
built on them: Part III generalizes the algebra half of Part~I's normal
spectral formula, and Part IV treats the case that Part~I sets aside as
requiring a different theorem (\cref{ihs:hh:sec:different}).
\Cref{ihs:app:provenance} records how the four manuscripts were assembled.

\subsection{The central hazard: several senses of ``spectrum''}\label{ihs:sec:hazard}
The parts use the word \emph{spectrum} for different things, on different
vector spaces and in different operator algebras. Quoted side by side without
their categories, the conclusions of Parts I and II about accumulation read
as a contradiction. They are not one. Parts I--III allow arbitrary nonzero
set-sized ordered abelian groups; the coefficient spaces and operator classes,
rather than divisibility, account for the different spectral conclusions.
The original Part I manuscript assumed divisibility. The proof review removes
that unused restriction, as explained in \cref{ihs:hh:sec:field}. Part IV's
spectral classification keeps divisibility; its determinant and Fredholm
alternative need no such restriction (\cref{ihs:fr:sec:determinant}).

There are exactly two kinds of spectrum in this report, and the parts use
three instances of them. A \emph{bijectivity spectrum} is the set of $\lambda$
for which $X-\lambda I$ fails to be bijective on a named vector space: this
is $\sigma^{\mathrm{alg}}_\Gamma$ of Part~I, also used in Part~IV. An
\emph{algebra-relative spectrum} is the set of $\lambda$ for which
$X-\lambda1$ has no two-sided inverse in a named unital algebra: this is
$\sigma^{\mathrm{adj}}_\Gamma$ (the algebra $\Ahh$) in Parts I and IV,
$\sigma_{\Arf}$ (the algebra $\Arf$) in Part~II, and $\SigA$ (the algebra
$\AlgG=\Alg((t^\Gamma))$, for an arbitrary unital $\Alg$) in Part~III. The
third notion, $\SigA$, is defined for constant elements $T\in\Alg$ and
contains the constant cases of the other two algebra-relative spectra:
$\Sigma^{\B}_\Gamma=\sigma^{\mathrm{adj}}_\Gamma$ and
$\Sigma^{\RCF_I(\C)}_\Gamma=\sigma_{\Arf}$ (\cref{ihs:prop:unify}). No
statement of this report identifies a bijectivity spectrum with an
algebra-relative one except where a theorem proves it for the operators in
question (\cref{ihs:hh:thm:spectrum,ihs:fr:thm:main,ihs:fr:thm:fredholm},
and the Volterra and block examples
\cref{ihs:dz:sec:volterra,ihs:dz:sec:blocks}).

\begin{table}[p]
\centering\footnotesize
\setlength{\tabcolsep}{3.5pt}
\begin{tabularx}{\textwidth}{@{}>{\raggedright\arraybackslash}p{.13\textwidth}YYYY@{}}
\toprule
 & Part I & Part II & Part III & Part IV\\
\midrule
Vector space &
$\Hil=H((t^\Gamma))$, $H$ an ordinary complex Hilbert space; every ordinary
square-summable vector is a constant of $\Hil$ &
$\Hrf=\C^{(I)}((t^\Gamma))$; each Hahn coefficient has \emph{finite}
coordinate support, but a vector may have infinitely many nonzero Hahn
coordinates &
none is needed; for $\Alg=\B$ the algebra acts on $\Hil$, for
$\Alg=\RCF_I(\C)$ on $\Hrf$ &
$\Hil$ of Part~I, with $H$ infinite-dimensional and separable\\
\addlinespace
Operator class &
$\Ahh=\B((t^\Gamma))$, Hahn series of \emph{ordinary bounded} operators on
$H$ &
$\Arf=\RCF_I(\C)((t^\Gamma))$, Hahn series of row- and column-finite
matrices, with no ordinary boundedness at all &
\emph{constant} elements $T$ of $\AlgG=\Alg((t^\Gamma))$, $\Alg$ any nonzero
unital complex algebra &
$A=T+E\in\Ahh$, $T$ injective compact self-adjoint, $E=E^*$ of positive
order, not assumed to commute with $T$\\
\addlinespace
``Spectrum'' means &
$\sigma^{\mathrm{alg}}_\Gamma(A)=\{\lambda\in\K:A-\lambda I$ is not
bijective on $\Hil\}$; for $A\in\Ahh$ also
$\sigma^{\mathrm{adj}}_\Gamma(A)$, failure of invertibility in $\Ahh$ &
$\sigma_{\Arf}(X)=\{z\in\K:X-z\Id$ has no two-sided inverse
\emph{inside} $\Arf\}$: a spectrum relative to a named unital
$\K$-algebra &
$\SigA(T)$: no two-sided inverse of $T-\lambda1$ inside $\AlgG$; equal to
$\sigma^{\mathrm{adj}}_\Gamma$ for $\Alg=\B$ and to $\sigma_{\Arf}$ for
$\Alg=\RCF_I(\C)$ &
both Part~I spectra, proved equal for this class; also the algebraic
Fredholm spectrum $\sigma_{\mathrm{ess,alg}}$\\
\addlinespace
Hypothesis on $\Gamma$ &
nonzero, set-sized, ordered abelian, \textbf{not assumed divisible};
only the metric-completeness subsection
\ref{ihs:hh:sec:rankone} assumes $\Gamma\subseteq\R$ &
nonzero, set-sized, ordered abelian, \textbf{not assumed divisible}; the
source manuscript records the absence of that hypothesis as a feature, and
only near-one binomial roots are used &
nonzero, set-sized, ordered abelian, \textbf{not assumed divisible}; the
polynomial image depends on divisibility (\cref{ihs:dz:thm:criterion}) &
nonzero, set-sized, ordered abelian and \textbf{divisible}\\
\addlinespace
What accumulation does &
accumulation points of $\sigma_\C(T)$ acquire whole complex infinitesimal
monads (\cref{ihs:hh:thm:spectrum}); isolated points do not thicken &
accumulation of the ordinary residues $d_i$ does not enlarge
$\sigma_{\Arf}(A)$ beyond $\{\lambda_i\}$
(\cref{ihs:rf:thm:resolvent,ihs:rf:ex:accum}) &
every point of the Drazin spectrum acquires a whole halo; in a Banach
algebra these are the accumulation points and the isolated non-poles
(\cref{ihs:dz:thm:halo}, \eqref{ihs:dz:eq:dictionary}) &
the accumulation point $0$ keeps the whole monad $\mK$, with no eigenvalues
in it; each nonzero eigenvalue becomes a finite Hermitian cluster
(\cref{ihs:fr:thm:main})\\
\bottomrule
\end{tabularx}
\caption{The senses of ``spectrum'' used in this report. Every spectral
statement below names its vector space or algebra and which of these
definitions it uses.}\label{ihs:tab:senses}
\end{table}

The Part IV column describes its spectral classification. Its separate
determinant and Fredholm alternative in \cref{ihs:fr:sec:determinant}
allow arbitrary complex Hilbert $H$ and set-sized ordered abelian $\Gamma$,
including the zero group.

\Cref{ihs:sec:interface} works the apparent conflict between Parts I and II
out on a single operator, says exactly why both conclusions hold, and shows
that Part~III's formula produces both of them.

The inner-product conventions also differ: Parts I and IV are linear in the
first variable, whereas Part II is conjugate-linear in the first. To
translate a scalar pairing, conjugate its value (equivalently, exchange its
arguments); do not copy a variational functional without translating its
linearity. Positivity and the adjoint identity are unchanged by this
conversion. Part III uses no inner product.

\subsection{Part I at a glance}\label{ihs:sec:glanceone}
Throughout Part I, $H\ne0$ is a complex Hilbert space, $\Gamma$ is nonzero,
ordered, abelian and set-sized, and
\[
 \K=\C((t^\Gamma)),\qquad \Hil=H((t^\Gamma)).
\]
The valuation is the least support exponent. A prime on an ordinary compact
set denotes its accumulation set, not a derivative and not an essential
spectrum.

\begin{theorem}[Part I theorem package]\label{ihs:hh:thm:package}
The following assertions hold in this model.
\begin{enumerate}[label=\textup{(\roman*)}]
\item Every everywhere-defined adjointable $\K$-linear operator on $\Hil$
has a unique expansion in $\B((t^\Gamma))$. Its coefficient operators are
ordinary bounded operators and have one common well-ordered support.
\item For an ordinary bounded normal $T$, its coefficientwise extension has
\begin{equation}\label{ihs:hh:eq:main-spectrum}
 \sigma_\Gamma(T)=\sigma_\C(T)\cup
       \bigl(\sigma_\C(T)'+\mK\bigr),
 \qquad \sigma_{p,\Gamma}(T)=\sigma_{p,\C}(T).
\end{equation}
Here $\sigma_\Gamma$ is the failure-of-bijectivity spectrum on $\Hil$; it
also equals the spectrum in the adjointable algebra.
\item If $S:V\to V$ is injective and $E$ is an endomorphism-valued Hahn
series of strictly positive order, then
\[
 \coker(S+E)\cong W((t^\Gamma)),\qquad
 V=\Ran S\oplus W,
\]
noncanonically. No commutation hypothesis is required.
\item For $D e_n=n^{-1}e_n$ and $s=t^\eta$, $\eta>0$, the operator
$C=D^2+s^2I$ on $\ell^2((t^\Gamma))$ is bounded, positive, self-adjoint and
coercive, but not surjective. Its cokernel has $\K$-dimension at least
$2^{\aleph_0}$. At every rank its range is closed and proper in the
valuation topology, with zero orthogonal complement.
\item The extension of an ordinary bounded $T$ has a least
$\Fld$-valued norm bound if and only if $T$ attains its ordinary Hilbert
operator norm. When the least bound exists, it is the ordinary norm itself.
\end{enumerate}
\end{theorem}

Each assertion is proved in Part I. The result about the least norm bound is
not a claim that these operators are unbounded: all operators in (i) have
field-valued bounds. The failure is the absence of a \emph{least} such
bound. Likewise, a spectral point in (ii) need not be an eigenvalue.

\subsection{Part II at a glance}\label{ihs:sec:glancetwo}
Throughout Part II, $\Gamma$ is a set-sized ordered abelian group with no
divisibility assumption, $I$ is a nonempty index set, and
\[
 \Fld=\R((t^\Gamma)),\quad \K=\C((t^\Gamma)),\quad
 \Arf=\RCF_I(\C)((t^\Gamma)),\quad \Hrf=\C^{(I)}((t^\Gamma)).
\]
Let $D=\diag(d_i)$ with the $d_i\in\R$ pairwise distinct, let
$B=B^*\in\Arf$ have positive support, and set $A=D+B$.

\Cref{ihs:rf:thm:unitary} constructs a unique pair $(U,\Lambda)$ with
\[
 AU=U\Lambda,\qquad U^*U=UU^*=\Id,\qquad
 U=\Id+\text{positive-support terms},
\]
where $\Lambda=\diag(\lambda_i)$, $\lambda_i-d_i$ has positive support, and
every $U_{ii}$ is positive real Hahn. All exponents of the corrections
belong to the additive monoid generated by $\supp B$; the construction uses
no extension of the value group. The distinct-real-label hypothesis entails
$\abs I\le\abs\R$; no claim of a larger residue index set is implicit.

\Cref{ihs:rf:thm:synthesis,ihs:rf:thm:commutant,ihs:rf:thm:resolvent} then
give exact eigenvector expansions, classify every commuting operator and
spectral projection, and prove
$\sigma_{\Arf}(A)=\{\lambda_i:i\in I\}$, the spectrum relative to the
specified unital $\K$-algebra and not a Banach-algebra spectrum.
\Cref{ihs:rf:thm:walks,ihs:rf:thm:marked,ihs:rf:thm:finite,ihs:rf:thm:sharp}
determine how much of an infinite operator a single coefficient can see: an
eigenvalue coefficient is supported by closed weighted walks, a change
outside a finite section can first affect it only through a walk that leaves
the section and returns, and the corresponding valuation threshold is exact
in a family of path examples.

\subsection{Part III at a glance}\label{ihs:sec:glancethree}
Throughout Part III, $\Alg$ is any nonzero unital complex algebra and
$\Gamma$ any nonzero set-sized ordered abelian group; no norm, normality or
divisibility is assumed. For a constant $T\in\Alg$,
\cref{ihs:dz:thm:halo} proves
\[
 \SigA(T)=\sigma_\Alg(T)\cup\bigl(\sigD(T)+\mK\bigr),
\]
where $\sigD(T)$ is the set of $a\in\C$ at which $T-a1$ is not
finite-index Drazin invertible. The mechanism is an exact inverse descent
(\cref{ihs:dz:thm:descent}): a strong automorphism of the Hahn field moves
any nonzero infinitesimal to a monomial, a cyclic-support projection reduces
a putative inverse to a finite-tail Laurent inverse, and such an inverse
exists exactly at Drazin points, with exact valuation $-\nu v(\eps)$ at index
$\nu$ (\cref{ihs:dz:thm:scale}). In a Banach algebra the Drazin spectrum
consists of the accumulation points and the isolated non-poles
\eqref{ihs:dz:eq:dictionary}, which recovers the adjointable-algebra half of
Part~I's normal formula (\cref{ihs:dz:cor:normal,ihs:dz:rem:recover}); the
Volterra operator has spectrum $\{0\}$, no accumulation point, and
$\Sigma^{\B}_\Gamma(\mathcal V)=\mK$. A polynomial germ of ramification index
$e$ maps a halo, after subtracting its ordinary image, onto zero and exactly
the nonzero infinitesimals with valuation in $e\Gamma$,
which gives a necessary and sufficient criterion for spectral mapping over
every value group (\cref{ihs:dz:thm:image,ihs:dz:thm:criterion}).

\subsection{Part IV at a glance}\label{ihs:sec:glancefour}
For Part IV's spectral classification, $H$ is infinite-dimensional and separable, $\Gamma$ is
\emph{divisible}, and the space and algebra are $\Hil$ and $\Ahh$ of Part~I.
For $A=T+E$ with $T$ ordinary injective compact self-adjoint and
$E=E^*\in\Ahh$ of positive order, not necessarily commuting with $T$,
\cref{ihs:fr:thm:main} proves
\[
 \sigma^{\mathrm{adj}}_\Gamma(A)=\sigma^{\mathrm{alg}}_\Gamma(A)=\mK\cup
 \bigcup_{c\in\sigma_\C(T)\setminus\{0\}}\{\lambda_{c,1},\ldots,\lambda_{c,m_c}\},
\]
with the second set exactly the point spectrum, and
$\sigma_{\mathrm{ess,alg}}(A)=\mK$. The finite clusters come from
coefficientwise Riesz projections and a direct rotation. They do not
assemble into a global spectral theorem in $\Ahh$: for
$A_\star=\diag(2^{-n})+t^\eta Y$ with a trace-class weighted path $Y$, no
Hahn unitary in $\Ahh$ diagonalizes $A_\star$ (\cref{ihs:fr:thm:nounitary}),
and the coefficientwise eigenvalue sum and product diverge although the
coefficientwise trace and Fredholm determinant exist
(\cref{ihs:fr:thm:nolidskii,ihs:fr:thm:noproduct,ihs:fr:thm:fredholm}). No
nonzero coherent Hahn analytic function detects the whole monad at the
accumulation point (\cref{ihs:fr:thm:nocharacteristic}).
The determinant and Fredholm alternative themselves hold for arbitrary
complex Hilbert spaces and set-sized ordered abelian exponent groups;
\cref{ihs:fr:sec:determinant} states this separate scope.

\subsection{What is and is not claimed as new}\label{ihs:sec:novelty}
No part claims a classical input as new, and the four parts have separate
novelty boundaries and separate literature comparisons.

For Part I: the ordinary spectral theorem, Baire category, the closed graph
theorem, Hilbert norm estimates and Hahn support lemmas are classical inputs
\cite{Douglas,Neumann}. The non-Archimedean failure of orthogonal
decomposition is known in other models, including work of Aguayo and Nova
\cite{AN} and Aguayo, Nova and Shamseddine \cite{ANS}; we do not claim to
have discovered that failure in general. The proposed contribution is the
particular combination of automatic adjoint structure, the exact
accumulation-monad spectrum, the arbitrary positive-support range-defect
formula, and the explicit coercive example with a continuum-dimensional
defect in $H((t^\Gamma))$. In the repository sections and primary literature
inspected, these statements were not found in this form. This is a qualified
report of the search, not a proof of priority; a theorem may have
predecessors under terminology not located by that search.

For Part II: simple-eigenvalue perturbation theory and its normalization
conventions are classical, and Greenbaum, Li and Overton give an accessible
account of both eigenvalue and eigenvector perturbation \cite{glo};
fixed-point and near-diagonal formulations are also established, including
the work of Kenmoe, Smerlak and Zadorin \cite{ksz}. The recursive
Rayleigh--Schr\"odinger mechanism and the usual first corrections are not
claimed as new. The support argument uses Neumann's ordered-series machinery
\cite{Neumann} and Higman's word-order theorem \cite{Higman}. Matrix
products and walks have a long-established connection; the path-sum
framework of Giscard, Thwaite and Jaksch is an explicit reference
\cite{pathsum}, and the walk theorem below is not a claim to have discovered
that connection. The coefficient algebra is likewise not new: Ara, Goodearl
and O'Meara study real and complex row- and column-finite matrices and
emphasize that positivity in that algebra does not force the existence of
eigenvalues or square roots \cite{ago}, which is why an infinite-dimensional
result must state both its vector space and its operator algebra. The
elementary no-eigenvector example in \cref{ihs:sec:failures} is a boundary
test, not a new phenomenon. The proposed contribution is the combined
arbitrary-rank Hahn theorem: a globally coherent normalized diagonalizer for
infinitely many simple residue levels, exact synthesis on a specified Hahn
module, coherent resolvents at every noneigenvalue, and marked-walk
stability with sharp finite-section thresholds. Neither the finite report
inspected in \cite{RepoPinB} nor the cited perturbation and row-finite
sources state that package. This is a targeted comparison, not an exhaustive
theorem of bibliographic absence; some components may have predecessors in
formal operator or valued-algebra literature not located in this search.

For Part III: Drazin inverses and their finite-index identities, the
resolvent-pole characterization, Banach-algebra Drazin spectra, the
empty-Drazin-spectrum characterization of algebraic elements, ordinary
polynomial spectral mapping and the nilpotent ceiling formula are classical
\cite{Drazin,BoassoOperators,Boasso}; generalized Drazin invertibility
\cite{Koliha} is a different condition. Strong Hahn-automorphism
constructions belong to the surrounding literature \cite{KaplanKrappSerra}.
The constant-normal formula is Part~I's. The proposed contribution is the
arbitrary-infinitesimal inverse descent in the original value group, the
resulting halo formula for every unital complex algebra, the exact inverse
valuation, and the ramification criterion identifying every missing
polynomial-image valuation. The algebraicity criterion is a Hahn
interpretation of a classical fact.

For Part IV: compact spectral theory, isolated-cluster perturbation
\cite{Kato}, the direct rotation of two projections \cite{SimonProjections},
ordinary trace-class determinant identities \cite{Bornemann}, Hahn-field
algebraic closedness \cite{Poonen} and non-Archimedean Fredholm theory in
Serre's different category \cite{Serre} are credited, and the range-defect
factorization is Part~I's. The proposed contribution is the exact
arbitrary-rank spectrum and algebraic Fredholm spectrum for the stated
noncommuting compact-residue class, the trace-class example with no
bounded-coefficient Hahn-unitary diagonalizer and no coefficientwise
eigenvalue trace sum or determinant product, the support-controlled Hahn
Fredholm alternative with its finite-coupling domain and accumulation
obstruction, and the separation of point-spectrum invariance from monad
growth under enlargement of the value group. Both comparisons were targeted,
at the pin recorded in \cref{ihs:app:auditC,ihs:app:auditD}, and certify no
priority.

\begin{table}[ht]
\centering\small
\begin{tabularx}{\textwidth}{@{}p{.25\textwidth}YY@{}}
\toprule
Strand & Established input or predecessor & Contribution here\\
\midrule
Finite surcomplex matrices & Repository spectral chapter & Infinite coefficient Hilbert spaces (Part I); arbitrary index sets (Part II)\\
Hahn formal calculus & Neumann support lemma & Automatic global support for adjoints (Part I)\\
Ordinary normal spectra & Hilbert spectral theorem & Exact infinitesimal thickening formula (Part I)\\
Formal perturbation & Neumann inversion near the identity & Persistence of arbitrary injective range defects (Part I)\\
Non-Archimedean orthogonality & Known nonorthomodular examples & A closed coercive range with continuum defect (Part I)\\
Operator norm bounds & Ordinary norm and attainment & Exact least-bound criterion over the Hahn field (Part I)\\
Perturbation recursions & Rayleigh--Schr\"odinger corrections & Global support certificate for all orders at once (Part II)\\
Row- and column-finite matrices & Ara--Goodearl--O'Meara & Hahn diagonalization under simple ordinary residues (Part II)\\
Matrix walks & Path-sum framework & Marked-walk stability and sharp section thresholds (Part II)\\
Drazin inverses & Finite-index identities, resolvent poles & Inverse descent at every Hahn infinitesimal; halo formula in any unital algebra (Part III)\\
Polynomial spectral mapping & Ordinary Drazin spectral mapping & Exact ramified image and value-group criterion (Part III)\\
Compact perturbation & Riesz projections, direct rotation & Exact spectrum under noncommuting Hahn perturbations (Part IV)\\
Fredholm determinants & Ordinary trace-ideal identities & Hahn Fredholm alternative; failure of spectral sums and products (Part IV)\\
\bottomrule
\end{tabularx}
\caption{The novelty boundary: classical tools are not renamed as new
theorems.}\label{ihs:tab:novelty}
\end{table}

\subsection{Relation to the nearby non-Archimedean literature}\label{ihs:sec:nearby}
Aguayo--Nova--Shamseddine work with free Banach spaces of countable type
over the complex Levi--Civita field, represented by weighted $c_0$-type
spaces \cite{ANS}. Ishizuka proves spectral results for self-adjoint
compactoid operators over valued fields with formally real residue field
\cite{Ishizuka}. These are nearby but different settings. The constant copy
of $H$ in the Part I space retains every ordinary square-summable vector,
while its valuation topology is discrete; the extension of an ordinary
compact operator is not thereby compactoid in the new topology.

The June 2026 preprint of Miabey treats diagonal operators and finite-rank
perturbations in a non-Archimedean Hilbert-space setting \cite{Miabey}. It
is used only to delimit related contemporary work; none of its theorems is
an input to the proofs here. Formal resolvent expansions also have a
classical relationship with pole and Drazin-spectrum questions; see Boasso
\cite{Boasso}. The normal-operator formula of Part I is proved directly, so
no Drazin-theoretic equivalence is needed as an unproved bridge. Part III
now proves that equivalence for constant elements of any unital algebra
(\cref{ihs:dz:thm:descent}), and so gives a second proof of the
algebra-relative half of Part~I's formula (\cref{ihs:dz:rem:recover}); the
Part~I proof does not depend on it.

For Part II, the relevant nearby object is the algebra of ordinary row- and
column-finite matrices itself \cite{ago}. The reference concerns ordinary
such matrices and is not a source for the Hahn diagonalization asserted
here.

\subsection{Set-sized mathematics and terminology}\label{ihs:sec:setsized}
All index sets, exponent groups, supports and families in the proofs are
sets, and we work in ordinary ZFC. The passages to $\No[i]$ in
\cref{ihs:sec:surcomplex,ihs:sec:transport} are transports of set-sized Hahn
objects, not formations of a set containing all surreal numbers. The full
surreal class offers additional scales but is not a set-sized coefficient
field in ZFC; the proofs below avoid class-indexed summations and
class-sized choices by fixing the workspace at the outset.

``Positive support'' means that every exponent is strictly greater than
zero; it does not mean that the complex coefficients are positive.
``Hermitian'' means $A^*=A$. In Part II, ``unitary'' means two-sided
algebraic unitarity in $\Arf$, and a ``Hahn-orthonormal basis'' means a
family with the strong-summation synthesis property proved in
\cref{ihs:rf:thm:synthesis}, not a Hamel basis and not an ordinary
Hilbert-space basis.

\section{Shared Hahn preliminaries}\label{ihs:sec:prelim}

This section fixes the Hahn conventions used by all four parts and states
the support lemma once. Each part specifies its own coefficient space,
operator algebra and, where present, inner product; these are not shared
merely because the scalar field is the same.

\subsection{Support conventions and strong summability}\label{ihs:sec:supportconv}
A Hahn series with coefficients in a complex vector space $V$ is a family
\[
 x=\sum_{\gamma\in\Gamma}x_\gamma t^\gamma,\qquad x_\gamma\in V,
\]
whose support $\{\gamma:x_\gamma\ne0\}$ is well ordered in the given order
of $\Gamma$. Write $V((t^\Gamma))$ for the resulting set. The zero series has
empty support. For $x\ne0$ put $v(x)=\min\supp x$, and put $v(0)=+\infty$.

A family $(x_a)_{a\in J}$ in $V((t^\Gamma))$ is \emph{Hahn summable}, also
called \emph{strongly Hahn summable}, when its union of supports is well
ordered and, for every $\gamma$, only finitely many $a$ have
$(x_a)_\gamma\ne0$. Its sum, written $\sH_a x_a$, is defined coefficientwise
by those finite sums. All unqualified infinite formal sums in this report
mean this notion. It is not convergence of ordinary partial sums in the full
surreal fine topology, and \cref{ihs:rf:prop:noconv} exhibits a strongly
summable family whose partial sums do not converge even in the intrinsic
valuation topology of $\K$.

\subsection{The support lemma, stated once}\label{ihs:sec:supportlemma}
The first two source manuscripts open with the Hahn--Neumann support lemma, one
citing it and supplying the two-set finiteness argument, the other
reproving Higman's word-order input by a minimal bad sequence. The lemma is
classical, and it is already printed with a full proof elsewhere in this
collection: the finite-dimensional spectral report proves the word-order
statement as its \texttt{spec:lem:words} and the positive-monoid statement
as its \texttt{spec:lem:neumann}, in its section ``A self-contained support
engine'' \cite{RepoSpectral}. We therefore state it here once, in the form
all four parts use, and do not reprint that proof.

\begin{lemma}[Classical Hahn--Neumann support lemma]\label{ihs:lem:support}
Let $\Gamma$ be an ordered abelian group.
\begin{enumerate}[label=\textup{(\roman*)}]
\item If $A,B\subseteq\Gamma$ are well ordered, then $A+B$ is well ordered
and each $\gamma$ has only finitely many representations $\gamma=a+b$ with
$a\in A$ and $b\in B$. A finite union of well-ordered subsets is well
ordered, and the corresponding statement holds for any finite number of
summands.
\item If $S\subseteq\Gamma_{>0}$ is well ordered, then the additive monoid
\[
 S^*=\{s_1+\cdots+s_n:n\ge0,\ s_j\in S\},
\]
including the empty sum $0$, is well ordered, and for each $\gamma\in\Gamma$
there are only finitely many \emph{ordered words} $(s_1,\ldots,s_n)$ over
$S$ with sum $\gamma$, the length $n$ being allowed to vary.
\end{enumerate}
\end{lemma}

This is the standard support lemma \cite{Neumann}, whose word form rests on
Higman's theorem that finite words over a well-quasi-ordered alphabet are
well-quasi-ordered \cite{Higman}. Part (i) is elementary and both source
manuscripts give the argument: infinitely many distinct $a$ in a fixed-sum
family yield an increasing subsequence in $A$ and hence a strictly
decreasing subsequence $\gamma-a$ in $B$, a contradiction; and every
sequence in a well-ordered set has an infinite nondecreasing subsequence,
obtained by successively choosing an occurrence of the minimum of the
remaining tail. Part (ii) is the usual Neumann lemma in its positive-support
form. Both are stated classical inputs here, not new combinatorial claims.

\emph{Part (ii) from Higman's theorem, as the Part~IV source deduces it.}
Order finite words over the well-ordered alphabet $S$ by subsequence
embedding with coordinatewise increase; Higman's theorem makes this a
well-quasi-order \cite{Higman}. In any infinite sequence of words there are
indices $i<j$ with the $i$th word embedded in the $j$th, and positivity of
the letters makes the total of the first at most that of the second; so no
infinite strictly decreasing sequence of totals exists. For a fixed total, a
strict embedding either inserts a positive letter or strictly increases a
letter, and so strictly increases the total; hence infinitely many distinct
words cannot share one total. This deduction does not assume that $n\min S$
is cofinal in $\Gamma$.

\begin{remark}[Why the full word form is needed]\label{ihs:rem:wordform}
All four parts use the finite-\emph{word} assertion of
\cref{ihs:lem:support}(ii), not only the assertion that each fixed power is
meaningful. The latter alone would not justify a geometric series. The
estimate $v(X^n)\ge n\eps$ is also not a substitute: in a value group of
higher rank the sequence $n\eps$ may be bounded above, as
\cref{ihs:sec:ranktwo} exhibits concretely.
\end{remark}

\Cref{ihs:lem:support}(i) licenses scalar multiplication on $V((t^\Gamma))$
and products of operator-valued Hahn series, even when the coefficient
algebra is noncommutative; a finite number of supports is combined by
iteration.

\subsection{Support-certified inversion and formal substitution}\label{ihs:sec:neumann}
\begin{lemma}[Support-certified Neumann inverse]\label{ihs:lem:neumann}
Let $\mathscr R$ be a unital complex algebra and let
$E\in\mathscr R((t^\Gamma))$ have $v(E)>0$. Then
\begin{equation}\label{ihs:eq:neumann}
 (I+E)^{-1}=\sum_{n\ge0}(-E)^n
\end{equation}
is a well-defined two-sided inverse in $\mathscr R((t^\Gamma))$, with
support in $(\supp E)^*$.
\end{lemma}
\begin{proof}
Expand each coefficient of $E^n$ as a sum over ordered words in $\supp E$.
\Cref{ihs:lem:support} makes the union well ordered and all coefficients
finite, including the sum over $n$; the products of the resulting series are
the formal convolution coefficients, and the geometric identity is a formal
one-variable identity. Noncommutativity causes no problem, because every
factor is a power of the same $E$.
\end{proof}

\begin{remark}[A second proof, from the first source]\label{ihs:rem:telescope}
The Part~I manuscript proves \cref{ihs:lem:neumann} by a different route,
which is recorded here because it uses no formal-identity bookkeeping. After
\cref{ihs:lem:support} has made every coefficient of $\sum_{n\ge0}(-E)^n$ a
finite sum, the finite identity
$(I+E)\sum_{n=0}^{N}(-E)^n=I+(-1)^N E^{N+1}$ telescopes coefficientwise
across all word lengths, since at a fixed exponent only finitely many word
lengths occur; right multiplication telescopes as well. This argument uses
associativity but not commutativity of $\mathscr R$.
\end{remark}

\begin{corollary}[Evaluation of homogeneous constructions]\label{ihs:cor:evaluate}
Let $\mathscr R$ be a unital complex algebra, let $S\subseteq\Gamma_{>0}$ be
well ordered, and let $X_n\in\mathscr R((t^\Gamma))$ satisfy
$\supp X_n\subseteq nS$ for every $n\ge1$. Then $(X_n)_{n\ge1}$ is strongly
summable. Finite sums, products, coefficientwise linear maps preserving
supports, and regroupings of such homogeneous constructions commute with
their evaluation $\sH_{n\ge1}X_n$. The same holds with $\mathscr R$ replaced
by a module over it, or by $\K$ itself.
\end{corollary}
\begin{proof}
The union of supports lies in $S^*$. At a fixed exponent,
\cref{ihs:lem:support}(ii) bounds the possible lengths and the possible
ordered exponent words. All coefficient operations within one word are
finite, so the claimed interchanges are finite rearrangements coefficient by
coefficient.
\end{proof}

\begin{corollary}[Formal power series at a positive-support argument]\label{ihs:cor:binomial}
If $X\in\mathscr R((t^\Gamma))$ has positive support, then any one-variable
formal power series with complex coefficients can be evaluated at $X$, and
formal identities between such series remain valid whenever their
substitutions have zero constant term. In particular $(I+X)^{-1}$ is given
by \eqref{ihs:eq:neumann} and $(I+X)^{-1/2}$ is defined by its binomial
series.
\end{corollary}
\begin{proof}
Take $S=\supp X$ in \cref{ihs:cor:evaluate}.
\end{proof}

\subsection{The surcomplex realization}\label{ihs:sec:surcomplex}
Suppose an ordered additive embedding $\jmath:\Gamma\hookrightarrow\No$ is
fixed; when $\Gamma$ is itself a set-sized ordered subgroup of the surreal
additive group, take $\jmath$ to be the inclusion. Conway normal forms give
the maps
\begin{equation}\label{ihs:eq:surreal-embedding}
 \sum_\gamma r_\gamma t^\gamma
 \longmapsto\sum_\gamma r_\gamma\omega^{-\jmath(\gamma)},
 \qquad r_\gamma\in\R,
\end{equation}
and their coefficientwise complexification. These are an ordered-field
embedding $\Fld\hookrightarrow\No$ and a field embedding
$\K\hookrightarrow\No[i]$, preserving conjugation and the prescribed strong
sums. The classical normal-form and multiplication facts are standard
surreal theory \cite{Gonshor}.

The sign in \eqref{ihs:eq:surreal-embedding} matters. A well-ordered set of
Hahn exponents becomes a reverse-well-ordered set of Conway exponents,
exactly the normal-form condition, and positive Hahn exponents become
infinitesimal monomials. All supports are sets, so the normal forms define
actual surreal or surcomplex numbers, and the vector spaces below never
become proper classes. In particular, $\Gamma=\Q$ and $t=\omega^{-1}$
suffice for all explicit counterexamples of Part I.

\begin{warning}[Distinct completions and topologies]\label{ihs:warn:three}
The Part I construction $H((t^\Gamma))$ is neither the algebraic tensor
product $H\otimes_\C\K$ nor a silently chosen sequence space $\ell^2(\K)$.
The Part II space $\C^{(I)}((t^\Gamma))$ is likewise not ordinary $\ell^2$,
and its vectors have finite-support \emph{coefficients} without having
finitely many nonzero Hahn coordinates. Nor does the intrinsic valuation
topology of either space equal the topology induced by embedding its scalar
field into the full surreal fine topology, in which the set-sized scalar
workspace is discrete \cite{RepoMap}. All of these distinctions matter to
the results.
\end{warning}

The closure of $H\otimes_\C\K$ in $H((t^\Gamma))$ is determined exactly by
\texttt{duals:thm:completion} and \texttt{duals:thm:cofinality-completion} of
\path{docs/surcomplex/three-duals-of-hahn-vector-spaces/}.

\section{The interface between the parts}\label{ihs:sec:interface}

\subsection{What is shared and what is not}\label{ihs:sec:sharednot}
Parts I and II share the scalar field $\K$, the support conventions of
\cref{ihs:sec:supportconv}, the support lemma \cref{ihs:lem:support}, the
inversion tools of \cref{ihs:sec:neumann}, and the surcomplex embedding
\eqref{ihs:eq:surreal-embedding}. Nothing else is shared. In particular no
theorem of Part I is used in the proof of any theorem of Part II, and no
theorem of Part II is used in Part I. Parts III and IV use the same shared
tools. Part IV works in Part~I's space and algebra and uses Part~I's
range-defect theorem \cref{ihs:hh:thm:defect}, for which its source gives a
second proof (\cref{ihs:fr:rem:secondproof}). The main theorems of Part III
use no theorem of Parts I, II or IV; its comparison with them is
\cref{ihs:sec:unify}. No theorem of Part III or IV is used in Part I or
Part II. The coefficient-space and algebra requirements remain local to
each part. Parts I and II prove positivity for their respective inner
products; Part III requires no inner product, and Part IV inherits Part~I's.
In particular,
\cref{ihs:hh:prop:inner} constructs the positive square root
$\norm x=\sqrt{\ip xx}$ from the even leading exponent of $\ip xx$ and a
near-one binomial root; this works for every ordered value group.
Likewise, \cref{ihs:rf:prop:inner} proves positivity without assuming that
every positive scalar has a square root. Part II uses no norm, and
\cref{ihs:rf:cor:scales} explicitly allows nondivisible groups.

Neither construction needs divisibility, and neither does Part III, whose
straightening automorphism adjoins no fractional exponent
(\cref{ihs:dz:lem:straighten}). Part IV's spectral classification assumes
divisibility to diagonalize finite Hermitian clusters inside the field
(\cref{ihs:fr:lem:hermitian}); its determinant and Fredholm alternative
do not (\cref{ihs:fr:sec:determinant}). In Part I all field-valued roots
used in norms and moduli have even leading exponent, and all other scalar
roots are near-one binomial roots; ordinary positive operator square roots
are taken in $\B$, before Hahn extension. General positive Hahn scalars
and squared norms have different root requirements:
$t$ has no square root in $\R((t^\Z))$, whereas $\norm{te}=t$ for any
ordinary unit vector $e$.

\subsection{The accumulation statements are about different objects}\label{ihs:sec:accum}
Read without their categories, the two parts appear to disagree about
accumulation. Part I proves that accumulation points of the ordinary
spectrum thicken into infinitesimal monads. Part II proves that the
algebra-relative spectrum is exactly the labeled eigenvalue set even when
the real residues accumulate. One operator makes the difference visible.

Let $I=\N_{\ge1}$ and consider the diagonal data $d_n=1/n$, whose residues
accumulate at $0$.
\begin{itemize}
\item As an \emph{ordinary bounded operator} on $H=\ell^2(\N_{\ge1})$, the
map $De_n=n^{-1}e_n$ is compact, normal and injective, and its ordinary
spectrum is $\{1/n:n\ge1\}\cup\{0\}$ with $0$ as its only accumulation
point. \Cref{ihs:hh:cor:compact} then gives, for the failure-of-bijectivity
spectrum on $\Hil=\ell^2((t^\Gamma))$,
\[
 \sigma_\Gamma(D)=\mK\cup\{1/n:n\ge1\}.
\]
The whole infinitesimal monad of $0$ is in the spectrum. The witness is
concrete: a putative solution of $Dy=f$ with $f=(n^{-1})_n$ would need
leading coefficient $(1)_n$, which is not in $\ell^2$.
\item As an \emph{element of $\Arf$} with $I=\N_{\ge1}$ and $B=0$, the same
matrix $\diag(1/n)$ has $\diag(n)\in\RCF_I(\C)\subseteq\Arf$ as a two-sided
inverse, because a diagonal matrix of unbounded complex entries is a
perfectly good row- and column-finite matrix. Hence
$0\notin\sigma_{\Arf}(D)$, and \cref{ihs:rf:ex:accum} records this.
\end{itemize}
Both statements are true. They are not statements about the same operator:
one is about a bounded operator on a Hilbert space whose vectors are
square-summable, the other about a matrix acting on finite-support
coefficient vectors, and invertibility is being tested in two different
algebras. The unbounded diagonal inverse is excluded in Part I precisely
because \cref{ihs:hh:thm:adjoint} forces the coefficients of an adjointable
operator to be ordinary \emph{bounded} operators; it is admitted in Part II
precisely because row- and column-finiteness replaces boundedness there.
\Cref{ihs:sec:finitesections} makes the same point from the finite-section
side, and \cref{ihs:rf:prop:bounded} makes the converse point: the Part II
diagonalizer has a first coefficient with unbounded entries, so that theorem
does not secretly solve the ordinary bounded-operator problem.
\Cref{ihs:fr:thm:nounitary} of Part~IV makes the same point for an arbitrary
Hahn unitary in $\Ahh$ and a trace-class perturbation.

\subsection{One formula for the constant algebra-relative spectra}\label{ihs:sec:unify}
Part III's spectrum is defined for any unital algebra, and two of this
report's algebra-relative spectra are instances of it. The following
statement is an observation made in assembling this report: it follows from
the definitions and \cref{ihs:dz:thm:halo} and is not a theorem of any of
the four source manuscripts.

\begin{proposition}[The algebra-relative spectra of constant elements]\label{ihs:prop:unify}
\begin{enumerate}[label=\textup{(\roman*)}]
\item For $\Alg=\B$ one has $\AlgG=\Ahh$, so for every ordinary bounded $T$
\[
 \sigma^{\mathrm{adj}}_\Gamma(T)=\Sigma^{\B}_\Gamma(T)
 =\sigma_\C(T)\cup\bigl(\sigma_{\mathrm D,\B}(T)+\mK\bigr)
 \supseteq\sigma^{\mathrm{alg}}_\Gamma(T).
\]
For normal $T$ the inclusion is an equality by \cref{ihs:hh:thm:spectrum}.
\item For $\Alg=\RCF_I(\C)$ one has $\AlgG=\Arf$, so for every constant
row- and column-finite matrix $T$, not necessarily diagonal or Hermitian,
$\sigma_{\Arf}(T)=\Sigma^{\RCF_I(\C)}_\Gamma(T)$ is given by
\eqref{ihs:dz:eq:halo}.
\end{enumerate}
\end{proposition}
\begin{proof}
Both algebra identities are the definitions of $\Ahh$ and $\Arf$, and the
spectra are then defined by the same condition. The inclusion in (i) holds
because a two-sided inverse in $\Ahh$ acts as an inverse map on $\Hil$.
\end{proof}

Applied to the operator of \cref{ihs:sec:accum}, $d_n=1/n$, the one formula
\eqref{ihs:dz:eq:halo} produces both of that section's conclusions, and
shows that the difference lies entirely in which algebra is named.
\begin{itemize}
\item In $\Alg=\B$, $\sigma_\C(D)=\{1/n\}\cup\{0\}$ and $0$ is an
accumulation point, hence not a finite pole, so
$\sigma_{\mathrm D,\B}(D)=\{0\}$ by \eqref{ihs:dz:eq:dictionary}. Then
$\Sigma^{\B}_\Gamma(D)=\{1/n:n\ge1\}\cup\mK$, which is
\cref{ihs:hh:cor:compact}.
\item In $\Alg=\RCF_I(\C)$, $D-a$ is invertible for every
$a\notin\{1/n\}$, including $a=0$, because every diagonal matrix is row- and
column-finite; so $\sigma_{\RCF_I(\C)}(D)=\{1/n\}$. Each $D-1/m$ is Drazin
invertible of index one, with $P_{\mathrm{nil}}=E_{\{m\}}$ and the diagonal
Drazin inverse with entries $(1/n-1/m)^{-1}$ for $n\ne m$ and $0$ at $m$. So
the Drazin spectrum is empty and $\Sigma^{\RCF_I(\C)}_\Gamma(D)=\{1/n\}$,
which is \cref{ihs:rf:ex:accum}.
\end{itemize}
The same computation applies to any constant diagonal $D=\diag(d_i)$ in
$\RCF_I(\C)$, distinct labels or not, and gives $\{d_i:i\in I\}$; for
distinct real labels this is the case $B=0$ of
\cref{ihs:rf:thm:resolvent}. Part III does not treat the perturbed,
nonconstant operators $D+B$ of Part~II, nor the nonconstant operators of
Part~IV. Its general theorem classifies algebra inversion; its Volterra
and block examples additionally prove equality with the bijectivity spectrum
(\cref{ihs:dz:sec:volterra,ihs:dz:sec:blocks}). This report also proves
$\sigma^{\mathrm{alg}}_\Gamma=\sigma^{\mathrm{adj}}_\Gamma$ for constant
normal operators (\cref{ihs:hh:thm:spectrum}) and the compact-residue class
of \cref{ihs:fr:thm:main}. Separately, \cref{ihs:fr:thm:fredholm} equates
bijectivity and algebra inversion for $I+C$ with $C\in\TT_+$.

\subsection{Reading one part alone}\label{ihs:sec:readingalone}
Each part is self-contained given \cref{ihs:sec:prelim}, except that Part~IV
uses Part~I's space, algebra and range-defect theorem, and Part III's
comparison remarks refer to Parts I and II. A reader interested in the
Hilbert-space obstructions can read Part~I and stop; a reader interested in
exact infinite diagonalization can read Part~II and stop; a reader
interested in nonnormal constant operators in an arbitrary algebra can read
Part~III; a reader interested in noncommuting compact perturbations and
determinants can read Part~I's \cref{ihs:sec:construction,ihs:hh:sec:defect}
and then Part~IV. \Cref{ihs:sec:checks} describes the four verification
suites together, and
\cref{ihs:app:auditA,ihs:app:auditB,ihs:app:auditC,ihs:app:auditD} record
the separate repository and literature audits, which were carried out
against three different pinned revisions.

\clearpage
\part*{Part I. Hahn--Hilbert spaces: infinitesimal spectral thickening and closed-range defects}
\addcontentsline{toc}{part}{Part I. Hahn--Hilbert spaces}

\emph{Standing hypotheses for Part I.} $H\ne0$ is an ordinary complex
Hilbert space and $\Gamma$ is a nonzero set-sized ordered abelian group.
Divisibility is not assumed. Only \cref{ihs:hh:sec:rankone}, which uses a
real-valued metric, additionally assumes $\Gamma\subseteq\R$; the subsequent
closed-range and duality statements use the valuation topology at any rank. ``Spectrum'' in this part means failure
of bijectivity on $\Hil=H((t^\Gamma))$, or failure of invertibility in
$\Ahh$, and \cref{ihs:hh:thm:spectrum} proves that for constant normal
operators the two agree.

\section{Why infinite dimension is the next substantial question}\label{ihs:hh:sec:why}

This part asks what happens when an ordinary infinite-dimensional Hilbert
space is retained as the coefficient space, rather than replacing a finite
matrix by a larger finite matrix. Over $\C$, positivity, completeness and
adjoints support a familiar chain of implications: self-adjoint spectra are
real, coercive operators are invertible, and closed subspaces admit
orthogonal projections. The construction below retains a positive-definite
inner product, all ordinary Hilbert vectors as constants, and a large
adjointable operator algebra; in rank one it is complete. Nevertheless
several of those implications fail, for an exactly identifiable reason:
\emph{a leading coefficient must itself be an ordinary Hilbert vector}.

The central results are not a list of isolated anomalies. An automatic
operator-structure theorem makes the admissible algebra intrinsic. An
injective-leading-term theorem explains the persistence of range defects.
Together with ordinary normal spectral theory, these give an exact spectral
classification and quantitative counterexamples to coercive inversion.

\section{The Hahn--Hilbert space}\label{ihs:sec:construction}

\subsection{The field and the roots actually needed}\label{ihs:hh:sec:field}
Define
\[
 \Fld=\R((t^\Gamma)),\qquad \K=\C((t^\Gamma))=\Fld[i].
\]
Order $\Fld$ by its first nonzero coefficient and conjugate $\K$
coefficientwise. Let $\Ocal=\{z:v(z)\ge0\}$ and $\mK=\{z:v(z)>0\}$,
including zero in both sets. Coefficient extraction at zero gives
$\st:\Ocal\to\C$. Every finite scalar has a unique expression
$z=a+\eps$ with $a\in\C$ and $\eps\in\mK$.

The source manuscript assumed divisibility to obtain square roots of
all positive elements of $\Fld$. That assumption is stronger than the
proofs need. For a positive scalar
$r=ct^\delta(1+u)$ with $c>0$ and $v(u)>0$, a square root exists in $\Fld$
if and only if $\delta\in2\Gamma$. Necessity follows by taking valuations;
if $\delta=2\beta$, then
\[
 \sqrt r=\sqrt c\,t^\beta\sum_{n\ge0}\binom{1/2}{n}u^n.
\]
The support lemma and the binomial identity verify this positive root.
Squared vector lengths and scalar moduli always have even leading exponent,
as the next proposition proves. The roots of ordinary positive operators
used later are formed in the ordinary bounded-operator algebra. Thus no
Part I proof requires roots of arbitrary positive Hahn scalars, and the
standing divisibility restriction can be removed.

\subsection{The Hahn--Hilbert inner product}\label{ihs:hh:sec:inner}
We take ordinary and extended inner products to be linear in the first
variable and conjugate-linear in the second. On $\Hil=H((t^\Gamma))$ define
\begin{equation}\label{ihs:hh:eq:inner}
 \ip{x}{y}=\sum_{\alpha,\beta}\ip{x_\alpha}{y_\beta}_H t^{\alpha+\beta}.
\end{equation}
Each ordinary Hilbert inner product in the coefficient is evaluated
\emph{before} the Hahn convolution. For example, a constant vector
$(1/n)_{n\ge1}\in\ell^2$ has its ordinary finite real squared norm, even
though the separate family of its constant-coordinate squared magnitudes is
not Hahn summable. We are not using that inadmissible family as a Hahn sum.

\begin{proposition}[Positive-definite scalar extension]\label{ihs:hh:prop:inner}
For any ordered abelian group $\Gamma$, formula
\eqref{ihs:hh:eq:inner} defines a positive-definite Hermitian form.
Its norm $\norm{x}=\sqrt{\ip{x}{x}}$ exists in $\Fld$, satisfies the
Cauchy--Schwarz and triangle inequalities, and
\begin{equation}\label{ihs:hh:eq:normvaluation}
 v(\ip{x}{x})=2v(x),\qquad v(\norm{x})=v(x)\quad(x\ne0).
\end{equation}
If $x=t^\delta(u+\text{positive-order terms})$ with $u\ne0$, then
$\norm{x}=t^\delta(\norm{u}_H+\text{positive-order terms})$.
\end{proposition}
\begin{proof}
\Cref{ihs:lem:support}(i) gives well-definedness and sesquilinearity. The
coefficient at $2\delta$, where $\delta=v(x)$, is exactly
$\norm{x_\delta}_H^2>0$; all other contributions have larger exponent. This
proves positivity. For $x\ne0$ write
\[
 \ip{x}{x}=\norm{x_\delta}_H^2 t^{2\delta}(1+u),
 \qquad u=0\ \text{or}\ v(u)>0.
\]
The binomial root of \cref{ihs:cor:binomial} gives
\[
 \norm{x}=\norm{x_\delta}_H t^\delta
       \sum_{n\ge0}\binom{1/2}{n}u^n.
\]
This is a positive square root in $\Fld$ for every $\Gamma$; set
$\norm0=0$. Its leading term proves both valuation assertions.
The same argument for $H=\C$ defines the scalar modulus
$\abs z=\sqrt{z\bar z}$ without divisibility. Completing a square for
$y\ne0$ (the case $y=0$ is immediate), using
$\ip{x-ay}{x-ay}\ge0$ with $a=\ip{x}{y}/\ip{y}{y}$, gives
$\abs{\ip{x}{y}}^2\le\ip{x}{x}\ip{y}{y}$. The usual expansion of
$\norm{x+y}^2$, followed by this inequality, gives the triangle inequality.
These are finite algebraic arguments over the ordered field $\Fld$.
\end{proof}

\section{Every adjoint has a global Hahn expansion}\label{ihs:sec:adjoints}

\subsection{The apparently smaller coefficient algebra}\label{ihs:hh:sec:algebra}
Let $\Ahh=\B((t^\Gamma))$. Thus an element $A=\sum_\gamma A_\gamma t^\gamma$
has bounded ordinary operator coefficients and one well-ordered support.
There is no uniform ordinary bound on the norms of all its coefficients. It
acts by
\begin{equation}\label{ihs:hh:eq:operator-convolution}
 (Ax)_\delta=\sum_{\gamma+\alpha=\delta}A_\gamma x_\alpha.
\end{equation}
\Cref{ihs:lem:support} proves this action is defined on every $x\in\Hil$.
The series $A^*=\sum_\gamma A_\gamma^*t^\gamma$ satisfies
$\ip{Ax}{y}=\ip{x}{A^*y}$, by finite coefficientwise rearrangement.

One might suspect that $\Ahh$ is only a convenient restricted operator
class. The next theorem says that the existence of an everywhere-defined
adjoint forces exactly this class, without any advance assumption that an
operator preserves Hahn sums.

\begin{theorem}[Automatic adjoint structure]\label{ihs:hh:thm:adjoint}
Let $A,B:\Hil\to\Hil$ be everywhere-defined $\K$-linear maps such that
\begin{equation}\label{ihs:hh:eq:adjoint-condition}
 \ip{Ax}{y}=\ip{x}{By}\qquad(x,y\in\Hil).
\end{equation}
There are unique bounded operators $A_\gamma\in\B$ with well-ordered
nonzero support such that $A$ acts by
\eqref{ihs:hh:eq:operator-convolution}. Moreover $B_\gamma=A_\gamma^*$.
Consequently
\[
 \End_{\mathrm{adj}}(\Hil)\cong\B((t^\Gamma))
\]
as $*$-algebras. No separability of $H$ is assumed.
\end{theorem}
\begin{proof}
\emph{Step 1: bounded coefficient operators.}
For constant vectors $u,w\in H$, define $A_\gamma u=[t^\gamma]Au$ and
$B_\gamma w=[t^\gamma]Bw$. Extracting the coefficient at $\gamma$ in
\eqref{ihs:hh:eq:adjoint-condition} gives
\begin{equation}\label{ihs:hh:eq:ordinaryadj}
 \ip{A_\gamma u}{w}_H=\ip{u}{B_\gamma w}_H.
\end{equation}
Both maps are defined on all of $H$. To see that $A_\gamma$ has closed
graph, let $u_n\to u$ and $A_\gamma u_n\to z$ in the ordinary Hilbert norm.
Testing against any fixed $w$ in \eqref{ihs:hh:eq:ordinaryadj} gives
$\ip{z}{w}=\ip{u}{B_\gamma w}=\ip{A_\gamma u}{w}$. Hence $z=A_\gamma u$. The
ordinary closed graph theorem makes $A_\gamma$ bounded. The same argument
applies to $B_\gamma$, and \eqref{ihs:hh:eq:ordinaryadj} identifies the two
as adjoints.

\emph{Step 2: one well-ordered support.}
Let $S=\{\gamma:A_\gamma\ne0\}$. If $S$ were not well ordered, ordinary
choice permits a strictly decreasing sequence
$\gamma_1>\gamma_2>\cdots$ in $S$. Each $\ker A_{\gamma_n}$ is a proper
closed subspace of $H$, and hence nowhere dense. Baire's theorem gives
$u\in H$ outside their countable union. Then the support of $Au$ contains
every $\gamma_n$, contradicting its being a Hahn series. Thus $S$ is well
ordered. Because an ordinary bounded operator is zero exactly when its
adjoint is zero, the coefficient family of $B$ has the same support.

\emph{Step 3: constants determine the full operator.}
Let $\widetilde A,\widetilde B$ be the convolution operators built from
these coefficient families. They are adjoints by \cref{ihs:lem:support} and
agree with $A,B$ on constants. For arbitrary $x\in\Hil$ and constant
$w\in H$,
\[
 \ip{Ax}{w}=\ip{x}{Bw}=\ip{x}{\widetilde B w}=\ip{\widetilde A x}{w}.
\]
If a Hahn vector pairs to zero with every constant $w$, then each of its
ordinary coefficients pairs to zero with all of $H$, so every coefficient
vanishes. Hence $Ax=\widetilde A x$. The analogous conclusion holds for $B$.
Uniqueness follows by evaluating on constant vectors and extracting
coefficients. Composition and the involution are the finite Hahn convolution
operations, proving the algebra assertion.
\end{proof}

\begin{corollary}[Automatic summability and boundedness]\label{ihs:hh:cor:autobounded}
Every everywhere-defined adjointable operator preserves Hahn-summable
families. It also admits a bound
\begin{equation}\label{ihs:hh:eq:orderbound}
 \norm{Ax}\le M\norm{x}\qquad(x\in\Hil)
\end{equation}
for some $M\in\Fld_{\ge0}$. If $A\ne0$ and $\gamma_0=\min\supp A$, any
ordinary real $c>\norm{A_{\gamma_0}}_H$ gives the valid bound
$M=c t^{\gamma_0}$.
\end{corollary}
\begin{proof}
For a Hahn-summable family, the finiteness in \cref{ihs:lem:support}(i),
applied with $\supp A$, allows the two sums to be interchanged, and their
union of supports remains well ordered. For the bound, write the leading
term of $x$ as $t^\delta u$. At exponent $2(\gamma_0+\delta)$, the
coefficient of $M^2\norm{x}^2-\norm{Ax}^2$ is
$c^2\norm{u}_H^2-\norm{A_{\gamma_0}u}_H^2>0$. This remains true when
$A_{\gamma_0}u=0$, since then the leading output moves to a higher exponent.
Thus the difference is positive for nonzero $x$.
\end{proof}

\begin{remark}\label{ihs:hh:rem:baire}
The Baire argument is a uniform-support theorem, not a uniform bound on
$\norm{A_\gamma}_H$. The proof also explains why unrestricted ordinary
algebraic endomorphisms, used later for choosing range complements, are not
automatically adjointable.
\end{remark}

Endomorphisms preserving strong sums, without adjointability, are classified
as $\mathrm{End}_\C(H)((t^\Gamma))$ in \texttt{duals:thm:strong} of
\path{docs/surcomplex/three-duals-of-hahn-vector-spaces/}.

\section{An injective leading term preserves the whole range defect}\label{ihs:hh:sec:defect}

\subsection{A theorem independent of Hilbert structure}\label{ihs:hh:sec:defectthm}
The next result applies to any complex vector space $V$. Write
$V_\Gamma=V((t^\Gamma))$. An ordinary linear map on $V$ extends
coefficientwise, even if $V$ also carries a topology and the map is
discontinuous in that topology.

\begin{theorem}[Positive-order defect rigidity]\label{ihs:hh:thm:defect}
Let $S:V\to V$ be injective. Let $E\in\End_\C(V)((t^\Gamma))$ have strictly
positive order, or be zero, and put $A=S+E$. Choose an algebraic complement
\begin{equation}\label{ihs:hh:eq:complement}
 V=\Ran S\oplus W.
\end{equation}
Then
\begin{equation}\label{ihs:hh:eq:defectdecomp}
 V_\Gamma=A(V_\Gamma)\oplus W((t^\Gamma)).
\end{equation}
In particular, $A$ is injective, $v(Ax)=v(x)$ for $x\ne0$, and
$\coker A\cong W((t^\Gamma))$ as $\K$-vector spaces. The isomorphism depends
on the chosen complement, but the nonvanishing of the defect does not.
\end{theorem}
\begin{proof}
At the least exponent $\delta$ of a nonzero $x$, the coefficient of $Ax$ is
$Sx_\delta\ne0$. All contributions from $E$ have greater exponent. This
proves injectivity and the valuation assertion.

Define $U:V\to V$ by $U(Sv+w)=v$ for $v\in V$, $w\in W$. Then $US=I$ and
$\ker U=W$. Extend $U$ coefficientwise and use \cref{ihs:lem:neumann} to set
\[
 R=(I+UE)^{-1},\qquad L=RU,\qquad \Pi=I-AL.
\]
The coefficients of $U$ need not be bounded in any additional topology; all
these formulas are legitimate in the ordinary algebraic endomorphism Hahn
algebra. The identities
\[
 LA=R(US+UE)=R(I+UE)=I
\]
and
\[
 U\Pi=U-UA R U=U-(I+UE)R U=0
\]
show that $\Pi$ kills the range of $A$ and takes values in
$\ker U=W((t^\Gamma))$. On that latter space $L$ vanishes, so $\Pi$ is the
identity. Finally, every $x$ decomposes as $ALx+\Pi x$; the intersection of
the two summands is zero because $LA=I$ and $L|_{W((t^\Gamma))}=0$.
\end{proof}

\begin{corollary}[No positive-order repair]\label{ihs:hh:cor:norepair}
An injective nonsurjective ordinary operator cannot become surjective on
$V_\Gamma$ by adding any globally supported positive-order Hahn
perturbation. This includes perturbations that do not commute with it and
perturbations with infinitely many positive exponents.
\end{corollary}

\begin{corollary}[Explicit independent obstruction classes]\label{ihs:hh:cor:independent}
Suppose $(f_j)_{j\in J}\subset V$ have linearly independent classes in
$V/\Ran S$ over $\C$. Their constant Hahn extensions have linearly
independent classes in $V_\Gamma/(S+E)V_\Gamma$ over $\K$. Consequently
\[
 \dim_\K\coker(S+E)\ge\dim_\C(V/\Ran S).
\]
\end{corollary}
\begin{proof}
In a nonzero finite sum $f=\sum_j c_j f_j$ with $c_j\in\K$, take the least
exponent among the nonzero $c_j$. The coefficient at that exponent is a
nonzero finite complex combination of the $f_j$, and is not in $\Ran S$. But
the leading coefficient of $(S+E)x$ is always in $\Ran S$.
\end{proof}

\begin{warning}\label{ihs:hh:warn:dimension}
\Cref{ihs:hh:thm:defect} gives the actual vector-space quotient, not only a
first-order approximation to it. It does \emph{not} assert
$\dim_\K W((t^\Gamma))=\dim_\C W$ for arbitrary infinite-dimensional $W$.
The safe dimension statement is the lower bound just proved in
\cref{ihs:hh:cor:independent}.
\end{warning}

\subsection{Why the result is more informative than a determinant test}\label{ihs:hh:sec:nodet}
There is no finite determinant for an arbitrary infinite-dimensional $S$. An
entrywise infinitesimal perturbation may even make every finite section
invertible, while the leading ordinary coefficient of a proposed inverse
vector is excluded from $H$. The exact splitting
\eqref{ihs:hh:eq:defectdecomp} shows that this is not an accidental choice
of data: the full algebraic cokernel persists. On the other hand, if $S$ is
an ordinary bijection, $S+E=S(I+S^{-1}E)$ is invertible by the same
support-certified Neumann formula. Thus the injective case has a sharp
dichotomy, not merely a negative example.

\section{The exact spectrum of an ordinary normal operator}\label{ihs:hh:sec:spectrum}

\subsection{Two spectral definitions and one ordinary reduction}\label{ihs:hh:sec:twodefs}
For a $\K$-linear endomorphism $A$ of $\Hil$, define
\[
 \sigma_\Gamma^{\mathrm{alg}}(A)
 =\{\lambda\in\K:A-\lambda I\text{ is not bijective on }\Hil\}.
\]
For $A\in\Ahh$, define $\sigma_\Gamma^{\mathrm{adj}}(A)$ by failure of
invertibility in the unital algebra $\Ahh$. These definitions need not be
identified in advance. In the theorem below, explicit inverses and explicit
range obstructions will prove their equality. Neither is the
algebra-relative spectrum $\sigma_{\Arf}$ of Part~II, which is
taken in a different algebra acting on a different vector space. For a
constant $A$, $\sigma^{\mathrm{adj}}_\Gamma(A)$ is the case $\Alg=\B$ of
Part~III's spectrum $\SigA$ (\cref{ihs:prop:unify}).

An ordinary bounded operator $T$ is regarded as a constant member of
$\Ahh$. We often suppress the extension subscript on $T$.

\begin{lemma}[Normal isolated-point decomposition]\label{ihs:hh:lem:normalgap}
Let $T\in\B$ be normal and $a\in\sigma_\C(T)$. Put $N=\ker(T-aI)$,
$M=N^\perp$, and $S=(T-aI)|_M$. Then $M,N$ reduce $T$ and $S$ is injective.
If $a$ is isolated in $\sigma_\C(T)$, then $N\ne0$ and $S$ is boundedly
invertible on $M$. If $a$ is an accumulation point, then $S$ is not
surjective.
\end{lemma}
\begin{proof}
These facts follow from the ordinary normal spectral theorem
\cite{Douglas}; here is the precise use of it. The spectral projection of
$\{a\}$ is the orthogonal projection onto $N$. For isolated $a$, the
function $(z-a)^{-1}$ on the rest of the compact spectrum is bounded, so
spectral functional calculus gives an inverse on $M$. The projection cannot
be zero, since otherwise $a$ would not belong to the spectrum. For
accumulation $a$, if $S$ were onto, its injectivity and the bounded inverse
theorem would give a bounded inverse. The spectrum of $T$ would then be
contained in $\{a\}$ together with a compact set separated from $a$,
contradicting accumulation. This also covers $N=0$.
\end{proof}

\subsection{Classification and all inverse formulas}\label{ihs:hh:sec:classif}
\begin{theorem}[Accumulation-monad spectrum]\label{ihs:hh:thm:spectrum}
For every ordinary bounded normal $T\in\B$,
\begin{equation}\label{ihs:hh:eq:exactspectrum}
 \sigma_\Gamma^{\mathrm{alg}}(T)
 =\sigma_\Gamma^{\mathrm{adj}}(T)
 =\sigma_\C(T)\cup\bigl(\sigma_\C(T)'+\mK\bigr).
\end{equation}
No scalar of negative valuation is in this spectrum. A nonzero infinitesimal
displacement of an isolated ordinary spectral point is in the resolvent,
whereas every infinitesimal displacement of an accumulation point remains in
the spectrum.
\end{theorem}
\begin{proof}
We partition all scalars $\lambda\in\K$ into exhaustive cases.

\emph{Infinite scalar.} If $v(\lambda)<0$, then
\begin{equation}\label{ihs:hh:eq:infiniteres}
 (T-\lambda I)^{-1}=-\sum_{n\ge0}\lambda^{-n-1}T^n.
\end{equation}
This is a support-certified Neumann inverse: the scalar $\lambda^{-1}$ has
positive valuation. It belongs to $\Ahh$, irrespective of the ordinary
growth of $\norm{T^n}$.

\emph{Finite scalar with nonspectral residue.} Write $\lambda=a+\eps$, where
$a\in\C$ and $\eps\in\mK$. If $a\notin\sigma_\C(T)$, let
$B=(T-aI)^{-1}\in\B$. Then
\begin{equation}\label{ihs:hh:eq:regularres}
 (T-aI-\eps I)^{-1}=\sum_{n\ge0}\eps^n B^{n+1}.
\end{equation}
The formula includes $\eps=0$. For nonzero $\eps$, use
\cref{ihs:lem:neumann} on $\eps B$; coefficients need not satisfy any
ordinary norm-convergence estimate.

\emph{Undisplaced ordinary spectral point.} If $\eps=0$ and
$a\in\sigma_\C(T)$, then $T-aI$ is not an ordinary bijection, by the bounded
inverse theorem. A nonzero ordinary kernel vector remains a Hahn kernel
vector. If there is no kernel but surjectivity fails, choose a constant $f$
outside its ordinary range. The coefficient at zero in any putative solution
would have to solve $(T-aI)y_0=f$, which is impossible. Thus the Hahn
operator is not bijective.

\emph{Displaced isolated point.} Suppose $a$ is isolated and $\eps\ne0$. Use
\cref{ihs:hh:lem:normalgap}; let $P,Q$ project onto $N,M$. The inverse is
\begin{equation}\label{ihs:hh:eq:isolatedres}
 (T-aI-\eps I)^{-1}=-\eps^{-1}P+\sum_{n\ge0}\eps^n S^{-n-1}Q.
\end{equation}
The first term is a single scalar Hahn series times a bounded projection;
the second is a positive-order Neumann series on $M$. Both have well-ordered
operator support and lie in $\Ahh$. The finite orthogonal splitting verifies
the inverse on both summands. If $M=0$, the second summand is simply zero.

\emph{Displaced accumulation point.} Suppose $a\in\sigma_\C(T)'$ and
$\eps\ne0$. On $M$ in \cref{ihs:hh:lem:normalgap}, the ordinary $S$ is
injective and not onto. The perturbation $-\eps I$ has strictly positive
order. By \cref{ihs:hh:thm:defect}, $S-\eps I$ is not onto
$M((t^\Gamma))$. The constant decomposition $H=N\oplus M$ extends to a
decomposition of Hahn spaces, and the full operator is block diagonal. It is
therefore not onto either. Explicitly, a constant $f\in M\setminus\Ran S$
cannot be its $M$-component output.

These cases prove exactly the stated set formula. In every nonspectral case
we constructed an inverse in $\Ahh$; in every spectral case we proved
failure of algebraic bijectivity. This proves equality of the two spectral
definitions as well.
\end{proof}

\subsection{Point spectrum is rigid, despite spectral thickening}\label{ihs:hh:sec:point}
\begin{theorem}[No new normal eigenvalues]\label{ihs:hh:thm:point}
For ordinary bounded normal $T$,
\[
 \sigma_{p,\Gamma}(T)=\sigma_{p,\C}(T),\qquad
 \ker(T-aI)=\ker_H(T-aI)((t^\Gamma))\quad(a\in\C).
\]
In particular, the new nonreal spectral points of a self-adjoint operator
are not eigenvalues.
\end{theorem}
\begin{proof}
Infinite scalars are in the resolvent by \cref{ihs:hh:thm:spectrum}. Suppose
$(T-aI-\eps I)x=0$ with $x\ne0$ and $\eps\in\mK$. At its leading exponent
$\delta$, one obtains $(T-aI)x_\delta=0$. Thus $a$ is an ordinary eigenvalue
and $x_\delta\in N=\ker_H(T-aI)$. If $\eps\ne0$, apply the ordinary
orthogonal projection $P$ onto $N$. Normality implies $P(T-aI)=0$, so the
full equation gives $-\eps Px=0$. Hence $Px=0$, contradicting
$Px_\delta=x_\delta\ne0$. Therefore $\eps=0$. For a constant $a$, the kernel
equation is coefficientwise, yielding the claimed eigenspace identity.
\end{proof}

\begin{example}[Normality is indispensable]\label{ihs:hh:ex:backshift}
Let $B$ be the backward shift on $\ell^2(\N_0)$: $Be_0=0$, $Be_n=e_{n-1}$
for $n\ge1$. For $s=t^\eta$ with $\eta>0$,
\[
 x=\sum_{n\ge0}s^n e_n\in\Hil,\qquad Bx=sx.
\]
The support $\{n\eta:n\ge0\}$ is well ordered. This produces a new
nonconstant eigenvalue for the nonnormal operator $B$, in direct contrast to
\cref{ihs:hh:thm:point}. No ordinary Hilbert convergence of the vector
series is asserted or needed.
\end{example}

\begin{remark}[The nonnormal constant case]\label{ihs:hh:rem:drazin}
Part~III computes $\sigma^{\mathrm{adj}}_\Gamma(T)$ for \emph{every}
ordinary bounded $T$, normal or not, and for constant elements of any unital
algebra: the accumulation set is replaced by the finite-index Drazin
spectrum (\cref{ihs:dz:thm:halo,ihs:dz:cor:normal}). Its proof is a second
route to the $\sigma^{\mathrm{adj}}$ half of \eqref{ihs:hh:eq:exactspectrum}
(\cref{ihs:dz:rem:recover}). It does not replace
\cref{ihs:hh:thm:spectrum}, which also determines
$\sigma^{\mathrm{alg}}_\Gamma$, nor \cref{ihs:hh:thm:point}; and it concerns
the spectral formula, not the new eigenvalue of the example above.
\end{remark}

\section{Consequences that separate finite and infinite spectral geometry}\label{ihs:hh:sec:consequences}

\subsection{Self-adjoint spectra are real exactly in the finite-spectrum case}\label{ihs:hh:sec:reality}
\begin{corollary}[A sharp spectral-reality criterion]\label{ihs:hh:cor:reality}
Let $T$ be ordinary bounded self-adjoint. The following are equivalent:
\begin{enumerate}[label=\textup{(\roman*)}]
\item $\sigma_\Gamma(T)\subseteq\Fld$;
\item $\sigma_\C(T)$ is finite;
\item $T$ has a finite ordinary orthogonal spectral resolution
$T=\sum_{j=1}^{r}a_jP_j$ with $a_j\in\R$.
\end{enumerate}
There is no finite-rank hypothesis on the projections $P_j$.
\end{corollary}
\begin{proof}
The usual self-adjoint spectrum is a nonempty compact subset of $\R$. A
finite set has no accumulation points, so \cref{ihs:hh:thm:spectrum} proves
(ii)$\Rightarrow$(i). An infinite compact subset has an accumulation point
$a$. For any $s=t^\eta$, $\eta>0$, the scalar $a+is$ belongs to
$\sigma_\Gamma(T)$ and does not belong to $\Fld$, proving the converse. The
equivalence with (iii) is the ordinary spectral theorem applied to a finite
spectral set.
\end{proof}

The criterion is finite \emph{spectrum}, not finite dimension. For example,
an infinite-rank orthogonal projection with infinite-dimensional kernel has
spectrum exactly $\{0,1\}$ even after Hahn extension. Infinite multiplicity
at an isolated point does not create a monad. The accumulation set in
\eqref{ihs:hh:eq:exactspectrum} therefore cannot be replaced by an essential
spectrum that includes such points.

\subsection{Ordinary compact operators and a continuous-spectrum example}\label{ihs:hh:sec:compact}
\begin{corollary}[Ordinary compact normal extension]\label{ihs:hh:cor:compact}
If $T$ is an ordinary compact normal operator of infinite rank, then
\[
 \sigma_\Gamma(T)=\mK\cup\{\text{nonzero ordinary eigenvalues of }T\}.
\]
If its ordinary spectrum is finite, there is no thickening.
\end{corollary}
\begin{proof}
The ordinary compact normal spectral theorem says that nonzero spectral
points are isolated eigenvalues of finite multiplicity and zero is the only
possible accumulation point. Infinite rank forces zero to be an accumulation
point. Apply \cref{ihs:hh:thm:spectrum}.
\end{proof}

For an injective compact self-adjoint $T$, Part~IV adds a self-adjoint Hahn
perturbation of positive order that need not commute with $T$: the monad
$\mK$ stays in the spectrum, and each nonzero eigenvalue is replaced by a
finite cluster of real Hahn eigenvalues (\cref{ihs:fr:thm:main}, under a
divisibility hypothesis).

Here ``compact'' describes the original operator on the ordinary Hilbert
space. It makes no claim of compactness or compactoidness in the Hahn
valuation topology. The sequence $(De_n)$ for $De_n=n^{-1}e_n$, for
instance, satisfies $v(De_n-De_m)=0$ for $n\ne m$, despite
$\norm{De_n}_H\to0$ in the ordinary norm. Thus the valuation ball
$\{x:v(x)>0\}$ separates these outputs at every rank; in the rank-one
metric below their pairwise distance is one.

For another concrete example, let $T$ be multiplication by the coordinate on
$L^2([0,1],\C)$. Its ordinary spectrum is $[0,1]$, it has no eigenvalues,
and every point of its spectrum is an accumulation point. Thus
\[
 \sigma_\Gamma(T)=[0,1]+\mK,\qquad \sigma_{p,\Gamma}(T)=\varnothing.
\]
The usual facts here can be seen directly: away from $[0,1]$ the reciprocal
multiplier is bounded; near any point of $[0,1]$, vectors supported in
arbitrarily small intervals prevent a bounded inverse; a singleton has
measure zero and cannot support a nonzero eigenvector. This gives full
infinitesimal fibers above a continuous real interval, without adding a
single eigenvector.

\subsection{Enlarging the scalar workspace does not repair the defect}\label{ihs:hh:sec:extension}
\begin{corollary}[Persistence under ordered-group extension]\label{ihs:hh:cor:extension}
Let $\Gamma\subseteq\Delta$ be nonzero ordered abelian groups and use the natural
coefficient-preserving embedding $K_\Gamma\subseteq K_\Delta$. For constant
normal $T$,
\[
 \sigma_\Delta(T)\cap K_\Gamma=\sigma_\Gamma(T).
\]
Moreover, the nonsurjectivity in \cref{ihs:hh:cor:norepair} persists after
such an extension for the same $S$ and $E$.
\end{corollary}
\begin{proof}
An element of $K_\Gamma$ has the same leading valuation, residue and
infinitesimal status in the larger group. Apply the exact spectral formula.
For the second assertion use the same complement $W$ in
\cref{ihs:hh:thm:defect}, now forming $W((t^\Delta))$.
\end{proof}

This is a genuine surcomplex-workspace statement. Adding further surreal
scales does not turn a prohibited leading ordinary Hilbert coefficient into
an allowed one. It does not assert that an unspecified larger class-sized
notion of vector could never behave differently.
\Cref{ihs:rf:prop:extension} is the corresponding naturality statement for
Part II; it is about a different construction and is proved separately.

\section{Which constant data have an infinitesimal resolvent?}\label{ihs:hh:sec:resolventdata}

\subsection{The inverse-smooth domain}\label{ihs:hh:sec:smooth}
Even when $S-sI$ is not onto, some constant right-hand sides have a unique
solution. The exact criterion is stronger than belonging merely to $\Ran S$.

\begin{theorem}[All inverse moments are necessary and sufficient]\label{ihs:hh:thm:smooth}
Let $S:V\to V$ be injective and $s=t^\eta$ with $\eta>0$. For constant
$f\in V$, the equation
\begin{equation}\label{ihs:hh:eq:constres}
 (S-sI)y=f,\qquad y\in V((t^\Gamma)),
\end{equation}
has a solution if and only if
\begin{equation}\label{ihs:hh:eq:inversesmooth}
 f\in\bigcap_{k\ge1}\Ran S^k.
\end{equation}
In that case the solution is unique and is
\begin{equation}\label{ihs:hh:eq:smoothsolution}
 y=\sum_{n\ge0}s^n S^{-n-1}f.
\end{equation}
The inverse powers denote the unique preimages, not an everywhere-defined
inverse operator on $V$.
\end{theorem}
\begin{proof}
Sufficiency follows by coefficientwise telescoping of
\eqref{ihs:hh:eq:smoothsolution}; its support is contained in the
well-ordered set $\{n\eta:n\ge0\}$. Injectivity follows from
\cref{ihs:hh:thm:defect}.

For necessity, partition the support of a proposed solution according to
cosets of $\Z\eta$ in $\Gamma$. Restricting a Hahn series to a coset
preserves well-ordering, and both $S$ and multiplication by $s$ preserve
that coset. On a coset other than $\Z\eta$, the restricted equation is
homogeneous. Its least nonzero coefficient would be killed by $S$, which is
impossible. Thus the support lies in $\Z\eta$.

A well-ordered subset of $\Z\eta$ is bounded below in its integer index. For
$f\ne0$, the valuation-preservation assertion of \cref{ihs:hh:thm:defect}
forces the least index to be zero. Write $y=\sum_{n\ge0}s^n y_n$, allowing
zero coefficients. The equation gives
\[
 Sy_0=f,\qquad Sy_n=y_{n-1}\quad(n\ge1).
\]
Consequently $S^{n+1}y_n=f$ for every $n$, proving
\eqref{ihs:hh:eq:inversesmooth} and the formula. The zero-data case is
immediate.
\end{proof}

\begin{example}[The diagonal inverse-smooth space]\label{ihs:hh:ex:smoothD}
For $De_n=n^{-1}e_n$ on $H=\ell^2(\N_{\ge1})$, the admissible constant data
are exactly
\begin{equation}\label{ihs:hh:eq:rapidspace}
 H_\infty=\left\{f=(f_n):
       \sum_{n\ge1}n^{2k}|f_n|^2<\infty
       \text{ for every integer }k\ge1\right\}.
\end{equation}
The formula follows because $D^{-k}f=(n^k f_n)$. Finitely supported vectors
belong to $H_\infty$, so this subspace is dense in the \emph{ordinary}
Hilbert topology. Yet $f_n=n^{-2}$ does not belong to it, although
$f\in\Ran D$. Thus even some data with an ordinary first preimage lack an
infinitesimal resolvent.
\end{example}

The same statement for $D-isI$ follows by replacing the recurrence with
$Dy_n=i y_{n-1}$. Its solution, when it exists, is
$\sum_{n\ge0}i^n s^nD^{-n-1}f$. No uniform bound on the ordinary norms of
these coefficients is required, only that each one actually belongs to $H$.

\subsection{Why individual scalar inverses are insufficient}\label{ihs:hh:sec:scalarinverses}
Every diagonal scalar $n^{-1}-s$ is invertible in $\K$ and has expansion
\[
 (n^{-1}-s)^{-1}=\sum_{k\ge0}n^{k+1}s^k.
\]
But applying these inverses coordinatewise to $f$ requires the entire
coefficient vector $(n^{k+1}f_n)_n$ to lie in $\ell^2$ at each $k$. The Hahn
vector space imposes this condition before any comparison of higher scales.
There is therefore no contradiction between invertibility of every diagonal
entry and failure of surjectivity of the infinite operator. A putative
coordinatewise inverse has left the chosen vector space, rather than found a
missing scalar inverse. In Part~II, where the vector space and
the operator algebra are different, the corresponding coordinatewise inverse
is admissible; \cref{ihs:sec:accum} compares the two situations directly.

\section{Coercive operators with continuum-dimensional defects}\label{ihs:hh:sec:coercive}

\subsection{An explicit large ordinary cokernel}\label{ihs:hh:sec:largecoker}
Continue with $H=\ell^2(\N_{\ge1})$ and $De_n=n^{-1}e_n$. The operator $D$ is
bounded, positive, self-adjoint and injective. It is not onto because
$(n^{-1})_n\in\ell^2$ has no preimage in $\ell^2$.

\begin{lemma}[Continuum many independent defect vectors]\label{ihs:hh:lem:continuum}
For $r\in(1/2,3/2]$, let $f^{(r)}=(n^{-r})_{n\ge1}\in H$. Their classes in
$H/\Ran D$ are linearly independent over $\C$.
\end{lemma}
\begin{proof}
Take distinct $r_1<\cdots<r_k$ and complex coefficients with $c_1\ne0$. For
$g=\sum_{j=1}^k c_j f^{(r_j)}$, one has
\[
 n g_n=c_1 n^{1-r_1}\bigl(1+o(1)\bigr).
\]
Its squared magnitude is bounded below by a positive real multiple of
$n^{2-2r_1}$ for sufficiently large $n$. Since $r_1\le3/2$, the exponent is
at least $-1$, and the series of those squared magnitudes diverges. Thus
$(n g_n)\notin\ell^2$, which is exactly $g\notin\Ran D$. There are continuum
many choices of $r$.
\end{proof}

\begin{theorem}[Every positive-order perturbation retains a large defect]\label{ihs:hh:thm:continuum}
For every $E\in\B((t^\Gamma))$ of positive order, including zero,
\[
 \dim_\K\coker(D+E)\ge2^{\aleph_0}.
\]
The same lower bound holds with leading operator $D^2$.
\end{theorem}
\begin{proof}
Apply \cref{ihs:hh:cor:independent} and \cref{ihs:hh:lem:continuum}. Since
$\Ran D^2\subseteq\Ran D$, the same vectors remain independent modulo
$\Ran D^2$.
\end{proof}

\subsection{A nonreal spectral point with a strong lower bound}\label{ihs:hh:sec:lowerbound}
For $s=t^\eta$, $\eta>0$, set $A=D-isI$. Since $D=D^*$,
\begin{equation}\label{ihs:hh:eq:lowerA}
 A^*A=D^2+s^2I,\qquad
 \norm{Ax}^2=\norm{Dx}^2+s^2\norm{x}^2\ge s^2\norm{x}^2.
\end{equation}
Nevertheless, $A$ is not onto by \cref{ihs:hh:thm:defect}. Thus the usual
lower-bound argument for excluding nonreal self-adjoint spectrum stops
precisely at surjectivity, not at injectivity.

The operator even has an attained upper norm bound:
\begin{equation}\label{ihs:hh:eq:normA}
 \norm{Ax}\le\sqrt{1+s^2}\,\norm{x},\qquad
 \norm{Ae_1}=\sqrt{1+s^2}.
\end{equation}
Indeed $D$ is an ordinary contraction and its extension remains a
contraction; for this diagonal operator the inequality also follows by
writing $I-D^2=R^*R$ with ordinary bounded diagonal
$Re_n=\sqrt{1-n^{-2}}e_n$. Equations
\eqref{ihs:hh:eq:lowerA}--\eqref{ihs:hh:eq:normA} show that neither lack of
a bound nor lack of an attained norm explains the failure of inversion.

\subsection{Positive coercivity does not imply surjectivity}\label{ihs:hh:sec:coercivethm}
\begin{theorem}[A coercive noninvertible self-adjoint operator]\label{ihs:hh:thm:coercive}
Let
\begin{equation}\label{ihs:hh:eq:C}
 C=D^2+s^2I.
\end{equation}
Then $C$ is everywhere-defined, adjointable and self-adjoint, with
\begin{align}
 \ip{Cx}{x}&\ge s^2\ip{x}{x},\label{ihs:hh:eq:coercive}\\
 \norm{Cx}&\le(1+s^2)\norm{x}.\label{ihs:hh:eq:Cbound}
\end{align}
The upper bound is attained at $e_1$. The operator is injective but not
surjective, and its cokernel has dimension at least $2^{\aleph_0}$ over
$\K$. Specifically, the constant vector $f=(n^{-2})_{n\ge1}$ is not in its
range.
\end{theorem}
\begin{proof}
Self-adjointness and \eqref{ihs:hh:eq:coercive} follow from
$\ip{D^2x}{x}=\ip{Dx}{Dx}\ge0$. The contraction property of $D^2$, the
triangle inequality and $Ce_1=(1+s^2)e_1$ give \eqref{ihs:hh:eq:Cbound} and
its attainment. The leading coefficient operator $D^2$ is injective, so
\cref{ihs:hh:thm:defect} gives injectivity, and
\cref{ihs:hh:thm:continuum} gives the cokernel bound.

For a direct witness, if $Cy=f$, valuation preservation forces $v(y)=0$. The
coefficient at zero must solve $D^2y_0=f$, that is, $(y_0)_n=1$ for every
$n$. This is not an $\ell^2$ vector. Thus no such $y$ exists.
\end{proof}

For clarity, coercivity is a bound by a strictly positive element $s^2$ of
the ordered scalar field. It is not a bound by a positive \emph{ordinary
real} independent of the infinitesimal scale. If instead $C\ge cI$ with
ordinary $c>0$ and an ordinary bounded positive leading operator with this
bound, the classical inverse and the positive-order Neumann construction
would apply. The infinitesimal scale is exactly where the distinction
becomes mathematically significant.

\begin{corollary}[Failure of a specified Lax--Milgram analogue]\label{ihs:hh:cor:lax}
On this Hahn inner-product space, a bounded coercive sesquilinear form need
not solve all variational problems, even when the right-hand side is
represented by an actual vector. In particular,
\[
 b(u,v)=\ip{Cu}{v},\qquad \ell(v)=\ip{f}{v},\qquad f=(n^{-2})_n,
\]
defines a bounded coercive form and a bounded conjugate-linear functional,
but there is no $u$ satisfying $b(u,v)=\ell(v)$ for all $v\in\Hil$.
\end{corollary}
\begin{proof}
Cauchy--Schwarz and \eqref{ihs:hh:eq:Cbound} bound the form;
\eqref{ihs:hh:eq:coercive} is coercivity, and
$\abs{\ell(v)}\le\norm{f}\norm{v}$. If the identity held for all $v$,
nondegeneracy of the inner product would give $Cu=f$, contradicting
\cref{ihs:hh:thm:coercive}.
\end{proof}

This does not contradict the ordinary complex Lax--Milgram theorem: the
scalar field, the topology and the duality structure have changed. It is an
explicit obstruction to transferring that theorem using only its
familiar-looking algebraic inequalities.

\section{Completeness does not restore orthogonal projection}\label{ihs:sec:topology}

\subsection{The rank-one topology and its completeness}\label{ihs:hh:sec:rankone}
For this subsection only, assume $\Gamma$ is a nonzero ordered
subgroup of $\R$. For $x\ne0$ define
\[
 \rho(x)=e^{-v(x)},\qquad \rho(0)=0,\qquad d(x,y)=\rho(x-y).
\]
The scalar absolute value here is $|a|_v=e^{-v(a)}$, not the $\Fld$-valued
modulus. The metric is ultrametric. Its topology agrees with the topology of
the field-valued norm balls on $\Hil$: a sufficiently higher valuation makes
a vector smaller than any prescribed positive radius, and choosing a still
higher monomial radius gives the converse inclusion of neighborhood bases.

It is not the ordinary Hilbert topology on the constant copy of $H$. Every
nonzero constant vector has $\rho$-size one. Nor is it the fine subspace
topology from the full surcomplex class, where the set-sized scalar
workspace is discrete \cite{RepoMap}.

\begin{proposition}[Rank-one valuation completeness]\label{ihs:hh:prop:complete}
For any complex vector space $V$, the metric space $V((t^\Gamma))$ with
metric $d$ is complete. In particular, $\Hil$ is complete in this topology.
\end{proposition}
\begin{proof}
Let $(x_n)$ be Cauchy. For each real cutoff $R$, all sufficiently late $x_n$
have exactly the same coefficients at exponents $\gamma\le R$, since
eventually $v(x_n-x_m)>R$. Define $x_\gamma$ to be the eventually stable
coefficient.

The support of this coefficient family is well ordered. Otherwise it would
contain a strictly decreasing sequence $\gamma_1>\gamma_2>\cdots$. Choose a
cutoff at least $\gamma_1$. Below that cutoff the entire stable family
agrees with one sufficiently late Hahn vector $x_N$, whose well-ordered
support cannot contain that sequence. This is a contradiction. Thus
$x=\sum x_\gamma t^\gamma$ belongs to $V((t^\Gamma))$. By the same
stabilization statement, $v(x_n-x)>R$ eventually for every real $R$, proving
convergence. Notice that no coefficient-space topology was used:
stabilization is exact, not approximate.
\end{proof}

\subsection{A proper closed subspace with zero orthogonal complement}\label{ihs:hh:sec:closedrange}
Return to an arbitrary nonzero ordered abelian group $\Gamma$. The valuation
topology has neighborhoods $\{x:v(x)>\gamma\}$, $\gamma\in\Gamma$.
These define the same topology as field-valued norm balls by the leading
norm formula. No real-valued metric or completeness hypothesis is needed
for the next two results.

\begin{theorem}[Closed-range orthogonality failure]\label{ihs:hh:thm:closedrange}
For the coercive $C=D^2+s^2I$, its range $M=C\Hil$ is closed and proper in
the valuation topology at every rank, but
\begin{equation}\label{ihs:hh:eq:orthdefect}
 M^\perp=\{0\},\qquad M^{\perp\perp}=\Hil\ne M.
\end{equation}
In particular, $M$ is not the range of a self-adjoint idempotent.
\end{theorem}
\begin{proof}
Use the splitting in \cref{ihs:hh:thm:defect} with $S=D^2$ and $E=s^2I$.
Its projection $\Pi=I-C(I+UE)^{-1}U$ has kernel exactly $M$:
it kills $M$, and $\Pi x=0$ gives $x=C(I+UE)^{-1}Ux$.
Every coefficient of $\Pi$ has nonnegative exponent, since $U,S$ are
constant and $E$ has positive order. Therefore $v(\Pi x)\ge v(x)$,
so $\Pi$ is valuation-continuous. The valuation topology is Hausdorff;
hence $M=\Pi^{-1}(\{0\})$ is closed. It is proper by
\cref{ihs:hh:thm:coercive}.

In rank one there is also the original metric proof:
$v(Cx)=v(x)$ makes $C$ an isometry, and its range is complete and closed
by \cref{ihs:hh:prop:complete}. The projection proof explains why rank one
is unnecessary for closedness.

If $z\in M^\perp$, then $\ip{Cx}{z}=0$ for every $x$. Since $C=C^*$,
$\ip{x}{Cz}=0$ for every $x$, so $Cz=0$. Injectivity gives $z=0$. This
proves \eqref{ihs:hh:eq:orthdefect}. A self-adjoint idempotent $P$ has
$\Hil=\Ran P\oplus\ker P$ and $\ker P=(\Ran P)^\perp$, by the finite
algebraic projection identities. If its range were $M$, the complement would
be zero and the range all of $\Hil$, a contradiction.
\end{proof}

The displayed range is not orthogonally closed, although it is closed in the valuation
topology. Thus the usual Hilbert identification of these two closure notions
fails. Known non-Archimedean failures of orthomodularity also warn against
unconditional projection theorems \cite{AN,ANS}. The theorem above gives a
particularly rigid witness: it is the closed range of a bounded, coercive,
self-adjoint injection and retains a continuum-sized algebraic deficit.

\subsection{A continuous functional outside inner-product duality}\label{ihs:hh:sec:riesz}
\begin{proposition}[Explicit failure of valuation Riesz representation]\label{ihs:hh:prop:riesz}
There is a nonzero $\K$-linear, valuation-continuous functional
$F:\Hil\to\K$ that vanishes on $M=C\Hil$ but is not of the form
$F(x)=\ip{x}{z}$ for any $z\in\Hil$.
\end{proposition}
\begin{proof}
Choose an algebraic complement $W$ of $\Ran D^2$ in $H$, and take the
projection $\Pi$ in \cref{ihs:hh:thm:defect} for $S=D^2$, $E=s^2I$. Its
coefficients have nonnegative exponents: $U$ and $S$ are constant, $E$ has
positive order, and $(I+UE)^{-1}$ has nonnegative order. Consequently
$v(\Pi x)\ge v(x)$. Choose a nonzero complex linear functional
$\psi:W\to\C$, and extend it coefficientwise to
$\psi_\Gamma:W((t^\Gamma))\to\K$. Set $F=\psi_\Gamma\Pi$. Coefficient
extraction gives $v(Fx)\ge v(x)$, so $F$ is valuation-continuous. It is
nonzero because $\Pi$ is the identity on $W((t^\Gamma))$, and it kills $M$.
If $F(x)=\ip{x}{z}$, then vanishing on $M$ implies $z\in M^\perp=0$, which
would force $F=0$.
\end{proof}

This functional is strong in the sense of the report
\path{docs/surcomplex/three-duals-of-hahn-vector-spaces/}
(\texttt{duals:thm:strong}); its leading coefficient is discontinuous, so it
lies in the first gap of \texttt{duals:thm:riesz}
(\texttt{duals:rem:ihs-riesz}).

The algebraic maps $U$ and $\psi$ need not be continuous for the ordinary
Hilbert norm. That does not prevent the displayed Hahn maps from being
valuation-continuous. Confusing those topologies would incorrectly apply an
ordinary closed graph or Riesz argument. Conversely, this proposition does
not contradict \cref{ihs:hh:thm:adjoint}: the functional just constructed
has no adjoint represented by a Hahn vector.

\section{When does an operator have a least field-valued norm?}\label{ihs:hh:sec:norm}

\subsection{Bounds versus the existence of their least element}\label{ihs:hh:sec:bounds}
For $A\in\Ahh$, let
\[
 \mathcal U(A)=\{M\in\Fld_{\ge0}:
             \norm{Ax}\le M\norm{x}\text{ for all }x\in\Hil\}.
\]
\Cref{ihs:hh:cor:autobounded} makes this set nonempty. By a field-valued
\emph{operator norm} we mean a least element of this set, not an arbitrary
chosen bound and not the valuation norm of the operator.

\begin{lemma}[Ordinary contraction bounds survive]\label{ihs:hh:lem:contract}
If $T\in\B$ and $n=\norm{T}_H$, then $\norm{Tx}\le n\norm{x}$ for every
$x\in\Hil$. For nonzero $x$ with leading coefficient $u$, the finite scalar
$\norm{Tx}/\norm{x}$ has standard part
\[
 \st\!\left(\frac{\norm{Tx}}{\norm{x}}\right)
       =\frac{\norm{Tu}_H}{\norm{u}_H}.
\]
The right side is zero if $Tu=0$.
\end{lemma}
\begin{proof}
The ordinary positive operator $n^2I-T^*T$ has a bounded positive square
root $R$ by ordinary Hilbert spectral theory. The coefficientwise identity
$R^*R=n^2I-T^*T$ implies $n^2\norm{x}^2-\norm{Tx}^2=\norm{Rx}^2\ge0$. For
the standard part, use the leading norm formula in
\cref{ihs:hh:prop:inner}. If $Tu=0$, the output has strictly larger
valuation or is zero, so the quotient has residue zero.
\end{proof}

\begin{theorem}[Norm-attainment criterion]\label{ihs:hh:thm:norm}
For an ordinary bounded operator $T$ on nonzero $H$, the set $\mathcal U(T)$
has a least element if and only if $T$ attains its ordinary operator norm on
a unit vector. When it does, the least element is $\norm{T}_H$ and is
attained by a constant unit vector in $\Hil$.
\end{theorem}
\begin{proof}
Put $n=\norm{T}_H$. If $T$ attains $n$, the constant maximizing vector
forces every upper bound to be at least $n$, while
\cref{ihs:hh:lem:contract} proves that $n$ is a bound. This includes $T=0$,
which attains its norm on every unit vector.

Now suppose $T$ does not attain its norm; then $n>0$. For each nonzero
$x\in\Hil$, \cref{ihs:hh:lem:contract} gives
\[
 \st\!\left(\frac{\norm{Tx}}{\norm{x}}\right)<n.
\]
Therefore every positive finite $M$ with $\st(M)=n$ is an upper bound,
including $n-\eps$ for every positive infinitesimal $\eps$. Every finite $M$
with residue greater than $n$ and every positive infinite $M$ is a bound as
well.

No finite nonnegative $M$ with residue less than $n$ can be a bound: choose
an ordinary unit vector $u$ with $\norm{Tu}_H>\st(M)$ using the definition
of the ordinary supremum. Its constant extension violates that bound. We
have thus described all upper bounds. A bound with residue greater than $n$,
or an infinite bound, can be decreased to $n$. A bound with residue $n$ can
be decreased by any fixed positive infinitesimal and remains positive with
the same residue. There is no least upper bound.
\end{proof}

\begin{example}[A bounded self-adjoint operator without a field-valued norm]\label{ihs:hh:ex:noleast}
Let $Te_n=(1-1/(n+1))e_n$ on $\ell^2(\N_{\ge1})$. Its ordinary norm is one
but is not attained. Indeed, for every nonzero $u$,
\[
 \norm{u}_H^2-\norm{Tu}_H^2
 =\sum_{n\ge1}\left(1-\left(1-\frac1{n+1}\right)^2\right)|u_n|^2>0,
\]
where this is an ordinary convergent real sum. The diagonal values approach
one, proving the supremum is one. Hence its Hahn extension has many
field-valued bounds, including $1-s$ and $1-2s$, but no least one.
\end{example}

\subsection{The two mechanisms are independent}\label{ihs:hh:sec:independent}
The operator $D^2+s^2I$ in \cref{ihs:hh:thm:coercive} has an attained least
upper norm bound $1+s^2$ but is not invertible. The operator in
\cref{ihs:hh:ex:noleast} fails norm attainment for a different reason: every
individual Hahn vector has a strictly submaximal \emph{ordinary leading
ratio}, so infinitesimally smaller bounds cannot be distinguished by any one
vector. Neither phenomenon should be used as an explanation for the other.

\section{Why finite invertibility is misleading}\label{ihs:sec:finitesections}

Let $D_N=\diag(1,1/2,\ldots,1/N)$ and $C_N=D_N^2+s^2I$. Every $C_N$ is
invertible over $\K$. For the truncated right-hand side
$f_N=(n^{-2})_{1\le n\le N}$,
\begin{equation}\label{ihs:hh:eq:finitesolution}
 (C_N^{-1}f_N)_n=\frac{1}{1+n^2s^2}=\sum_{k\ge0}(-1)^k n^{2k}s^{2k}.
\end{equation}
The constant coefficient of this finite solution is the $N$-vector of ones,
whose ordinary Hilbert norm is $\sqrt N$. No single ordinary $\ell^2$ vector
can serve as the stable constant coefficient as $N\to\infty$. This directly
identifies the obstruction missed by every finite matrix determinant.

For $A_N=D_N-isI$, all diagonal resolvents likewise exist, with
\[
 (n^{-1}-is)^{-1}=\sum_{k\ge0}i^k n^{k+1}s^k.
\]
The leading inverse coefficient has ordinary norm $N$. Infinite extension
would require the unbounded ordinary operator $\diag(n)$ as a bounded
coefficient, contradicting \cref{ihs:hh:thm:adjoint}. More strongly,
particular constant data fail the vector-space condition of
\cref{ihs:hh:thm:smooth}; thus the problem is actual nonsurjectivity, not
just exclusion from a preferred inverse algebra.

\begin{remark}[The same matrix in the other part]\label{ihs:hh:rem:othersection}
The operator $\diag(n)$ excluded here is admitted in Part~II,
where the coefficient algebra is $\RCF_I(\C)$ rather than $\B$, and where
the finite-section question has a positive answer of a different kind:
\cref{ihs:rf:thm:finite} gives an exact finite spatial certificate for each
individual coefficient. The two statements are about different algebras and
do not conflict; see \cref{ihs:sec:accum}.
\end{remark}

\section{Part I: interpretation and remaining boundaries}\label{ihs:hh:sec:boundaries}

\subsection{What has been settled in this construction}\label{ihs:hh:sec:settled}
The questions raised by this particular scalar-extension model now have
sharp answers. Adjointability is not an opaque class condition: it is
exactly global Hahn support with bounded ordinary coefficients. For constant
normal operators, infinitesimal resolvent behavior is completely classified
by the ordinary spectral accumulation set. Infinitesimal perturbations of
injective leading operators cannot remove any ordinary range defect. The
leading ordinary coefficient gives a direct and robust witness, even where
every finite section is invertible.

These statements distinguish spectral membership from eigenvalue existence,
compactness in the original Hilbert topology from compactness in the
valuation topology, upper norm bounds from a least such bound, and metric
completeness from orthogonal completeness. None of those pairs can be
silently identified after extending the scalars.

\subsection{What would require a different theorem}\label{ihs:hh:sec:different}
The formula \eqref{ihs:hh:eq:exactspectrum} concerns \emph{constant normal}
operators. General Hahn operators with multiple coefficient scales can have
more complicated spectra; \cref{ihs:hh:ex:backshift} already shows why the
point-spectrum rigidity fails without normality. The perturbation theorem
handles injective leading coefficients and proves a large class of
obstructions, but is not a classification of all spectra of all Hahn
operator series.

Two later parts of this report supply such different theorems for two
classes, and leave the general classification open.
\begin{itemize}
\item \emph{Constant nonnormal operators.} Part~III computes the
 algebra-relative spectrum $\sigma^{\mathrm{adj}}_\Gamma(T)$ of every
 constant bounded $T$, with the Drazin spectrum in place of the accumulation
 set (\cref{ihs:dz:thm:halo}). The Volterra operator has no ordinary
 accumulation point and still acquires the whole monad
 (\cref{ihs:dz:sec:volterra}), so the accumulation-only formula is false
 without normality. The bijectivity spectrum $\sigma^{\mathrm{alg}}_\Gamma$
 of nonnormal constant operators is determined there only in examples.
\item \emph{A nonconstant class.} Part~IV treats $T+E$ with $T$ injective
 compact self-adjoint and $E$ a self-adjoint Hahn series of positive order
 that need not commute with $T$, over a divisible $\Gamma$: both spectra
 agree and equal the monad $\mK$ together with finite Hermitian clusters
 (\cref{ihs:fr:thm:main}). The eigenvalues move, unlike those of a constant
 normal operator (\cref{ihs:hh:thm:point}), but only within their clusters.
 This local answer does \emph{not} assemble into a global spectral theorem
 with bounded coefficients: an explicit trace-class example has no
 Hahn-unitary diagonalizer in $\Ahh$ and no coefficientwise eigenvalue
 trace sum or determinant product
 (\cref{ihs:fr:thm:nounitary,ihs:fr:thm:nolidskii,ihs:fr:thm:noproduct}).
\end{itemize}
Spectra of general Hahn operator series with several coefficient scales,
including nonnormal nonconstant ones, remain unclassified.

One could change the vector construction to allow coefficientwise objects
outside $H$, choose a different sequence-space completion, or impose a
stronger orthogonal-completeness axiom. Such changes must be specified and
analyzed rather than treated as harmless notational variations. Enlarging
only the exponent group is not enough, by \cref{ihs:hh:cor:extension}.

\subsection{Why this matters beyond a counterexample}\label{ihs:hh:sec:beyond}
In finite dimension, replacing $\R/\C$ by a real-closed field and its
complexification preserves much of spectral algebra. In infinite dimension,
that replacement cannot by itself preserve the analytic role of Hilbert
space. An infinitesimal imaginary displacement may fail to produce a
resolvent, even when it gives a strict lower norm bound. An infinitesimal
positive regularizer may leave a continuum of unsolved linear equations.
These are precise obstructions to using the model as an automatic
replacement for complex Hilbert space, not a rejection of surcomplex
analysis in every possible formulation.

The constructive side is equally useful. The automatic adjoint theorem gives
a well-defined algebra for exact symbolic manipulation. The inverse formulas
identify every ordinary normal resolvent that actually exists. The
inverse-smooth criterion identifies exactly which constant data remain
solvable at a monomial infinitesimal shift. The defect splitting gives
explicit noncanonical projections and an exact algebraic quotient. Together
these form a workable theory with both positive structure and sharp limits.

\clearpage
\part*{Part II. Row- and column-finite Hahn matrices: exact synthesis and algebra-relative spectra}
\addcontentsline{toc}{part}{Part II. Row- and column-finite Hahn matrices}

\emph{Standing hypotheses for Part II.} $\Gamma$ is a set-sized ordered
abelian group, \emph{not assumed divisible}; $I$ is a nonempty index set.
``Spectrum'' in this part means $\sigma_{\Arf}$, failure of two-sided
invertibility inside the named unital $\K$-algebra
$\Arf=\RCF_I(\C)((t^\Gamma))$, acting on $\Hrf=\C^{(I)}((t^\Gamma))$. It is
not a Banach-algebra spectrum, and it is not the failure-of-bijectivity
spectrum of Part~I, which lives on a different vector space.

\section{The question, the answer, and its scope}\label{ihs:rf:sec:question}

Passing from finitely many to infinitely many coordinates changes the
problem: products can cease to be defined, formal eigenvectors can escape the
chosen vector space, and infinitely many small denominators need not define
a bounded operator.

We address the following precise question. Suppose an infinite diagonal
matrix has distinct ordinary real entries and is perturbed by genuinely
infinitesimal surcomplex coefficients. Is there a natural associative
operator algebra, together with a vector space carrying a positive
surreal-valued inner product, in which \emph{all} eigenvectors can be
constructed coherently, summed against arbitrary vectors, and used to
recover exact spectral projections and resolvents?

The answer is affirmative in a coefficientwise locally finite Hahn category.
The category is part of the theorem, not an incidental convention. The
operators need not be bounded on an ordinary Hilbert space; the vector space
is not ordinary $\ell^2$; and Hahn summation is not shorthand for a sequence
of topological limits. Those distinctions allow accumulation of real energy
levels while retaining coefficientwise exactness.

The repository was checked for this part at commit \texttt{a3124af}; the
complete commit identifier appears in the bibliography \cite{RepoPinB} and
in \cref{ihs:app:auditB}. In particular, a still-unfinished Lean
implementation of a normal-form bridge is not being treated as an open
mathematical conjecture. The surcomplex interpretation in
\cref{ihs:sec:transport} uses the established normal-form theorem
\cite{Gonshor}, not an unproved declaration in a software project.

\section{The coefficient algebra and its Hahn vector space}\label{ihs:sec:category}

\subsection{Row- and column-finite coefficients}\label{ihs:rf:sec:rcf}
Let $I$ be a nonempty set. A matrix $M=(M_{ij})_{i,j\in I}$ is
\emph{row- and column-finite} if each of its rows and each of its columns
has finite support. The number of nonzero entries need not be uniformly
bounded as the row or column varies. Set
\[
 R_I=\RCF_I(\C),\qquad V_I=\C^{(I)}.
\]
Every diagonal matrix with entries in $\C$ belongs to $R_I$, even when its
entries are unbounded in the ordinary absolute value.

\begin{lemma}\label{ihs:rf:lem:rcf}
The usual finite row-by-column multiplication makes $R_I$ an associative
unital $\C$-algebra. Conjugate transpose is an involution, and $R_I$ acts
faithfully on $V_I$. With the ordinary finite-support inner product,
$\ip{Mx}{y}=\ip{x}{M^*y}$.
\end{lemma}
\begin{proof}
For a fixed row of $MN$, only the finitely many rows of $N$ selected by that
row of $M$ can contribute. Their supports have finite union. The analogous
column argument proves column-finiteness. Every entry in an iterated product
involves only finitely many intermediate indices, so associativity follows
from distributivity of finite sums. The identity matrix belongs to $R_I$,
and conjugate transpose interchanges the two finiteness conditions. A
finite-support vector is taken to a finite-support vector because the
relevant columns have finite union of supports. The adjoint identity is a
finite computation. Finally, an operator acting as zero on every standard
vector $e_i$ has every column zero.
\end{proof}

The requirement is deliberately imposed on \emph{coefficients}, rather than
on the complete matrix of Hahn entries. The eventual diagonalizer will
generally have infinitely many nonzero entries in a row or column, spread
over increasingly complicated exponents.

\subsection{The Hahn algebra and module}\label{ihs:rf:sec:hahnobjects}
With the conventions of \cref{ihs:sec:supportconv}, define
\begin{equation}\label{ihs:rf:eq:spaces}
 \K=\C((t^\Gamma)),\qquad \Fld=\R((t^\Gamma)),\qquad
 \Arf=R_I((t^\Gamma)),\qquad \Hrf=V_I((t^\Gamma)).
\end{equation}
The usual Hahn convolution makes $\K$ a field and $\Fld$ an ordered
subfield; a nonzero real series is positive exactly when its leading
coefficient is positive. The field property does not require $\Gamma$ to be
divisible: a nonzero series factors as a nonzero constant times a monomial
times a unit of the form $1+{}$positive support, whose inverse is given by
\cref{ihs:lem:neumann}. Write
$\mF=\{f\in\Fld:\supp f\subset\Gamma_{>0}\}$; it includes zero and is not
the cone of positive-valued elements.

\begin{proposition}\label{ihs:rf:prop:algebra}
The space $\Arf$ is an associative unital $\K$-algebra with involution
\[
 \left(\sum M_\gamma t^\gamma\right)^*=\sum M_\gamma^*t^\gamma.
\]
It acts faithfully on the $\K$-vector space $\Hrf$ by convolution. For
products in this algebra and its module,
\[
 v(XY)\ge v(X)+v(Y),\qquad v(Xx)\ge v(X)+v(x).
\]
The order function on $\Arf$ need not be a field valuation: if $I$ contains
two distinct indices, $\Arf$ has zero divisors. For a singleton $I$,
$\Arf=\K$ is a field and its order function is the field valuation.
\end{proposition}
\begin{proof}
At a prescribed exponent, \cref{ihs:lem:support}(i) leaves finitely many
coefficient products; \cref{ihs:rf:lem:rcf} puts each in $R_I$ or $V_I$. The
same finite-factor support argument permits reassociation of every triple
product. Scalar Hahn series act through scalar matrices. The support of a
product is contained in the sum of the supports, proving the displayed
inequalities. Faithfulness follows by evaluating on each $e_i$ and then
comparing coefficients. The adjoint identity follows coefficientwise.
For distinct $i,j\in I$, the nonzero constant diagonal matrix units
$E_{ii}$ and $E_{jj}$ belong to $R_I$ and satisfy $E_{ii}E_{jj}=0$.
If $I$ is a singleton, identifying a matrix with its sole entry gives
$R_I=\C$ and hence $\Arf=\K$.
\end{proof}

\subsection{The inner product and its positivity}\label{ihs:rf:sec:inner}

\begin{proposition}\label{ihs:rf:prop:inner}
The formula
\begin{equation}\label{ihs:rf:eq:inner}
 \ip{x}{y}
 =\sum_{\alpha,\beta}\ip{x_\alpha}{y_\beta}_{\C^{(I)}}t^{\alpha+\beta}
 =\sH_{i\in I}\overline{x_i}y_i
\end{equation}
defines a conjugate-linear/linear $\K$-valued inner product on $\Hrf$.
For $x\ne0$,
\[
 \ip{x}{x}\in\Fld_{>0},\qquad v(\ip{x}{x})=2v(x).
\]
The adjoint relation $\ip{Xx}{y}=\ip{x}{X^*y}$ holds for $X\in\Arf$.
\end{proposition}
\begin{proof}
For each total exponent there are finitely many pairs $(\alpha,\beta)$, and
each pair involves finite-support vectors. Thus the family over $i$ in
\eqref{ihs:rf:eq:inner} satisfies both conditions for strong Hahn
summability. Sesquilinearity and the adjoint relation follow from finite
coefficient calculations. If $\gamma=v(x)$, the coefficient at $2\gamma$ in
$\ip{x}{x}$ is
\[
 \sum_i \abs{(x_\gamma)_i}^2>0.
\]
There can be no smaller exponent, and no other pair of support exponents
sums to $2\gamma$. This proves positivity and the valuation formula.
\end{proof}

\begin{remark}[No divisibility, and no norm]\label{ihs:rf:rem:nodiv}
\Cref{ihs:rf:prop:inner} is positive without assuming that every positive
element of $\Fld$ has a square root. Indeed, for arbitrary $\Gamma$ the
field $\Fld$ need not be real closed. The only roots used in the main
construction are roots of units with residue $1$, obtained by binomial Hahn
series through \cref{ihs:cor:binomial}. In particular no analogue of
$\norm{x}=\sqrt{\ip{x}{x}}$ is introduced here; the norm of
\cref{ihs:hh:prop:inner} is a Part I object, but its construction does
not require divisibility. The even-leading-exponent argument also supplies
a square root of the squared length in Part II; we simply do not use a norm
in this part's algebraic construction.
\end{remark}

\begin{lemma}[Strong sums are respected]\label{ihs:rf:lem:strongaction}
If $(x_a)$ is strongly summable in $\Hrf$ and $X\in\Arf$, then $(Xx_a)$ is
strongly summable and $X(\sH_a x_a)=\sH_a Xx_a$. The analogous statement
holds for multiplication by a fixed element in $\Arf$ and for taking the
inner product with a fixed vector.
\end{lemma}
\begin{proof}
The union of resulting supports lies in the sum of two well-ordered sets. At
any exponent there are finitely many pairs of coefficient exponents, and for
each such pair only finitely many family members can contribute. The claimed
identities are therefore identities between finite sums at each coefficient.
\end{proof}

\section{Exact infinite diagonalization}\label{ihs:sec:diagonal}

\subsection{The divided commutator}\label{ihs:rf:sec:divided}
Let $(d_i)_{i\in I}$ be pairwise distinct elements of $\C$, and put
$D=\diag(d_i)$. Define two coefficientwise linear maps on $R_I$ and hence on
$\Arf$:
\begin{equation}\label{ihs:rf:eq:maps}
 \pi(M)=\diag(M_{ii}),\qquad
 \mathcal R_D(M)_{ij}=
 \begin{cases}
 M_{ij}/(d_i-d_j),&i\ne j,\\
 0,&i=j.
 \end{cases}
\end{equation}
They do not enlarge Hahn support, and they preserve the row- and column-finiteness of
every coefficient. On off-diagonal matrices, $\mathcal R_D$ is the inverse
of $X\mapsto[D,X]$.

No uniform bound for $\mathcal R_D$ in a real operator norm is claimed.
Every individual denominator is a nonzero ordinary complex number and
therefore has Hahn valuation zero. Arbitrarily small ordinary gaps do not
change supports. This is exactly the point at which the pairwise-distinct
hypothesis does its work, and \cref{ihs:rf:rem:monoid} records what happens
when it is absent.

\subsection{A similarity theorem without self-adjointness}\label{ihs:rf:sec:similarity}

\begin{theorem}[Coherent simple-residue diagonalization]\label{ihs:rf:thm:similarity}
Let $D=\diag(d_i)\in R_I$ have pairwise distinct complex entries, and let
$B\in\Arf$ have positive support. There exist unique
\[
 V=\Id+W\in\Id+(\Arf)_{>0},\qquad \pi(W)=0,
 \qquad \Lambda=D+\Delta,
\]
with $\Delta$ diagonal of positive support, such that
\begin{equation}\label{ihs:rf:eq:AV}
 (D+B)V=V\Lambda.
\end{equation}
The matrix $V$ is invertible in $\Arf$. If $S=\supp B$, then
\[
 \supp W,\ \supp\Delta\ \subseteq S^*\setminus\{0\}.
\]
The construction uses only finite coefficient operations, the nonzero
numbers $d_i-d_j$, and strong Hahn summation.
\end{theorem}
\begin{proof}
Introduce a formal bookkeeping variable $h$, distinct from the Hahn
monomials, and seek
\[
 V(h)=\sum_{n\ge0}h^nV_n,\qquad
 \Lambda(h)=D+\sum_{n\ge1}h^n\Delta_n,
 \qquad V_0=\Id,\quad \pi(V_n)=0\ (n\ge1).
\]
At degree $n\ge1$, the equation $(D+hB)V(h)=V(h)\Lambda(h)$ is
\[
 [D,V_n]+BV_{n-1}=\Delta_n+\sum_{k=1}^{n-1}V_{n-k}\Delta_k.
\]
Thus the recursion is forced:
\begin{align}
 T_n&=BV_{n-1}-\sum_{k=1}^{n-1}V_{n-k}\Delta_k,\label{ihs:rf:eq:Trec}\\
 \Delta_n&=\pi(T_n)=\pi(BV_{n-1}),\label{ihs:rf:eq:Drec}\\
 V_n&=-\mathcal R_D(T_n).\label{ihs:rf:eq:Vrec}
\end{align}
The second equality in \eqref{ihs:rf:eq:Drec} follows because multiplying an
off-diagonal $V_{n-k}$ by a diagonal matrix does not create diagonal
entries. These formulas define every coefficient and prove the formal
eigenidentity.

Inductively, $V_n$ and $\Delta_n$ are homogeneous of degree $n$ in $B$, with
support contained in $nS$. Each coefficient is in $R_I$, by
\cref{ihs:rf:lem:rcf} and the support preservation of
\eqref{ihs:rf:eq:maps}. \Cref{ihs:cor:evaluate} therefore licenses the
substitutions
\[
 V=\Id+\sH_{n\ge1}V_n,\qquad \Delta=\sH_{n\ge1}\Delta_n.
\]
It also licenses substitution in the products in the formal equation,
proving \eqref{ihs:rf:eq:AV}. The support claim follows at once. The inverse
of $V$ exists by \cref{ihs:cor:binomial}.

For uniqueness, suppose $(V',\Lambda')$ is another pair with the stated
normalizations. Then $T=V^{-1}V'$ belongs to $\Id+(\Arf)_{>0}$ and satisfies
$\Lambda T=T\Lambda'$. The diagonal equation is
$(\lambda_i-\lambda_i')T_{ii}=0$. Since $T_{ii}$ has residue $1$, it is
nonzero, so $\lambda_i=\lambda_i'$. For $i\ne j$, the element
$\lambda_i-\lambda_j$ has nonzero residue $d_i-d_j$, and hence $T_{ij}=0$.
Thus $T$ is diagonal. Finally $V_{ii}=V'_{ii}=1$ gives $T_{ii}=1$, proving
$T=\Id$ and the asserted uniqueness.
\end{proof}

The bookkeeping construction does not assert that $h=1$ is a permissible
substitution in arbitrary formal power series. It is permissible here
because the degree-$n$ coefficient is supported in $nS$, and
\cref{ihs:lem:support}(ii) supplies the missing summability proof.

\subsection{Positivity turns the similarity into a unitary}\label{ihs:rf:sec:unitary}

\begin{theorem}[Normalized infinite Hahn spectral theorem]\label{ihs:rf:thm:unitary}
Assume now that every $d_i$ is real, that the $d_i$ are pairwise distinct,
and that $B=B^*\in\Arf$ has positive support. With $A=D+B$, there is a
unique pair $(U,\Lambda)$ satisfying
\begin{equation}\label{ihs:rf:eq:unitarymain}
 AU=U\Lambda,\qquad U^*U=UU^*=\Id,\qquad U\in\Id+(\Arf)_{>0},
\end{equation}
where $\Lambda=\diag(\lambda_i)$, $\lambda_i\in\Fld$,
$\lambda_i-d_i\in\mF$, and every diagonal entry $U_{ii}$ is real and
strictly positive. Both $U-\Id$ and $\Lambda-D$ are supported in
$S^*\setminus\{0\}$, where $S=\supp B$.

Here $\mF$ includes zero and is not the cone of positive-valued elements.
In particular, the correction to an eigenvalue may have either sign.
\end{theorem}
\begin{proof}
Use \cref{ihs:rf:thm:similarity} and write $v_i=Ve_i$. Since
$Av_i=\lambda_i v_i$ and $A=A^*$,
\[
 \lambda_i\ip{v_i}{v_i}=\ip{v_i}{Av_i}=\overline{\ip{v_i}{Av_i}}.
\]
The norm square is positive and nonzero by \cref{ihs:rf:prop:inner};
therefore $\lambda_i=\overline{\lambda_i}$. For distinct $i,j$,
self-adjointness similarly gives
\[
 (\lambda_j-\lambda_i)\ip{v_i}{v_j}=0.
\]
The eigenvalue difference has nonzero real residue, so the columns are
orthogonal.

Consequently $N=V^*V$ is diagonal. Because $V=\Id+W$ and $W_{ii}=0$,
\begin{equation}\label{ihs:rf:eq:normalization}
 N_{ii}=1+\ip{We_i}{We_i},\qquad N=\Id+\pi(W^*W).
\end{equation}
The diagonal entries are positive with residue $1$. The binomial series
\[
 H=N^{-1/2}=\sH_{m\ge0}\binom{-1/2}{m}(N-\Id)^m
\]
exists by \cref{ihs:cor:binomial}, is real diagonal, and satisfies
$HNH=\Id$. Its diagonal entries have leading coefficient $1$, hence are
positive. Define $U=VH$. Since $H$ commutes with $\Lambda$, $AU=U\Lambda$
and $U^*U=HNH=\Id$. Both $V$ and $H$ are invertible in $\Arf$, so $U$ is
invertible and $U^*=U^{-1}$; this proves the other unitary identity as well.
Notice that in an infinite algebra a left inverse alone would not have
sufficed.

Because $V_{ii}=1$, one has $U_{ii}=H_{ii}>0$. The support claims follow
from \cref{ihs:lem:support}(ii) applied to the homogeneous expansions; no
exponent division is involved in the binomial normalization.

For uniqueness, let $(U',\Lambda')$ satisfy the same requirements. The
intertwining matrix $T=U^*U'$ has residue $\Id$. As in the uniqueness proof
of \cref{ihs:rf:thm:similarity}, first $\Lambda'=\Lambda$ and then $T$ is
diagonal. It is unitary. From $U'_{ii}=U_{ii}T_{ii}$ and positivity of both
diagonal entries, $T_{ii}$ is positive real. The identity
$\overline{T_{ii}}T_{ii}=1$ then forces $T_{ii}=1$.
\end{proof}

\begin{corollary}[No new scales]\label{ihs:rf:cor:scales}
Assume the hypotheses of \cref{ihs:rf:thm:unitary}: in particular that
$D=\diag(d_i)$ has \emph{pairwise distinct ordinary real} entries, so that
every divisor $d_i-d_j$ used in \eqref{ihs:rf:eq:maps} is a nonzero ordinary
complex constant of valuation zero, and that $B=B^*$ has globally
well-ordered positive support $S$. Then, under those hypotheses, all
eigenvalues $\lambda_i$ and all entries of the normalized diagonalizer $U$
lie in the original Hahn field $\K$, and their correction supports lie in
the additive monoid $S^*$ generated by $S$; consequently the result holds
over an arbitrary ordered exponent group, with no passage to its divisible
hull. If in addition the coefficients of $A$ are real Hahn, the
diagonalizer can be taken, and by uniqueness is, real Hahn orthogonal.
\end{corollary}
\begin{proof}
The support assertion was proved in
\cref{ihs:rf:thm:similarity,ihs:rf:thm:unitary}. When all coefficients are
real, the recursion and the binomial series use real operations throughout.
\end{proof}

\begin{remark}[The monoid statement is not true without the distinct-residue hypothesis]\label{ihs:rf:rem:monoid}
The support conclusion of \cref{ihs:rf:cor:scales} must not be lifted out of
its hypothesis. The collection's finite-dimensional spectral report contains
an example, its \texttt{spec:ex:monoid}, which shows that the corresponding
unconditional statement is false: with $\Gamma=\Z$ and
\[
 A=\begin{pmatrix}t^3&t^4\\t^4&0\end{pmatrix},
 \qquad
 \lambda_+=t^3+t^5-t^7+2t^9-\cdots,
\]
the exponent $5$ is not in the nonnegative monoid generated by the raw input
exponents $3$ and $4$, and that report concludes that a theorem asserting
that eigenvalue supports stay in the raw input monoid ``would therefore be
false'' \cite{RepoSpectral}.

There is no conflict. Writing that matrix as $D+B$ with $B$ of positive
support forces $D=\diag(0,0)$: both diagonal entries, $t^3$ and $0$, have
residue zero. The residues are therefore not pairwise distinct, the divided
commutator $\mathcal R_D$ of \eqref{ihs:rf:eq:maps} does not exist, and
\cref{ihs:rf:thm:similarity,ihs:rf:thm:unitary,ihs:rf:cor:scales} do not
apply. \Cref{ihs:rf:cor:scales} divides only by differences $d_i-d_j$ of
pairwise distinct ordinary complex constants, each of valuation zero; that
is precisely what keeps every correction exponent inside $S^*$. The report
cited above states the correct unconditional invariant for the general
finite case --- the input \emph{subgroup}, supplemented by the positive
monoids of the normalized perturbations at each node --- and that is a
different statement from the one proved here.
\end{remark}

\subsection{Initial coefficients and their meaning}\label{ihs:rf:sec:initial}

For a one-scale perturbation $B=t^\eta B_1+t^{2\eta}B_2+\cdots$, $\eta>0$,
the first terms are
\begin{align}
 \lambda_i={}&d_i+t^\eta(B_1)_{ii}
 +t^{2\eta}\left((B_2)_{ii}
 +\sum_{j\ne i}\frac{\abs{(B_1)_{ji}}^2}{d_i-d_j}\right)
 +\bigo(t^{3\eta}),\label{ihs:rf:eq:lambda2}\\
 U_{ji}={}&t^\eta\frac{(B_1)_{ji}}{d_i-d_j}
 +\bigo(t^{2\eta})\quad(j\ne i),\label{ihs:rf:eq:U1}\\
 U_{ii}={}&1-\frac{t^{2\eta}}2
 \sum_{j\ne i}\frac{\abs{(B_1)_{ji}}^2}{(d_i-d_j)^2}
 +\bigo(t^{3\eta}).\label{ihs:rf:eq:Udiag2}
\end{align}
Here the remainder notation is a valuation statement within the one-scale
series under discussion. Every displayed sum is finite for a fixed $i$.
Equations
\eqref{ihs:rf:eq:lambda2}--\eqref{ihs:rf:eq:Udiag2} follow directly from
\eqref{ihs:rf:eq:Trec}--\eqref{ihs:rf:eq:normalization}; they are the
familiar perturbation coefficients, not independent novelty claims.

What changes in the present theorem is that all orders and all coordinates
fit into one operator with globally controlled Hahn support. A uniform bound
for the real numbers in \eqref{ihs:rf:eq:U1} is unnecessary, but a finite
coefficient support in every row and column remains essential to the chosen
algebra.

\section{Spectral synthesis, commutants, and projections}\label{ihs:sec:synthesis}

Throughout this section the hypotheses of \cref{ihs:rf:thm:unitary} hold,
and $u_i=Ue_i$ denotes its normalized eigenvector.

\begin{theorem}[Exact Hahn spectral synthesis]\label{ihs:rf:thm:synthesis}
The family $(u_i)_{i\in I}$ is orthonormal. Every $x\in\Hrf$ has the unique
strongly summable expansion
\begin{equation}\label{ihs:rf:eq:synthesis}
 x=\sH_{i\in I}u_i\ip{u_i}{x}.
\end{equation}
More precisely, the coordinate map
$x\mapsto(\ip{u_i}{x})_{i\in I}$ identifies $\Hrf$ with itself when the
right side is interpreted as the Hahn vector $U^*x$. One has the identities
\begin{align}
 \ip{x}{x}&=\sH_{i\in I}\abs{\ip{u_i}{x}}^2,\label{ihs:rf:eq:parseval}\\
 Ax&=\sH_{i\in I}\lambda_i u_i\ip{u_i}{x}.\label{ihs:rf:eq:spectralaction}
\end{align}
The eigenvalues of $A$ on $\Hrf$ are exactly the $\lambda_i$, and each
corresponding eigenspace is $\K u_i$.
\end{theorem}
\begin{proof}
Orthonormality is the entrywise meaning of $U^*U=\Id$. Put $y=U^*x$. At each
exponent, $y$ has finite coordinate support. Therefore the family
$(e_i y_i)_{i\in I}$ is strongly summable in $\Hrf$ and sums to $y$.
\Cref{ihs:rf:lem:strongaction} gives
\[
 x=Uy=\sH_i Ue_i y_i=\sH_i u_i\ip{u_i}{x},
\]
proving existence and summability. Taking inner products with each $u_i$
proves uniqueness whenever such a strong expansion is given. The diagonal
action of $\Lambda$ takes $\Hrf$ to itself because $\Lambda\in\Arf$.
Applying $U$ to the coordinate expansion of $\Lambda y$ proves
\eqref{ihs:rf:eq:spectralaction}. The norm identity follows from
$\ip{x}{x}=\ip{y}{y}$ and \cref{ihs:rf:prop:inner}.

Finally, if $Ax=zx$ with $x\ne0$, then $(\lambda_i-z)y_i=0$ for all $i$. At
least one $y_i$ is nonzero, so $z=\lambda_i$ for that index. Distinct
residues make the $\lambda_i$ pairwise distinct; all other coordinates of
$y$ vanish. This proves the eigenspace statement.
\end{proof}

The expansion need not be finite over $\K$, so $(u_i)$ is generally not a
Hamel basis. Nor does \eqref{ihs:rf:eq:parseval} involve the ordinary
convergence of a series of nonnegative real numbers: its terms and sum lie
in $\Fld$, and the assertion is strong Hahn summability.

\subsection{The coherent diagonal algebra}\label{ihs:rf:sec:diagalg}

Let
\begin{equation}\label{ihs:rf:eq:diagalg}
 \Dg=(\C^I)((t^\Gamma))\subseteq\Arf
\end{equation}
be the algebra of diagonal Hahn operators. Thus a family $(a_i)_{i\in I}$ of
elements of $\K$ defines an element of $\Dg$ precisely when
$\bigcup_i\supp a_i$ is well ordered. There is no finite-incidence condition
on the indices $i$ here: an arbitrary diagonal matrix of complex numbers is
one legitimate coefficient in $R_I$.

For finite $I$, every family is coherent, so $\Dg=\K^I$; the same equality
holds for the trivial value group. If $I$ is infinite and $\Gamma\ne0$,
the inclusion is strict. Choose distinct indices $i_n$ and $\eta>0$, and
set $a_{i_n}=t^{-n\eta}$ for $n\ge1$, with all other entries zero. Their
support union contains the decreasing sequence $-\eta,-2\eta,\ldots$,
so this family does not belong to $\Dg$.

Coherence places no ordinary bound on a fixed diagonal coefficient:
the constant matrix with entries $n$ at $i_n$ and zero elsewhere belongs
to $\Dg$ when $I$ is infinite. These are separate conditions, and the
resolvent proof below uses precisely global support coherence.

\begin{theorem}[The complete commutant]\label{ihs:rf:thm:commutant}
Within the specified Hahn operator algebra,
\begin{equation}\label{ihs:rf:eq:commutant}
 \Cent_{\Arf}(A)=U\Dg U^*.
\end{equation}
In particular this commutant is commutative. Its idempotents are exactly
\begin{equation}\label{ihs:rf:eq:PJ}
 P_J=UE_JU^*,\qquad J\subseteq I,
\end{equation}
where $E_J=\diag(\mathbf 1_{i\in J})$. Every such idempotent is
self-adjoint. The correspondence $J\mapsto P_J$ is a complete Boolean
algebra isomorphism, with complement $\Id-P_J$ and product
$P_JP_L=P_{J\cap L}$.
\end{theorem}
\begin{proof}
If $X$ commutes with $A$, then $Y=U^*XU$ commutes with $\Lambda$. At entry
$(i,j)$ this says $(\lambda_i-\lambda_j)Y_{ij}=0$. For $i\ne j$ the first
factor is nonzero in the field $\K$, so $Y$ is diagonal. Its membership in
$\Arf$ is precisely the support coherence in \eqref{ihs:rf:eq:diagalg}. The
converse is immediate, proving \eqref{ihs:rf:eq:commutant}.

An idempotent of $\Dg$ has entries satisfying $a_i^2=a_i$ in a field, so
each is $0$ or $1$. Every choice is allowed, since its complete operator
support is contained in $\{0\}$. These are exactly the $E_J$. Conjugating
proves the projection statement and the finite Boolean identities. Arbitrary
unions and intersections of subsets give the joins and meets for the partial
order on idempotents defined by $P\le Q$ when $PQ=P$; no assertion about an
infinite operator-valued series is needed for completeness of this Boolean
algebra.
\end{proof}

\begin{corollary}[Projection sums are pointwise Hahn sums]\label{ihs:rf:cor:projsums}
For every subset $J\subseteq I$ and every $x\in\Hrf$,
\[
 P_Jx=\sH_{i\in J}u_i\ip{u_i}{x}.
\]
If $(J_a)_{a\in L}$ are pairwise disjoint, then $(P_{J_a}x)_{a\in L}$ is
strongly summable in $\Hrf$ and its sum is $P_{\bigcup_a J_a}x$.
\end{corollary}
\begin{proof}
In the coordinates $y=U^*x$, these are the sums of disjoint coordinate
pieces of a Hahn vector. At any exponent, only finitely many of the pieces
can be nonzero. Apply \cref{ihs:rf:lem:strongaction} to transport the sums
by $U$.
\end{proof}

\begin{remark}[Operator sums and vectorwise additivity]\label{ihs:rf:rem:notopsum}
For infinite $I$, the operator family $(P_{\{i\}})_{i\in I}$ is \emph{not}
strongly Hahn summable in $\Arf$: its coefficient at exponent zero is
$E_{\{i\}}\ne0$ for every $i$. Thus $\sH_iP_{\{i\}}$ is undefined in that
algebra, whereas $\sH_iP_{\{i\}}x=x$ holds for each $x\in\Hrf$.
This distinction already occurs for $B=0$.

It does not by itself exclude a spectral-measure interpretation.
Classical projection-valued measures require countable additivity in
the strong operator topology, which tests convergence on each Hilbert
vector, rather than an operator-norm sum; see
\cite[Definition~5.1 and footnote~18]{Williams}.
Here $J\mapsto P_J$ on the power set of $I$ is a projection-valued set
function with vectorwise \emph{Hahn} additivity by
\cref{ihs:rf:cor:projsums}. Such additivity does not imply convergence
in the intrinsic valuation topology.

For an explicit failure, take $\Gamma=\Q$, $I=\N_{>0}$, $D=\diag(n)$,
and $B=0$, so $U=\Id$ and $P_{\{n\}}=E_{\{n\}}$. Put
\[
 \gamma_n=1-\frac1{n+1},\qquad
 x=\sH_{n\ge1}t^{\gamma_n}e_n.
\]
The exponents increase strictly, hence form a well-ordered set, and
each coefficient has one coordinate. Thus $x\in\Hrf$, but
\[
 v\!\left(x-\sum_{n=1}^N P_{\{n\}}x\right)
 =\gamma_{N+1}<1
\]
for every $N$. The partial sums never enter the neighborhood of $x$
with error valuation greater than $1$. Consequently even this rank-one
example lacks countable additivity in the topology of pointwise
valuation convergence. The established object is the Hahn-additive
projection calculus; no general topologically countably additive
spectral-measure or Borel integration theorem is asserted.
\end{remark}

\subsection{An unrestricted ordinary label calculus}\label{ihs:rf:sec:labels}
Every function $f:I\to\C$ defines the operator
\[
 T_f=U\diag(f(i))U^*\in\Cent_{\Arf}(A).
\]
This is a unital involution-preserving homomorphism from $\C^I$ into
$\Arf$. No continuity or ordinary boundedness of $f$ is required. It is a
\emph{label multiplier calculus}: it does not assign values to an ordinary
function at arbitrary surreal arguments. More generally, a $\K$-valued label
family is allowed exactly when it belongs to the coherent diagonal algebra
$\Dg$. Polynomial expressions $p(A)$ have labels $p(\lambda_i)$ because
finite polynomial operations preserve coherence.

\section{The exact algebra-relative spectrum and all resolvents}\label{ihs:sec:resolvent}

\begin{definition}\label{ihs:rf:def:spectrum}
For $X\in\Arf$, define
\[
 \sigma_{\Arf}(X)=\{z\in\K:X-z\Id\text{ has no two-sided inverse in }\Arf\}.
\]
This is a spectrum relative to a unital algebra, with no norm or compactness
assumption. It is the second of the two senses recorded in
\cref{ihs:tab:senses}.
\end{definition}

Diagonalization alone does not automatically identify this spectrum.
Entrywise invertibility of infinitely many diagonal entries does not, in
general, guarantee that their inverses have globally well-ordered support.
The simple-residue hypothesis gives the additional coherence. For example,
on $I=\N_0$ and with $\eta>0$, the diagonal
$Q=\diag(t^{n\eta})$ belongs to $\Dg$, but its entrywise inverse has
support $\{-n\eta:n\ge0\}$ and is not coherent. Any matrix inverse would
have to have exactly those diagonal entries and zero off-diagonal entries,
by multiplying entrywise against $Q$. Thus $Q$ is not invertible in $\Arf$.
It is not onto $\Hrf$ either: the Hahn vector
$\sum_{n\ge0}t^{n\eta}e_n$ would require the preimage $(1,1,\ldots)$,
which has an inadmissible infinite-support coefficient at zero.

\begin{theorem}[Coherent resolvents at every noneigenvalue]\label{ihs:rf:thm:resolvent}
Under the hypotheses of \cref{ihs:rf:thm:unitary},
\begin{equation}\label{ihs:rf:eq:spectrum}
 \sigma_{\Arf}(A)=\{\lambda_i:i\in I\}.
\end{equation}
For $z\in\K$ outside that set,
\begin{equation}\label{ihs:rf:eq:resolvent}
 (A-z\Id)^{-1}=U\diag\bigl((\lambda_i-z)^{-1}\bigr)U^*\in\Arf.
\end{equation}
The conclusion and the proof also hold for the non-Hermitian
simple-complex-residue setting of \cref{ihs:rf:thm:similarity}, with
$U,U^*$ replaced by $V,V^{-1}$.
\end{theorem}
\begin{proof}
The eigenvector $u_i$ prevents $A-\lambda_i\Id$ from being invertible. It
remains to prove coherence of the diagonal inverse when every
$\lambda_i-z$ is nonzero. Write $\Lambda=D+\Delta$ with
$\supp\Delta\subset\Gamma_{>0}$.

First suppose $v(z)<0$. Then $\Lambda-z\Id=-z(\Id-z^{-1}\Lambda)$. Since
$v(\Lambda)\ge0$, the operator $z^{-1}\Lambda$ has positive support. The
inverse therefore belongs to $\Dg$ by \cref{ihs:lem:neumann}.

Now suppose $v(z)\ge0$, also allowing $z=0$. Write $z=c+\zeta$, where
$c\in\C$ is the coefficient at exponent zero and $\zeta$ has positive
support. If $c\ne d_i$ for every $i$, the ordinary diagonal matrix
$C=\diag((d_i-c)^{-1})$ belongs to $R_I$. It need not have bounded complex
entries, and none are required. The factorization
\[
 \Lambda-z\Id=(D-c\Id)\bigl(\Id+C(\Delta-\zeta\Id)\bigr)
\]
has a positive-support second factor and an invertible constant first
factor. Its inverse is again in $\Dg$.

Finally, if $c=d_j$, then there is exactly one such index $j$. Split the
diagonal algebra into its $E_{\{j\}}$ and $\Id-E_{\{j\}}$ parts. On the
second part the preceding argument applies, using arbitrary harmless values
at the omitted index to write the factorization. On the first part, the
inverse is the single scalar Hahn series $(\lambda_j-z)^{-1}$ placed at
index $j$. It exists by hypothesis. The union of the supports of these two
pieces is a finite union of well-ordered sets. Adding them gives the desired
diagonal inverse. Conjugation by $U$ proves \eqref{ihs:rf:eq:resolvent}.

Nothing in the coherence argument used reality of the distinct constants
$d_i$; the last assertion follows in the same way.
\end{proof}

\begin{example}[An accumulation point need not be in this spectrum]\label{ihs:rf:ex:accum}
Take $I=\N_{>0}$, $D=\diag(1/n)$, and $B=0$. Then $0$ is an ordinary
accumulation point of the labels, but
\[
 D^{-1}=\diag(n)\in R_I\subset\Arf.
\]
Consequently $0\notin\sigma_{\Arf}(D)$. This does not contradict the
ordinary bounded-operator spectrum on $\ell^2$, where the displayed inverse
is unbounded. It records the different algebra in which invertibility is
being tested. It also does not contradict \cref{ihs:hh:cor:compact}, which
puts the whole infinitesimal monad $\mK$ into the failure-of-bijectivity
spectrum of the same diagonal data \emph{as a compact operator on
$\ell^2((t^\Gamma))$}; \cref{ihs:sec:accum} sets the two side by side, and
\cref{ihs:sec:unify} derives both from Part~III's single formula.
\end{example}

\begin{corollary}[The inhomogeneous equation]\label{ihs:rf:cor:forcing}
If $z\notin\{\lambda_i:i\in I\}$, every equation $(A-z\Id)x=f$ with
$f\in\Hrf$ has the unique solution
\[
 x=\sH_i u_i\frac{\ip{u_i}{f}}{\lambda_i-z}.
\]
The series is strongly summable in $\Hrf$.
\end{corollary}
\begin{proof}
The coherent inverse diagonal operator takes $U^*f$ to a Hahn vector. Apply
$U$ and \cref{ihs:rf:thm:synthesis}. Invertibility gives uniqueness.
\end{proof}

\section{Weighted walks control every coefficient}\label{ihs:sec:locality}

\subsection{The graph attached to a Hahn perturbation}\label{ihs:rf:sec:graph}
Let $S=\supp B\subset\Gamma_{>0}$ and write $B=\sum_{s\in S}B_s t^s$.
Construct an edge-labeled directed graph on $I$ by putting an edge
\[
 j\xrightarrow{\ s\ }i \quad\text{when}\quad (B_s)_{ij}\ne0.
\]
Loops are allowed. Since $B_s=B_s^*$, every nonloop edge has a reverse edge
with the same label. For a fixed $s$, every vertex has finitely many
incoming and outgoing edges. The graph obtained by forgetting all labels
need not be locally finite when $S$ is infinite.

A walk has a finite sequence of edges; its weight is the sum of the positive
labels. The empty walk has weight zero. Although an edge label can be
infinitesimal relative to another in $\Gamma$, no edge has zero or negative
weight.

\begin{lemma}[Finite walks at a prescribed weight]\label{ihs:rf:lem:finitewalks}
For fixed starting vertex $i$ and fixed $\gamma\in\Gamma$, there are only
finitely many walks beginning at $i$ and having total weight $\gamma$. In
particular their visited vertices form a finite set.
\end{lemma}
\begin{proof}
\Cref{ihs:lem:support}(ii) gives finitely many ordered label words with
total weight $\gamma$. For any one such word, there are finitely many
choices of the first edge and, recursively, finitely many choices at each
subsequent step. The word length is finite, so the complete set of walks is
finite. Taking a finite union over words proves the assertion.
\end{proof}

\begin{theorem}[Coefficientwise walk dependence]\label{ihs:rf:thm:walks}
For the intermediate-normalized $V$ and the normalized unitary $U$, the
coefficient at a positive exponent $\gamma$ in $V_{ji}$ or $U_{ji}$ is a
finite sum of expressions supported on walks from $i$ to $j$ of weight
$\gamma$. The coefficient of $\lambda_i-d_i$ is a finite sum of expressions
supported on closed walks at $i$ of weight $\gamma$.

Each expression is a product of edge coefficients and their conjugates,
multiplied by a rational number and by inverse powers of nonzero real label
differences at visited vertices. Thus those coefficients are completely
determined by the finitely many relevant walks and the ordinary labels at
their vertices.

This is a dependence statement, not a claim that every permitted walk
contributes a nonzero term, or that its contribution is just its raw edge
product.
\end{theorem}
\begin{proof}
Work first with the homogeneous recursion
\eqref{ihs:rf:eq:Trec}--\eqref{ihs:rf:eq:Vrec}. A coefficient of $V_0=\Id$
is represented by an empty walk. Multiplication by $B$ attaches an edge to
the end of a walk, with matrix multiplication summing over the intermediate
vertex. The projection $\pi$ retains closed walks. In the term
$(V_{n-k}\Delta_k)_{ji}$, the scalar $(\Delta_k)_{ii}$ is represented by a
closed walk at $i$; traverse that closed walk and then a walk from $i$ to
$j$ representing $(V_{n-k})_{ji}$. Division by $d_j-d_i$ changes no walk or
exponent. This proves the assertion for $V$ and $\Lambda$ by induction on
the homogeneous degree.

For normalization, \eqref{ihs:rf:eq:normalization} expresses $N_{ii}-1$
through products $\overline{W_{ki}}W_{ki}$. More generally, at a fixed
coefficient one factor can come from one walk from $i$ to $k$ and the other
from another. Reverse the conjugated walk, using Hermitian symmetry of the
labeled edges. Their concatenation is a closed walk at $i$. Powers of
$N-\Id$ concatenate such closed walks; the binomial coefficients contribute
only rational scalars. Finally $U_{ji}=V_{ji}H_{ii}$ concatenates those
closed walks with a walk from $i$ to $j$.

\Cref{ihs:lem:support}(ii) bounds the homogeneous degrees at any prescribed
exponent, and \cref{ihs:rf:lem:finitewalks} bounds every index sum within
them. All coefficient formulas are therefore genuinely finite.
\end{proof}

\subsection{Marked-edge stability}\label{ihs:rf:sec:marked}

Suppose $B$ and $\widetilde B$ satisfy the same hypotheses with the same
$D$. Their support union $S$ is well ordered and positive. Use the union of
the two edge graphs, and mark an edge of label $s$ whenever
$(\widetilde B_s-B_s)_{ji}\ne0$. For a vertex $i$, let $\kappa_i$ be the
minimum weight of a closed walk at $i$ containing a marked edge, or
$+\infty$ if no such walk exists. The minimum exists whenever the set is
nonempty, since its weights form a subset of the well-ordered monoid $S^*$.

\begin{theorem}[Exact marked-walk stability]\label{ihs:rf:thm:marked}
The eigenvalues labeled by $i$ satisfy
\begin{equation}\label{ihs:rf:eq:marked}
 v(\widetilde\lambda_i-\lambda_i)\ge\kappa_i.
\end{equation}
More generally, the coefficient of $\widetilde U_{ji}-U_{ji}$ at exponent
$\gamma$ vanishes unless there is a walk from $i$ to $j$ of weight $\gamma$
containing a marked edge.

In particular, set $E=\widetilde B-B$. Then
\[
 v(\widetilde U-U)\ge v(E),\qquad v(\widetilde\Lambda-\Lambda)\ge v(E).
\]
If $E$ has zero diagonal and $S\ne\varnothing$, the second estimate sharpens
to
\begin{equation}\label{ihs:rf:eq:offdiagprecision}
 v(\widetilde\Lambda-\Lambda)\ge v(E)+\min S.
\end{equation}
Zero differences are interpreted as having valuation $+\infty$.
\end{theorem}
\begin{proof}
The recursions and the normalization use the same rational expressions in
the edge coefficients for both inputs, with identical denominators because
$D$ is fixed. Subtract a product of coefficients from its counterpart by
successively replacing its factors. Each resulting term contains a factor
from $E$ or its conjugate. It therefore has a marked edge in its associated
walk. This argument is finite at every exponent by
\cref{ihs:rf:thm:walks}, so it proves both dependence assertions and
\eqref{ihs:rf:eq:marked}.

A term containing a factor from $E$ has weight at least $v(E)$, since all
its other edges have positive weights. If $E$ has zero diagonal, a closed
walk containing a marked edge cannot have length one. It has at least one
further edge with weight at least $\min S$. This proves
\eqref{ihs:rf:eq:offdiagprecision}. The same bounds hold uniformly for every
matrix coefficient and hence for the global operator order function.
\end{proof}

These statements are valuation precision bounds, not analytic continuity in
a real operator norm. They are also stronger than merely saying that finite
jets depend on finite jets: they identify \emph{which spatial changes} can
first become visible at a given Hahn scale.

\subsection{Finite sections with exact coefficient certificates}\label{ihs:rf:sec:finitesec}

Let $F_0\subset I$ be finite and contain $i$, and let $A_{F_0}$ denote the
principal compression of $A$ to $F_0$. Its diagonalizer and eigenvalues are
normalized by the same rules. Write $\lambda_i^{F_0}$ for its eigenvalue
with residue $d_i$. Define the escape weight
\[
 \kappa_i(F_0)=\min\{\mathrm{wt}(w):w\text{ is a closed walk at }i
 \text{ and visits }I\setminus F_0\},
\]
again using $+\infty$ for an empty set.

\begin{theorem}[Finite-section exactness]\label{ihs:rf:thm:finite}
For every finite $F_0\ni i$,
\begin{equation}\label{ihs:rf:eq:finitebound}
 v(\lambda_i-\lambda_i^{F_0})\ge\kappa_i(F_0).
\end{equation}
For a fixed exponent $\gamma$, the coefficient of $\lambda_i$ is exactly
that of a finite principal section containing all vertices of closed walks
at $i$ of weight $\gamma$. The analogous statement for $U_{ji}$ holds with
walks from $i$ to $j$ and with both endpoints included. Consequently every
single requested coefficient has a finite spatial certificate, even when $I$
and $\supp B$ are infinite.
\end{theorem}
\begin{proof}
For the coefficient assertion, all relevant walks are finite in number by
\cref{ihs:rf:lem:finitewalks}. If their vertices lie in $F_0$, every finite
expression in \cref{ihs:rf:thm:walks} is unchanged by compression. In
particular, every coefficient strictly below $\kappa_i(F_0)$ agrees, which
proves \eqref{ihs:rf:eq:finitebound}.

Alternatively, remove all edges crossing between $F_0$ and its complement
while keeping the diagonal and all internal edges. The resulting infinite
operator is a direct sum on those two index sets. Its recursion restricted
to $F_0$ is exactly the recursion for $A_{F_0}$. A closed walk at $i$
contains a removed crossing edge exactly when it leaves $F_0$. Thus
\cref{ihs:rf:thm:marked} gives the same inequality.
\end{proof}

The finite section for one coefficient need not work for \emph{all}
coefficients below a prescribed arbitrary-rank exponent. Such an initial
segment can be infinite and can involve infinitely many spatial vertices.
The theorem guarantees finiteness at each exact weight; it does not replace
an infinite Hahn support by a finite truncation without additional data.

\begin{corollary}[One-scale radius bound]\label{ihs:rf:cor:radius}
Let $A=D+t^\eta C$, where $\eta>0$ and $C=C^*\in R_I$, and let $F_0$ contain
the graph ball of radius $r$ about $i$ in the undirected support graph of
$C$. Then
\begin{equation}\label{ihs:rf:eq:radiusbound}
 v(\lambda_i-\lambda_i^{F_0})\ge(2r+2)\eta.
\end{equation}
The coefficient of $t^{n\eta}$ in $\lambda_i$ depends only on the ball of
radius $\lfloor n/2\rfloor$ about $i$. The coefficient of $t^{n\eta}$ in
$U_{ji}$ vanishes when $\dist(i,j)>n$.
\end{corollary}
\begin{proof}
A walk that leaves the radius-$r$ ball requires at least $r+1$ edges to
reach the outside and at least $r+1$ more to return. Every edge has weight
$\eta$. A closed walk of length $n$ cannot visit a vertex farther than
$\lfloor n/2\rfloor$ from its starting point, while a walk from $i$ to $j$
has length at least their graph distance. Apply
\cref{ihs:rf:thm:walks,ihs:rf:thm:finite}.
\end{proof}

\subsection{The threshold is sharp}\label{ihs:rf:sec:sharp}

\begin{theorem}[A nonzero first boundary coefficient]\label{ihs:rf:thm:sharp}
Let $L\ge1$ and let $A_L(t)$ be the tridiagonal matrix on vertices
$0,\ldots,L$ with
pairwise distinct real diagonal entries $d_0,\ldots,d_L$ and off-diagonal
entries $t b_k$ and $t\overline{b_k}$ between $k$ and $k+1$, where
$b_k\in\C\setminus\{0\}$. Let $\lambda_0^{[L]}$ be its branch with residue
$d_0$, and compare it with the corresponding branch for the section
$0,\ldots,L-1$. Then
\begin{equation}\label{ihs:rf:eq:sharp}
 \lambda_0^{[L]}-\lambda_0^{[L-1]}
 =\frac{\prod_{k=0}^{L-1}\abs{b_k}^2}
 {(d_0-d_L)\prod_{k=1}^{L-1}(d_0-d_k)^2}\,t^{2L}
 +\bigo(t^{2L+2}).
\end{equation}
In particular the leading coefficient is nonzero, and the exponent in
\eqref{ihs:rf:eq:radiusbound} cannot be increased in general.
\end{theorem}
\begin{proof}
Use $q=t^2$ and put $c_k=\abs{b_k}^2$. Schur elimination on the tail gives
the scalar equation
\begin{equation}\label{ihs:rf:eq:continued}
 x-d_0=q c_0 H_1^{[L]}(x,q),\qquad
 H_k^{[L]}(x,q)=\frac{1}{x-d_k-q c_k H_{k+1}^{[L]}(x,q)},
\end{equation}
with $H_L^{[L]}=(x-d_L)^{-1}$. For $x$ near $d_0$, all constant terms of the
denominators are nonzero, so these are formal power series in $q$ with
coefficients rational in $x$, regular at $d_0$.

For $L\ge2$, subtract the full and truncated tail fractions. At the last
common vertex one obtains
\[
 H_{L-1}^{[L]}-H_{L-1}^{[L-1]}
 =\frac{q c_{L-1}}{(x-d_{L-1})^2(x-d_L)}+\bigo(q^2).
\]
At each preceding vertex, the identity
\[
 \frac1{a-qcH}-\frac1{a-qc\widetilde H}
 =\frac{qc(H-\widetilde H)}{(a-qcH)(a-qc\widetilde H)}
\]
adds one power of $q$, one factor $c$, and a squared leading denominator.
Thus
\[
 q c_0\bigl(H_1^{[L]}-H_1^{[L-1]}\bigr)
 =\frac{\prod_{k=0}^{L-1}c_k}
 {(x-d_L)\prod_{k=1}^{L-1}(x-d_k)^2}\,q^L+\bigo(q^{L+1}).
\]
To justify passage from the tail difference to the root difference, let
\[
 G_L(x,q)=x-d_0-qc_0H_1^{[L]}(x,q),\qquad
 x_L=\lambda_0^{[L]},\quad x_{L-1}=\lambda_0^{[L-1]}.
\]
Both roots are $d_0+\bigo(q)$. The formal divided difference of $G_L$
between these two roots is a unit $1+\bigo(q)$, since
$G_L(x,0)=x-d_0$. If $C(x)$ denotes the displayed coefficient of $q^L$,
then
\[
 G_L(x_{L-1},q)=-C(d_0)q^L+\bigo(q^{L+1}),\qquad
 G_L(x_L,q)=0.
\]
Subtracting and dividing by that unit gives
$x_L-x_{L-1}=C(d_0)q^L+\bigo(q^{L+1})$, as claimed.
For $L=1$ the truncated branch is exactly $d_0$, and the original Schur
equation gives the same formula directly. This proves \eqref{ihs:rf:eq:sharp}. The nonzero
denominator and the nonzero positive numerator prove sharpness.
\end{proof}

\section{An accumulating-level chain with an exact eigenpair}\label{ihs:sec:chain}

Consider $I=\N_0$, $\Gamma=\Z$, and the real symmetric operator
\begin{equation}\label{ihs:rf:eq:chain}
 (Ax)_n=\frac{x_n}{n+1}+t(x_{n-1}+x_{n+1}),\qquad x_{-1}=0.
\end{equation}
Its diagonal labels $d_n=1/(n+1)$ are pairwise distinct but accumulate at
zero, and the neighboring gaps
\[
 d_n-d_{n+1}=\frac1{(n+1)(n+2)}
\]
have no uniform positive real lower bound. \Cref{ihs:rf:thm:unitary}
nevertheless constructs all its normalized eigenvectors simultaneously in
$\C^{(\N_0)}((t))$.

\subsection{The first eigenvalue can be written in closed form}\label{ihs:rf:sec:closed}

\begin{proposition}[An exact branch and its eigenvector]\label{ihs:rf:prop:closed}
For the operator \eqref{ihs:rf:eq:chain}, the eigenvalue with residue
$d_0=1$ is
\begin{equation}\label{ihs:rf:eq:closedlambda}
 \lambda_0=\sqrt{1+4t^2},
\end{equation}
where the square root has residue $1$. Define
\begin{equation}\label{ihs:rf:eq:qdef}
 q=\frac{\lambda_0-1}{2t}=\frac{2t}{\lambda_0+1}.
\end{equation}
Then $v(q)=1$, the intermediate-normalized eigenvector has coordinates
$(v_0)_n=(n+1)q^n$, and the normalized eigenvector is
\begin{equation}\label{ihs:rf:eq:closedvector}
 (u_0)_n=(n+1)q^n\left(\frac{(1-q^2)^3}{1+q^2}\right)^{1/2}.
\end{equation}
All square roots in these formulas are near-one binomial Hahn roots.
\end{proposition}
\begin{proof}
The binomial series defines \eqref{ihs:rf:eq:closedlambda}; it gives
$q=t+\bigo(t^3)$ and the exact identities
\begin{equation}\label{ihs:rf:eq:qrelations}
 t=\frac{q}{1-q^2},\qquad
 \lambda_0=\frac{1+q^2}{1-q^2}=t(q+q^{-1}).
\end{equation}
For $n\ge0$, direct substitution gives
\begin{align*}
 \frac{(n+1)q^n}{n+1}+t\bigl(nq^{n-1}+(n+2)q^{n+1}\bigr)
 &=q^n\bigl(1+t(nq^{-1}+(n+2)q)\bigr)\\
 &=\lambda_0(n+1)q^n.
\end{align*}
For $n=0$, the term with the factor $n$ is zero, agreeing with the boundary
condition. Because $v(q^n)=n$, at any exponent only finitely many vector
coordinates can contribute. The displayed sequence is therefore a Hahn
vector, not merely an arbitrary sequence of elements of $\K$. It has zeroth
coordinate $1$ and residue $e_0$. The eigenspace statement in
\cref{ihs:rf:thm:synthesis}, together with its residue and zeroth
coordinate, identifies it with the unique intermediate-normalized branch
constructed in \cref{ihs:rf:thm:similarity}.

The elementary formal identity
\[
 \sum_{n\ge0}(n+1)^2z^n=\frac{1+z}{(1-z)^3}
\]
can be evaluated at $z=q^2$ by \cref{ihs:cor:binomial}. Hence
\[
 \ip{v_0}{v_0}=\frac{1+q^2}{(1-q^2)^3}.
\]
Taking its inverse near-one square root gives
\eqref{ihs:rf:eq:closedvector}.
\end{proof}

\Cref{ihs:rf:prop:closed} is a directly proved explicit example; no separate
priority claim is made for this solvable tridiagonal model. It also checks
the perturbation construction by a method independent of the matrix
recursion. Its eigenvalue begins
\begin{equation}\label{ihs:rf:eq:chainjet}
 \lambda_0=1+2t^2-2t^4+4t^6-10t^8+28t^{10}+\bigo(t^{12}).
\end{equation}
The included exact computation obtains these same coefficients from the
six-vertex principal section. \Cref{ihs:rf:cor:radius} shows that the
section $0,\ldots,5$ gives the infinite branch correctly through degree
$10$; its first possible error is degree $12$.

\subsection{The first finite-section error is explicit}\label{ihs:rf:sec:chainerror}

For the path section $0,\ldots,L$ with unit edge coefficients, the
coefficient in \cref{ihs:rf:thm:sharp} simplifies to
\[
 \frac1{(1-1/(L+1))\prod_{k=1}^{L-1}(1-1/(k+1))^2}=L(L+1).
\]
Indeed, $1-1/(k+1)=k/(k+1)$ and the product telescopes. Therefore
\begin{equation}\label{ihs:rf:eq:chainerror}
 \lambda_0^{[L]}-\lambda_0^{[L-1]}=L(L+1)t^{2L}+\bigo(t^{2L+2}).
\end{equation}
The infinite branch and the section $0,\ldots,L$ agree through degree $2L$,
so the same first coefficient describes the error made by truncating the
infinite chain at $L-1$.

This is stronger than a qualitative finite-section limit: it specifies the
first scale at which the omitted boundary is visible and gives its nonzero
coefficient. It makes sense without any ordinary numerical choice of a small
real $t$.

\section{A higher-rank example: exact sums without cofinal powers}\label{ihs:sec:ranktwo}

Let $\Gamma=\Z^2$ with lexicographic order, and set $a=(0,1)$, $b=(1,0)$.
Both are positive, but $na<b$ for every positive integer $n$. Consider
\begin{equation}\label{ihs:rf:eq:ranktwo}
 A=\begin{pmatrix}
 0&t^a&0\\
 t^a&1&t^b\\
 0&t^b&3
 \end{pmatrix}.
\end{equation}
The normalized branch at vertex $0$ has the coefficient identities
\begin{equation}\label{ihs:rf:eq:ranktwocoeffs}
 \coeff{2a}{\lambda_0}=-1,\qquad
 \coeff{a+b}{U_{20}}=\frac13,\qquad
 \coeff{2a+2b}{\lambda_0}=-\frac13.
\end{equation}
The first uses the edge $0\leftrightarrow1$. The distant eigenvector
coordinate uses the path $0\to1\to2$, and the last eigenvalue coefficient is
the first return through the distant edge $1\leftrightarrow2$. These
coefficients follow from the recursion, or from the shortest corresponding
walk expressions; the included bivariate calculation checks them exactly.

There can be infinitely many powers of $t^a$ before the scale $t^b$.
Accordingly, a finite total-degree computation in two formal variables is
\emph{not} the same as computing every coefficient below $b$ in the
lexicographic valuation. The verification program checks specified bidegrees
and explicitly does not conflate these truncations.

\begin{proposition}[Why geometric partial sums need not converge]\label{ihs:rf:prop:noconv}
In this Hahn field, the family $((t^a)^n)_{n\ge0}$ is strongly summable with
sum $(1-t^a)^{-1}$, but its finite partial sums do not converge to that sum
in the valuation topology of $\K$.
\end{proposition}
\begin{proof}
\Cref{ihs:lem:neumann} gives the sum and its algebraic identity. The exact
remainder after degree $N$ is
\[
 \frac{t^{(N+1)a}}{1-t^a},
\]
of valuation $(N+1)a<b$. Thus it never belongs to the neighborhood of zero
defined by $v(x)>b$. This prevents convergence to the Hahn sum.
\end{proof}

In a surreal realization sending $a$ to $1$ and $b$ to $\omega$, the
monomials become $\omega^{-1}$ and $\omega^{-\omega}$. The same exact
remainder remains larger in magnitude than the latter infinitesimal scale.
The spectral theorem uses the word-finiteness proof, so it remains valid
despite this failure of sequential convergence.

\section{Three boundaries that cannot be omitted}\label{ihs:sec:failures}

\subsection{No bounded Hilbert-space diagonalizer with these jets}\label{ihs:rf:sec:boundedfail}

\begin{proposition}[The inverse-gap obstruction]\label{ihs:rf:prop:bounded}
On ordinary $\ell^2(\N_{>0})$, let $D=\diag(1/n)$ and let $C$ be the
nearest-neighbor adjacency operator. Both $D$ and $C$ are bounded and
self-adjoint. There is no formal diagonalization of $D+tC$ of the form
\[
 (D+tC)U(t)=U(t)\Lambda(t),\qquad U(t)=\Id+tX+\cdots,
\]
with $X$ a bounded operator and $\Lambda(t)$ diagonal with constant term
$D$. In particular no ordinary operator-norm analytic near-identity unitary
diagonalizer can have that property.
\end{proposition}
\begin{proof}
The order-$t$ off-diagonal equation is $[D,X]+C=\text{diagonal}$. At the
neighboring entry $(n,n+1)$ it gives
\[
 X_{n,n+1}=-\frac1{1/n-1/(n+1)}=-n(n+1).
\]
Every matrix entry of a bounded operator has absolute value at most its
operator norm. The displayed sequence is unbounded, a contradiction. An
operator-norm analytic family would have a bounded derivative at zero, so it
is ruled out as well.
\end{proof}

The diagonalizer of \cref{ihs:rf:thm:unitary} does have this first
coefficient, and that coefficient belongs to $R_I$: it is supported on
neighboring diagonals, even though its real entries are unbounded. Thus the
theorem does not secretly solve the ordinary small-divisor or
bounded-operator problem. It solves a different, explicitly algebraic
problem with exact infinitesimal scales. The same contrast appears from the
other side in \cref{ihs:sec:finitesections}, where the excluded object is an
unbounded diagonal \emph{coefficient} of an adjointable operator.

\subsection{Repeated residues can destroy every eigenvector}\label{ihs:rf:sec:repeated}

Let $T$ be the bilateral adjacency operator on finite-support vectors
indexed by $\Z$:
\[
 (Tx)_n=x_{n-1}+x_{n+1}.
\]
It is real symmetric and row- and column-finite. A nonzero finite-support
vector cannot be an eigenvector of $T$: at the index immediately above the
largest nonzero coordinate, $Tx$ is nonzero while the vector itself is zero.
The same observation shows that $T$ is injective on $\C^{(\Z)}$.

\begin{proposition}[A positive Hahn operator with no eigenvectors]\label{ihs:rf:prop:degenerate}
For any $\eta\in\Gamma_{>0}$, the operator
\[
 A=\Id+t^\eta T\in\RCF_{\Z}(\C)((t^\Gamma))
\]
is self-adjoint and satisfies $\ip{x}{Ax}>0$ for every nonzero
$x\in\C^{(\Z)}((t^\Gamma))$. Nevertheless it has no eigenvalue in $\K$ for
its action on that Hahn vector space.
\end{proposition}
\begin{proof}
If $\gamma=v(x)$, the leading term of $\ip{x}{x}$ has exponent $2\gamma$ and
positive coefficient. The term $t^\eta\ip{x}{Tx}$ has valuation at least
$2\gamma+\eta$. Thus $\ip{x}{Ax}$ has the same positive leading term as the
norm square.

Suppose $Ax=zx$ with $x\ne0$. Injectivity of $T$ on the leading
finite-support coefficient gives $v(Tx)=v(x)=\gamma$. Hence $z\ne1$ and
$v(z-1)=\eta$ in the equation $t^\eta Tx=(z-1)x$. If $c$ is the nonzero
leading coefficient of $t^{-\eta}(z-1)$, comparison at exponent
$\gamma+\eta$ gives $Tx_\gamma=cx_\gamma$. That is a forbidden
finite-support eigenvector of the bilateral adjacency operator. The
contradiction proves the assertion.
\end{proof}

All residue labels in this example are $1$. Thus replacing ``pairwise
distinct ordinary residues'' in \cref{ihs:rf:thm:unitary} merely by
``positive self-adjoint'' would be false. The general phenomenon of positive
infinite row- and column-finite matrices without eigenvalues is already
emphasized in \cite{ago}; the proof here shows explicitly how it survives
the chosen Hahn extension.

\subsection{Entrywise Hahn matrices are not enough}\label{ihs:rf:sec:incoherent}

\begin{proposition}[Failure without global support coherence]\label{ihs:rf:prop:incoherent}
Let $I=\N_0$ and $\Gamma=\Z$. The entrywise Hahn diagonal matrix
$M=\diag(t^{-n})_{n\ge0}$ does not define an operator from $\Hrf$ to itself
by coordinatewise multiplication, although every entry belongs to $\K$ and
every individual exponent coefficient is row- and column-finite.
\end{proposition}
\begin{proof}
Its union of supports is $\{0,-1,-2,\ldots\}$, which is not well ordered;
thus $M\notin\Arf$. More concretely,
\[
 x=\sum_{n\ge0}t^n e_n\in\Hrf
\]
would be sent to the sequence $(1,1,1,\ldots)$. This would require an
infinite-support coefficient at exponent zero, so it is not an element of
$\Hrf$.
\end{proof}

Likewise, a coefficient matrix with an infinite nonzero column can send one
$e_i$ outside $\C^{(I)}$; requiring row-finiteness as well makes the adjoint
remain in the same algebra. These hypotheses are simple sufficient
conditions for the whole framework. The counterexamples do not claim that no
larger useful operator categories exist.

\section{Transport to surreal numbers and extension of scales}\label{ihs:sec:transport}

\begin{corollary}[An infinite surcomplex spectral system]\label{ihs:rf:cor:surcomplex}
Whenever the input is given in a workspace embedded by
\eqref{ihs:eq:surreal-embedding}, every scalar entry of $U$, every
eigenvalue $\lambda_i$, and every scalar coordinate of a vector in $\Hrf$ is
an actual surreal or surcomplex number. All identities in
\cref{ihs:sec:category,ihs:sec:diagonal,ihs:sec:synthesis,ihs:sec:resolvent,ihs:sec:locality}
remain valid after this embedding, with sums interpreted as the transported
strong sums. The eigenvalues of a Hermitian input are surreal, and its inner
products have positive surreal norm squares.
\end{corollary}
\begin{proof}
The embedding preserves every finite coefficient operation and every strong
sum used in the construction. It preserves order and conjugation, so it
preserves positivity, adjoints, and the normalization conditions.
\end{proof}

The matrices and vectors remain set-indexed structured objects; an infinite
matrix is not being identified with one surcomplex scalar. Nor is this a
theorem about every conceivable class-sized operator over $\No[i]$. It
provides a rigorous infinite-dimensional theory in a fixed, explicitly
embedded workspace.

\begin{proposition}[Naturality and absence of new spectral points]\label{ihs:rf:prop:extension}
Let $\iota:\Gamma\hookrightarrow\Gamma'$ be a strictly increasing additive
embedding. Apply it to all Hahn exponents of $A$, keeping the same index set
and complex coefficients. The normalized diagonalizer and eigenvalues in the
$\Gamma'$-workspace are exactly the images of $U$ and $\lambda_i$. Moreover,
\[
 \sigma_{\mathcal A_{I,\Gamma'}}(\iota_*A)=\{\iota_*\lambda_i:i\in I\}.
\]
In particular, enlargement of the value group creates neither new labeled
eigenvalues nor additional algebraic spectral points for this simple-residue
class.
\end{proposition}
\begin{proof}
An order embedding takes well-ordered supports to well-ordered supports and
preserves the finiteness of coefficient incidences. Hence it commutes with
all strong sums, products, divided commutators, and near-one binomial
operations in the proof of \cref{ihs:rf:thm:unitary}. The transported pair
satisfies the same normalization, and uniqueness identifies it with the pair
constructed in the larger workspace. \Cref{ihs:rf:thm:resolvent} applied
there, including to parameters $z\in K_{\Gamma'}$ outside the smaller field,
proves the displayed equality.
\end{proof}

This is an algebraic invariance statement. A scale extension can change
whether a sequence of partial sums converges in the intrinsic valuation
topology, as \cref{ihs:rf:prop:noconv} illustrates. The exact Hahn
identities are stable even when such sequential behavior changes.
\Cref{ihs:hh:cor:extension} is the corresponding statement in
Part~I; it concerns a different construction and is proved separately,
also without divisibility.

\section{Effective coefficient calculation}\label{ihs:rf:sec:computation}

\subsection{What a certified computation needs}\label{ihs:rf:sec:contract}
The proof gives an exact coefficient algorithm once the input presentation
supplies the required finite data. To compute a coefficient at a weight
$\gamma$, one needs a finite list of all ordered input-weight words summing
to $\gamma$, the coefficient matrices on the finitely many vertices those
words can visit, and exact arithmetic for the ordinary labels and edge
coefficients. \Cref{ihs:rf:thm:walks} then reduces the problem to finite
rational operations and finite matrix jets.

For a one-scale locally finite graph and an integer degree $n$, the spatial
part is particularly explicit: use the radius-$\lfloor n/2\rfloor$ ball for
an eigenvalue coefficient, and a sufficiently large radius-$n$ ball for an
eigenvector coefficient. Equations
\eqref{ihs:rf:eq:Trec}--\eqref{ihs:rf:eq:Vrec} compute the intermediate
normalization, after which \eqref{ihs:rf:eq:normalization} computes the
unitary normalization.

For an arbitrary abstract ordered group, the existence of finitely many
weight decompositions does \emph{not} provide an algorithm that finds them
from an order-comparison oracle. Likewise, arbitrary complex coefficients
need not have a decidable exact equality test. The theorem is therefore an
exact mathematical construction with a stated computation contract, not a
universal effective representation of all surreal numbers.

\subsection{A coefficient-level description of the implementation}\label{ihs:rf:sec:impl}
The Part II program uses exact SymPy matrices with rational and Gaussian
rational entries. It implements finite multivariate Taylor jets, indexed by
tuples $\alpha\in\N^r$ of bounded total degree. With
$B=\sum_{\beta\ne0}B_\beta z^\beta$, it computes
\begin{align*}
 T_\alpha&=\sum_{0<\beta\le\alpha}B_\beta V_{\alpha-\beta}
 -\sum_{0<\beta<\alpha}V_{\alpha-\beta}\Delta_\beta,\\
 \Delta_\alpha&=\pi(T_\alpha),\qquad V_\alpha=-\mathcal R_D(T_\alpha).
\end{align*}
The second sum means componentwise $\beta\le\alpha$ with both $\beta$ and
$\alpha-\beta$ nonzero. It next forms $N=V^*V$ and evaluates the truncated
binomial series for $N^{-1/2}$. All coefficients retained are checked
against the full eigenidentity and both unitary identities.

The program does not numerically approximate eigenvalues. It also does not
pretend to enumerate arbitrary Hahn supports or implement infinite matrices.
The infinite claims rest on the proofs in this part; the code tests finite
algebraic instances and the advertised coefficient formulas.

\section{Part II: further questions and a formalization route}\label{ihs:rf:sec:further}

\subsection{Beyond simple ordinary residue labels}\label{ihs:rf:sec:beyondsimple}
The repeated-residue counterexample \cref{ihs:rf:prop:degenerate} shows that
infinite degeneracy cannot be ignored. A natural next class consists of
diagonal residue blocks that are finite, perhaps with unbounded block sizes,
together with an explicit hierarchy of positive splitting scales. A useful
theorem would need a global support certificate ensuring that the block
refinements and their unitary transformations remain coherent across all
blocks. Finite-dimensional splitting by itself does not provide that uniform
certificate.

Another direction replaces ordinary gaps by Hahn-valued gaps. Then dividing
by $d_i-d_j$ can lower valuations by index-dependent amounts. A possible
hypothesis would control the total gap losses along weighted walks so that
the resulting effective weights are still well ordered and have finite word
multiplicity. The present proof does not establish that extension: its
divided commutator is support preserving exactly because the nonzero gaps
are ordinary constants. \Cref{ihs:rf:rem:monoid} is the concrete warning
attached to this point.

\subsection{Larger vector and operator categories}\label{ihs:rf:sec:largercat}
The vector space $\Hrf$ admits only finite-support coefficients. Allowing
ordinary $\ell^2$ coefficients would include many more states, but real
small divisors can then invalidate the recurrence, as
\cref{ihs:rf:prop:bounded} shows. A weighted Hilbert-coefficient or
Banach-coefficient Hahn category would have to replace row/column finiteness
by analytic estimates compatible with the inverse commutator. Such an
extension could connect the exact algebraic results to ordinary spectral
approximation, but it is not a consequence of the current construction.
Part~I is a Hilbert-coefficient theory in that spirit, but for
a different class of operators, and it does not supply the missing
diagonalization theorem either. Part~IV shows that in Part~I's
bounded-coefficient algebra the corresponding global statement is false
even for a trace-class tridiagonal perturbation with simple residues
(\cref{ihs:fr:thm:nounitary}), so an extension of this kind would need
hypotheses beyond simple residues and small coefficient ideals.

Similarly, the marked-walk bounds identify first possible errors, not their
noncancellation in every graph. It is worthwhile to classify graphs and
signs for which the minimum escaping-walk weight always survives in the
coefficient difference. The path formula \eqref{ihs:rf:eq:sharp} proves this
only for a particular family.

\subsection{A proof-assistant decomposition}\label{ihs:rf:sec:formalization}
A formalization can be organized without initially constructing the entire
surreal class. The foundational objects are the set-sized algebra $R_I$, its
involution and finite-support module, followed by their Hahn extensions. The
essential reusable ingredients are: finite row/column support under
products; the inner-product leading-term lemma; ordered-word summability;
coefficientwise divided commutators; the homogeneous diagonalization
recursion; and the uniqueness argument.

After those, normalization, synthesis, and commutant classification are
short algebraic developments. The resolvent theorem needs a separate
coherence lemma for infinitely many inverse diagonal entries, while locality
requires an inductive relation between coefficient expressions and finite
labeled walks. Only the final transport step needs a proved normal-form
embedding into actual surreal numbers. No claim is made that these new
infinite-dimensional statements have been added to, compiled in, or checked
by the repository.

\subsection{Conclusion of Part II}\label{ihs:rf:sec:conclusion}
The principal result is not that every infinite Hermitian surcomplex matrix
has a spectral theorem. That assertion is false in the category defined
here, by \cref{ihs:rf:prop:degenerate}. The result is that simple ordinary
residues and global positive Hahn support provide a complete, exact,
infinite-dimensional spectral mechanism with unusually strong locality and
extension properties. Its infinitesimal series are actual algebraic objects,
its vector expansions are licensed sums, its resolvents are coherent
operators, and its finite-section errors have verifiable scale thresholds.

\clearpage
\part*{Part III. Drazin halos: constant elements of an arbitrary Hahn operator algebra}
\addcontentsline{toc}{part}{Part III. Drazin halos}

\emph{Standing hypotheses for Part III.} $\Alg$ is a nonzero associative
unital complex algebra, with the complex scalars central, and $\Gamma$ is a
nonzero set-sized ordered abelian group. Divisibility, a norm on $\Alg$,
normality and rank one are \emph{not} assumed. Among general results, only
\cref{ihs:dz:sec:banach} and \cref{ihs:dz:prop:polemap} add a Banach norm.
``Spectrum'' in this part means $\SigA(T)$, failure of two-sided
invertibility of $T-\lambda1$ inside the named algebra $\AlgG$, for a
\emph{constant} $T\in\Alg$ (\cref{ihs:dz:def:spectrum}). It is the third
sense recorded in \cref{ihs:tab:senses}. It is a spectrum relative to an
algebra, like $\sigma^{\mathrm{adj}}_\Gamma$ of Part~I and $\sigma_{\Arf}$ of
Part~II, and \cref{ihs:sec:unify} proves that it contains both of them as
special cases for constant elements. It is not a failure-of-bijectivity
spectrum on a chosen vector space.

\emph{Notation.} The source manuscript of this part wrote $A$ for the
algebra, $D$ for a Drazin inverse, $P,Q$ for its idempotents, $B$ for the
Drazin spectrum, $V$ for the Volterra operator, $h$ for a positive-order
correction and $r$ for a ramification index. Each of these letters already
has a different meaning in Parts~I and~II, so this part writes instead
$\Alg$, $S^{\mathrm D}$, $P_{\mathrm{nil}},P_{\mathrm{inv}}$, $Z$,
$\mathcal V$, $\kappa$ and $e_a$. Its $K_\Gamma$ and $\mathfrak m_\Gamma$
are the $\K$ and $\mK$ of this report.

\section{The question, the answer, and its scope}\label{ihs:dz:sec:intro}

\subsection{Why nonnormal constant elements are the next step}\label{ihs:dz:sec:why}
The distinction between a nilpotent operator and a quasinilpotent but
nonnilpotent operator is easy to blur at infinitesimal scales. Over $\C$,
every nonzero scalar displacement of either operator can be invertible.
Over a Hahn field, an exact inverse must satisfy a support condition. The
nilpotent case requires only finitely many negative powers of the
displacement. A nonnilpotent singularity can demand an infinite descending
sequence of exponents, which is not a Hahn support. Small numerical size is
not a replacement for this support requirement.

For constant \emph{normal} operators, \cref{ihs:hh:thm:spectrum} gives the
ordinary spectrum together with infinitesimal halos of the ordinary
accumulation points. This part asks what replaces accumulation points when
normality is dropped and the Hilbert operator algebra is replaced by an
arbitrary unital algebra, and exactly how a nondivisible value group affects
polynomial spectral mapping. The answers are finite-index Drazin
invertibility and a ramification test. In a Banach algebra, an isolated
spectral point can acquire
an entire infinitesimal halo when the ordinary resolvent singularity is not a
finite pole. Conversely, every finite pole permits an inverse at every
nonzero infinitesimal displacement, with an exact valuation determined by its
order. These questions are formulated and answered here; they are not
represented as named open conjectures already posed in published literature.

\subsection{The theorem package at a glance}\label{ihs:dz:sec:glance}
Write $\sigma_\Alg(T)$ for the ordinary algebra-relative spectrum of
$T\in\Alg$ and $\sigD(T)$ for the set of $a\in\C$ such that $T-a1$ is not
Drazin invertible (\cref{ihs:dz:def:drazin}). The first main conclusion,
\cref{ihs:dz:thm:halo}, is
\begin{equation}\label{ihs:dz:eq:preview}
 \SigA(T)=\sigma_\Alg(T)\cup\bigl(\sigD(T)+\mK\bigr).
\end{equation}
The added halos contain all infinitesimals, including nonmonomial ones with
arbitrary well-ordered support. The inverse criterion is unchanged by
enlarging the value group.

For a nonconstant $p\in\C[X]$, set $Z=\sigD(T)$. For $b\in p(Z)$ define the
finite nonempty fiber and the ramification indices
\[
 Z_b=\{a\in Z:p(a)=b\},\qquad
 e_a=\operatorname{ord}_{X=a}(p(X)-b).
\]
The second conclusion, \cref{ihs:dz:thm:image}, is the exact local image
\begin{equation}\label{ihs:dz:eq:imagepreview}
 p(\SigA(T))\cap(b+\mK)=\{b\}\cup
 \Bigl\{b+\zeta:0\ne\zeta\in\mK,\ v(\zeta)\in\bigcup_{a\in Z_b}e_a\Gamma\Bigr\}.
\end{equation}
In contrast, $\SigA(p(T))$ contains the whole halo $b+\mK$. Spectral mapping
therefore holds exactly when, in each relevant fiber, at least one of the
integers $e_a$ acts surjectively on $\Gamma$ (\cref{ihs:dz:thm:criterion}).
This includes divisible groups but gives a sharp answer for every group.

\subsection{Classical ingredients and the proposed novelty}\label{ihs:dz:sec:novelty}
Drazin inverses and their finite-index algebraic identities are classical
\cite{Drazin}. Their relation to resolvent poles, their Banach-algebra
spectra, and ordinary polynomial spectral mapping are also classical
\cite{BoassoOperators,Boasso}. Generalized Drazin invertibility, which permits
a quasinilpotent singular part, is a different condition \cite{Koliha}; only
the finite-index notion is used here.

The candidate contribution is the combined arbitrary-rank classification:
straightening and descent for every infinitesimal in the original Hahn field,
the resulting nonnormal halo formula, and the complete ramification and
value-group description of the polynomial image. Independent coefficient
proofs are given, but re-proving a classical Laurent identity does not make
that identity new. The constant-normal case is an explicit overlap with
Part~I (\cref{ihs:dz:rem:recover}), not a new claim. The Banach algebraicity
criterion (\cref{ihs:dz:cor:algebraic}) is a Hahn interpretation of a
classical Drazin-spectrum characterization. \Cref{ihs:app:auditC} records the
source's repository and literature audit, and corrects the statements of it
that are stale.

\section{Hahn algebras and the third spectrum}\label{ihs:dz:sec:hahn}

\subsection{The coefficient algebra is part of the statement}\label{ihs:dz:sec:algebra}
Set $\AlgG=\Alg((t^\Gamma))$, with coefficientwise addition and convolution
product; \cref{ihs:lem:support}(i) makes it a unital $\K$-algebra with $\K$
central. No topology on $\Alg$ is needed. For nonzero $x\in\AlgG$ write
$v_\Alg(x)=\min\supp x$ and $\lc_\Alg(x)=x_{v_\Alg(x)}$, and put
$v_\Alg(0)=+\infty$; on scalar series write $v$ and $\lc$. Since $\Alg$ may
have zero divisors, in general only $v_\Alg(xy)\ge v_\Alg(x)+v_\Alg(y)$. For
a nonzero central scalar $s\in\K$ and nonzero $x\in\AlgG$, however,
\begin{equation}\label{ihs:dz:eq:scalarvaluation}
 v_\Alg(sx)=v(s)+v_\Alg(x),\qquad
 \lc_\Alg(sx)=\lc(s)\lc_\Alg(x),
\end{equation}
because a nonzero complex number cannot annihilate a nonzero element of
$\Alg$. Every $s\in\Ocal$ is uniquely $a+\eps$ with $a=\st(s)\in\C$ and
$\eps\in\mK$. A scalar of negative valuation is called infinite; the
\emph{halo} of $a\in\C$ is $a+\mK$.

\begin{definition}[The spectrum relative to a Hahn operator algebra]\label{ihs:dz:def:spectrum}
For a constant $T\in\Alg$ put
\[
 \SigA(T)=\{\lambda\in\K:T-\lambda1\text{ has no two-sided inverse in }\AlgG\}.
\]
The ordinary spectrum $\sigma_\Alg(T)$ is defined by the same condition for
$\lambda\in\C$ and inverses in $\Alg$.
\end{definition}

For $\Alg=\B$, convolution gives an action on $H((t^\Gamma))$, and $\AlgG$
is exactly the adjointable algebra $\Ahh$ of Part~I. The definition is
\emph{not} silently replaced by failure of bijectivity in the algebra of all
$\K$-linear endomorphisms, or by the spectrum of row- and column-finite
matrices on a finite-coordinate module; \cref{ihs:sec:unify} states exactly
which identifications hold. In the examples of Volterra and unbounded Jordan
lengths below, actual vector nonsurjectivity is proved directly.

\subsection{The support calculus used}\label{ihs:dz:sec:support}
All infinite sums in this part are strong Hahn sums in the sense of
\cref{ihs:sec:supportconv}. The positive-support lemma is
\cref{ihs:lem:support}; the source manuscript used it in the combined form:
\emph{if $S\subseteq\Gamma_{>0}$ and $B\subseteq\Gamma$ are well ordered,
then $B+S^*$ is well ordered and each exponent has only finitely many
representations as an element of $B$ plus an ordered word in $S$.} This is
\cref{ihs:lem:support}(i) applied to $B$ and $S^*$, followed by
\cref{ihs:lem:support}(ii) inside $S^*$; it is classical, not a new result
here. If $v_\Alg(h)>0$ then $\sum_{n\ge0}h^n$ inverts $1-h$
(\cref{ihs:lem:neumann}), and for a scalar $h$ of positive order and
$q\in\Q$ the binomial series $(1+h)^q=\sum_n\binom qnh^n$ is strongly defined
and obeys the exponent law (\cref{ihs:cor:binomial}).

\begin{warning}[No hidden rank-one assumption]\label{ihs:dz:warn:rank}
This part never assumes that $n\eta$ eventually exceeds every element of
$\Gamma$ for $\eta>0$; that fails at higher rank. All infinite sums are
justified by well ordering and coefficientwise local finiteness, not by
convergence of finite subsums in a valuation topology.
\end{warning}

\section{Drazin decomposition and finite Laurent tails}\label{ihs:dz:sec:laurent}

\begin{definition}[Finite-index Drazin invertibility]\label{ihs:dz:def:drazin}
An element $S\in\Alg$ is \emph{Drazin invertible} if there exist
$S^{\mathrm D}\in\Alg$ and an integer $\nu\ge0$ with
\begin{equation}\label{ihs:dz:eq:drazin}
 SS^{\mathrm D}=S^{\mathrm D}S,\qquad
 S^{\mathrm D}SS^{\mathrm D}=S^{\mathrm D},\qquad
 S^{\nu+1}S^{\mathrm D}=S^\nu.
\end{equation}
The least such $\nu$ is the Drazin index; index zero means ordinary
invertibility. The \emph{Drazin spectrum} is
$\sigD(T)=\{a\in\C:T-a1\text{ is not Drazin invertible}\}$; it is contained
in $\sigma_\Alg(T)$.
\end{definition}

Given the identities, put $P_{\mathrm{nil}}=1-SS^{\mathrm D}$ and
$P_{\mathrm{inv}}=SS^{\mathrm D}$. Indeed,
$(SS^{\mathrm D})^2=S(S^{\mathrm D}SS^{\mathrm D})=SS^{\mathrm D}$, so these
are complementary idempotents commuting with $S$. The second Drazin identity
gives $S^{\mathrm D}P_{\mathrm{nil}}=P_{\mathrm{nil}}S^{\mathrm D}=0$, and
the third gives $S^\nu P_{\mathrm{nil}}=0$. In the corner algebra
$P_{\mathrm{inv}}\Alg P_{\mathrm{inv}}$, with identity $P_{\mathrm{inv}}$,
the element $SP_{\mathrm{inv}}$ is invertible with inverse $S^{\mathrm D}$.
Thus $P_{\mathrm{nil}}$ selects the nilpotent part and $P_{\mathrm{inv}}$ the
invertible part; the corresponding elements are $SP_{\mathrm{nil}}$ and
$SP_{\mathrm{inv}}$. Either corner may be zero. If $\nu\ge1$ is least,
then $S^{\nu-1}P_{\mathrm{nil}}\ne0$: otherwise
$S^\nu S^{\mathrm D}=S^{\nu-1}$ would lower the index by one.
No centrality of $P_{\mathrm{nil}}$ in the whole algebra is required.

Let $u$ be a central formal variable; the ordinary Laurent algebra $\Alg((u))$
permits only finitely many negative integer exponents.

\begin{lemma}[Laurent inverse criterion]\label{ihs:dz:lem:laurent}
For $S\in\Alg$, the pencil $S-u1$ is invertible in $\Alg((u))$ if and only if
$S$ is Drazin invertible. If $S^{\mathrm D}$ and
$P_{\mathrm{nil}}=1-SS^{\mathrm D}$ are its data of index $\nu$, then
\begin{equation}\label{ihs:dz:eq:laurent}
 (S-u1)^{-1}=\sum_{n\ge0}u^n(S^{\mathrm D})^{n+1}
            -\sum_{j=0}^{\nu-1}u^{-j-1}S^jP_{\mathrm{nil}}.
\end{equation}
The second sum is empty for $\nu=0$; otherwise its exact pole order is $\nu$.
\end{lemma}
\begin{proof}
Assume the Drazin identities. Multiplying the first sum by $S-u1$ on either
side gives $P_{\mathrm{inv}}$: its constant term is $SS^{\mathrm D}$ and the
higher coefficients cancel. Multiplying the negative sum gives
$P_{\mathrm{nil}}-u^{-\nu}S^\nu P_{\mathrm{nil}}=P_{\mathrm{nil}}$. The
combined result is $1$. Its coefficient at $u^{-\nu}$ is
$-S^{\nu-1}P_{\mathrm{nil}}\ne0$ when $\nu\ge1$ is least.

Conversely, suppose $B(u)=\sum_{n\ge-m}B_nu^n$ is a two-sided inverse, with
$m\ge0$ and absent coefficients set to zero. Coefficient comparison in both
products gives
\begin{equation}\label{ihs:dz:eq:recurrence}
 SB_n-B_{n-1}=\delta_{n0}1=B_nS-B_{n-1}\qquad(n\in\Z).
\end{equation}
In particular each $B_n$ commutes with $S$. Set $S^{\mathrm D}=B_0$ and
$P_{\mathrm{nil}}=-B_{-1}=1-SS^{\mathrm D}$. If $m=0$, then
$SS^{\mathrm D}=S^{\mathrm D}S=1$. If $m\ge1$, the negative recurrences,
including the finite lower bound, give $S^mP_{\mathrm{nil}}=0$, and the
positive recurrences give $S^{\mathrm D}=S^mB_m$. Since $P_{\mathrm{nil}}$
commutes with $S$, this gives $P_{\mathrm{nil}}S^{\mathrm D}=0$; the
definition of $P_{\mathrm{nil}}$ and commutation give
$S^{\mathrm D}P_{\mathrm{nil}}=0$. More explicitly,
$P_{\mathrm{nil}}B_0=P_{\mathrm{nil}}S^mB_m=0$ and
$B_0P_{\mathrm{nil}}=P_{\mathrm{nil}}B_0$, since $B_0$ commutes with $S$.
It follows that
$P_{\mathrm{nil}}^2=P_{\mathrm{nil}}$, $S^{\mathrm D}SS^{\mathrm D}=S^{\mathrm D}$
and $S^{m+1}S^{\mathrm D}=S^m$, which are Drazin identities. Any Drazin data
give \eqref{ihs:dz:eq:laurent}; uniqueness of a two-sided inverse therefore
identifies its constant coefficient with $S^{\mathrm D}$ and proves
uniqueness of the Drazin inverse.
\end{proof}

This is the classical finite-Laurent mechanism, written out to expose the
exact obstruction that arbitrary Hahn supports must not be allowed to
conceal. No complex-analytic convergence was used.

\section{Straightening and inverse descent at arbitrary rank}\label{ihs:dz:sec:descent}

\begin{lemma}[Strong parameter straightening]\label{ihs:dz:lem:straighten}
Given $0\ne\eps\in\mK$ and $\eta=v(\eps)$, there is a strong,
valuation-preserving $\C$-algebra automorphism $\Phi$ of $\K$ with
$\Phi(t^\eta)=\eps$. The same construction is an $\Alg$-fixing automorphism
of $\AlgG$. No divisibility or enlargement of $\Gamma$ is required.
\end{lemma}
\begin{proof}
Write $\eps=c\,t^\eta(1+\kappa)$ with $c\in\C^*$ and $v(\kappa)>0$. Since
$\Gamma$ is torsion-free, it embeds in $\Gamma\otimes_\Z\Q$. Choose a
rational linear functional $\ell$ on that vector space with $\ell(\eta)=1$;
its restriction to $\Gamma$ is a homomorphism to $\Q$, not necessarily
order-preserving. Choose $L\in\C$ with $e^L=c$ and set $c^q=e^{qL}$ for
$q\in\Q$; these are compatible rational powers of the one constant $c$.
Define on monomials
\begin{equation}\label{ihs:dz:eq:straighten}
 \Phi(t^\gamma)=c^{\ell(\gamma)}t^\gamma(1+\kappa)^{\ell(\gamma)},
\end{equation}
and extend coefficientwise using strong sums. Let $S=\supp\kappa$. For an
input support $B$, every resulting exponent lies in $B+S^*$, so by the
combined support lemma of \cref{ihs:dz:sec:support} each output coefficient
has finitely many contributors. The same bound works simultaneously for any
strongly summable family of inputs, so the map is well defined and strong,
including for uncountable supports. Additivity of $\ell$ and the binomial
exponent law make it multiplicative on monomials, and the same support bounds
justify all rearrangements, hence multiplication of arbitrary series. This
remains true for noncommuting coefficients in $\Alg$, since every binomial
factor is central. The leading coefficient of a nonzero input $x$ is
multiplied by $c^{\ell(v_\Alg(x))}\ne0$, so valuation is preserved.

Surjectivity needs an argument: a triangular-looking map is not surjective
merely because its leading terms are nonzero. Write $\Phi=\Psi\circ C_c$,
where $C_c(t^\gamma)=c^{\ell(\gamma)}t^\gamma$ has an evident inverse and
$\Psi(t^\gamma)=t^\gamma(1+\kappa)^{\ell(\gamma)}$. Put $\Xi=\Psi-\mathrm{id}$.
One application of $\Xi$ shifts a starting exponent by a sum of at least one
member of $S$, so a term of $\Xi^n(x)$ is shifted by a word of length at
least $n$ in $S$. Hence
\begin{equation}\label{ihs:dz:eq:operatorinverse}
 \Theta(x)=\sum_{n\ge0}(-\Xi)^n(x)
\end{equation}
is a strong sum: at a fixed final exponent there are finitely many starting
exponents and words, and each word has only finitely many partitions into
nonempty blocks, so only finitely many iterations and intermediate terms
contribute. The same reasoning shows that $\Theta$ preserves strongly
summable families. Coefficientwise telescoping gives
$(\mathrm{id}+\Xi)\Theta=\Theta(\mathrm{id}+\Xi)=\mathrm{id}$. Hence $\Psi$ is
bijective, and its inverse is an algebra homomorphism. If $\kappa=0$, take
$\Xi=0$. Finally \eqref{ihs:dz:eq:straighten} sends $t^\eta$ to
$c\,t^\eta(1+\kappa)=\eps$ and fixes every coefficient in $\Alg$.
\end{proof}

The rational functional chooses a coordinate on the rational span of the
exponents; it does not adjoin any fractional exponent to the field. The
binomial corrections use only exponents already generated by
$\supp\kappa$. This is why nondivisibility causes no obstruction to
straightening a parameter that already exists. Related strong
Hahn-automorphism constructions belong to the surrounding literature
\cite{KaplanKrappSerra}; the general idea is not claimed as new here.

\begin{lemma}[Cyclic-support inverse descent]\label{ihs:dz:lem:projection}
For $S\in\Alg$ and $0<\eta\in\Gamma$, put $u=t^\eta$. The pencil $S-u1$ is
invertible in $\AlgG$ if and only if it is invertible in $\Alg((u))$. If it
is invertible, its inverse in $\AlgG$ already belongs to the embedded
Laurent subalgebra.
\end{lemma}
\begin{proof}
Retain the coefficients on the cyclic subgroup $\Z\eta$:
$\varpi_\eta\bigl(\sum B_\gamma t^\gamma\bigr)=\sum_{n\in\Z}B_{n\eta}u^n$.
The intersection of a well-ordered support with $\Z\eta$ is well ordered and
becomes, under $n\mapsto n\eta$, a subset of $\Z$ with a finite lower bound;
thus $\varpi_\eta B\in\Alg((u))$. Because the pencil is supported on
$\Z\eta$, coefficient comparison gives
$\varpi_\eta((S-u1)B)=(S-u1)\varpi_\eta B$ and
$\varpi_\eta(B(S-u1))=(\varpi_\eta B)(S-u1)$. A two-sided inverse therefore
projects to a two-sided Laurent inverse. Conversely a Laurent inverse embeds
into $\AlgG$, since a lower-bounded set of integer multiples of a positive
element is well ordered. Uniqueness of a two-sided inverse identifies it
with its projection.
\end{proof}

The map $\varpi_\eta$ is not generally a ring homomorphism: if
$\eta/2\in\Gamma$, it sends $t^{\eta/2}$ to zero but its square to $u$. The
proof used only left and right compatibility with the particular cyclically
supported pencil, and no claim that $n\eta$ is cofinal in $\Gamma$.
The restriction to $\Z\eta$ is essential to the finite negative Laurent
tail: in lexicographic $\Z\times\Z$, for example, the well-ordered set
$\{(-1,n):n\ge0\}$ contains infinitely many negative exponents. General
Hahn supports need not have only finitely many negative terms.

\begin{theorem}[Arbitrary-parameter inverse descent]\label{ihs:dz:thm:descent}
For $S\in\Alg$ and $0\ne\eps\in\mK$, the following are equivalent:
\begin{enumerate}[label=\textup{(\roman*)}]
\item $S-\eps1$ has a two-sided inverse in $\AlgG$;
\item its inverse is a Laurent series in $\eps$ with coefficients in $\Alg$
 and only finitely many negative powers;
\item $S$ is Drazin invertible in $\Alg$.
\end{enumerate}
When they hold, the unique inverse is
\begin{equation}\label{ihs:dz:eq:epsinverse}
 (S-\eps1)^{-1}=\sum_{n\ge0}\eps^n(S^{\mathrm D})^{n+1}
                -\sum_{j=0}^{\nu-1}\eps^{-j-1}S^jP_{\mathrm{nil}}.
\end{equation}
\end{theorem}
\begin{proof}
Apply \cref{ihs:dz:lem:straighten} with $\eta=v(\eps)$. It fixes $S$ and
takes $t^\eta$ to $\eps$, so it preserves the invertibility question.
\Cref{ihs:dz:lem:projection} descends an inverse of $S-t^\eta1$ to its
Laurent subalgebra, and \cref{ihs:dz:lem:laurent} identifies the condition as
Drazin invertibility and supplies the inverse. Applying the strong
automorphism to that formula gives \eqref{ihs:dz:eq:epsinverse}; its
nonnegative part is a strong sum by the support lemma and its negative part
is finite.
\end{proof}

The conclusion rules out more than one unsuccessful expansion: no inverse
with exotic extra exponents can exist when the finite-Laurent obstruction
remains. Although the straightening automorphism involved choices,
existence and the inverse formula are intrinsic.

\section{Spectral halos and exact resolvent scales}\label{ihs:dz:sec:halo}

\begin{theorem}[Constant-element spectrum]\label{ihs:dz:thm:halo}
For every $T\in\Alg$,
\begin{equation}\label{ihs:dz:eq:halo}
 \SigA(T)=\sigma_\Alg(T)\cup\bigl(\sigD(T)+\mK\bigr).
\end{equation}
No infinite scalar is spectral. A constant $a\in\C$ is spectral exactly when
$a\in\sigma_\Alg(T)$. For $0\ne\eps\in\mK$, the scalar $a+\eps$ is spectral
exactly when $a\in\sigD(T)$.
\end{theorem}
\begin{proof}
If $v(\lambda)<0$, then $\lambda^{-1}T$ has positive order unless $T=0$, and
the strong geometric series gives
\begin{equation}\label{ihs:dz:eq:infiniteinverse}
 (T-\lambda1)^{-1}=-\lambda^{-1}\sum_{n\ge0}\lambda^{-n}T^n;
\end{equation}
for $T=0$ the formula has just its first term. If $\lambda=a\in\C$ and $B$
is a two-sided Hahn inverse, its constant coefficient satisfies
$(T-a1)B_0=B_0(T-a1)=1$, because $T-a1$ is constant; so the ordinary
invertibility question is unchanged by Hahn extension. Every other scalar is
uniquely $a+\eps$ with $0\ne\eps\in\mK$, and \cref{ihs:dz:thm:descent}
applied to $S=T-a1$ decides exactly this case. Since
$\sigD(T)\subseteq\sigma_\Alg(T)$, the halo centres are included correctly.
\end{proof}

\begin{theorem}[Exact valuation and leading coefficient]\label{ihs:dz:thm:scale}
Suppose $S=T-a1$ is Drazin invertible of index $\nu\ge1$. For
$0\ne\eps\in\mK$ put $\eta=v(\eps)$ and $c=\lc(\eps)$. Then
\begin{equation}\label{ihs:dz:eq:scale}
 v_\Alg\bigl((T-a1-\eps1)^{-1}\bigr)=-\nu\eta,\qquad
 \lc_\Alg\bigl((T-a1-\eps1)^{-1}\bigr)=-c^{-\nu}S^{\nu-1}P_{\mathrm{nil}}.
\end{equation}
For index zero the inverse has valuation zero and leading coefficient
$S^{-1}$. For infinite $\lambda$ the inverse has valuation $-v(\lambda)$ and
leading coefficient $-\lc(\lambda)^{-1}1$.
\end{theorem}
\begin{proof}
In \eqref{ihs:dz:eq:epsinverse} the last nonzero negative term is
$-\eps^{-\nu}S^{\nu-1}P_{\mathrm{nil}}$, of valuation $-\nu\eta$ and with the
stated nonzero leading coefficient; every other term has strictly larger
valuation, so no cancellation can affect it. The other cases follow from the
initial term of the corresponding geometric series.
\end{proof}

This is a valuation of an operator-valued Hahn series, not a statement about
the least field-valued operator-norm bound of \cref{ihs:hh:sec:norm}. The
integer $\nu$ measures the largest nilpotent chain in the singular part.
Isolation of a spectral point does not by itself determine it.

\begin{corollary}[Workspace invariance]\label{ihs:dz:cor:workspace}
Let $\Gamma\hookrightarrow\Lambda$ be an order-preserving group embedding.
For every constant $T\in\Alg$ and every $\lambda\in K_\Gamma$, invertibility of
$T-\lambda1$ in $\Alg((t^\Gamma))$ is equivalent to invertibility of its
image in $\Alg((t^\Lambda))$; when an inverse exists, the latter is the image
of the former.
\end{corollary}
\begin{proof}
The embedding preserves the sign of the valuation, the constant coefficient,
and nonzeroness of the infinitesimal part, hence every case of
\cref{ihs:dz:thm:halo}. The explicit inverse formulas already exist in the
smaller workspace; uniqueness identifies them after extension.
\end{proof}

Adjoining fractional exponents cannot repair a non-Drazin constant pencil. It
can repair a scalar root problem; \cref{ihs:dz:sec:poly} shows why these two
facts are compatible with a change in the image of a spectrum under a
polynomial. \Cref{ihs:hh:cor:extension,ihs:rf:prop:extension} are the
corresponding statements of Parts~I and~II. Invariance concerns the status
of each old scalar: the halos themselves grow with the field, and
\cref{ihs:fr:cor:extension} of Part~IV makes this explicit for a nonconstant
operator, whose eigenvalues do not change while its non-eigenvalue monad
grows.

\section{Banach algebras: poles, nonnormality and algebraicity}\label{ihs:dz:sec:banach}

This section adds a Banach norm on $\Alg$; the preceding results did not
need one. In a Banach algebra, finite-index Drazin invertibility has a
familiar analytic meaning: a spectral point is Drazin invertible exactly when
it is a finite pole of the ordinary resolvent \cite{BoassoOperators,Boasso}.

To see the mechanism, take $S=T-a1$ with Drazin inverse $S^{\mathrm D}$ and
use the local resolvent convention $R(u)=(S-u1)^{-1}$. The
positive part of \eqref{ihs:dz:eq:laurent} converges in norm for
sufficiently small complex $u$, being a geometric series in $uS^{\mathrm D}$;
its negative part is finite. It supplies the ordinary resolvent on a
punctured disk, with pole order equal to the Drazin index if $a$ is
spectral. Conversely, the coefficients of a finite-pole Laurent expansion of
the analytic resolvent satisfy \eqref{ihs:dz:eq:recurrence}, so its constant
coefficient is a Drazin inverse. With the opposite convention
$(u1-S)^{-1}$ that coefficient would instead be $-S^{\mathrm D}$.
No accumulation point of the ordinary
spectrum can be a finite pole, and an isolated spectral singularity cannot be
removable, since continuity of the inverse identity would then invert
$T-a1$. Consequently
\begin{equation}\label{ihs:dz:eq:dictionary}
 \sigD(T)=\acc\sigma_\Alg(T)\cup
 \{a\in\sigma_\Alg(T):a\text{ is isolated and the resolvent has no finite
 pole at }a\}.
\end{equation}
Here $\acc$ is the ordinary accumulation set, written $\sigma_\C(T)'$ in
Part~I. For a normal operator an isolated spectral point is a simple pole.
Without normality, an isolated essential resolvent singularity is an
additional source of a whole infinitesimal halo. Generalized Drazin
invertibility is not an adequate substitute: it permits a quasinilpotent
singular corner \cite{Koliha}, whereas \cref{ihs:dz:thm:descent} requires
that corner to be nilpotent of \emph{finite} index. The Volterra example of
\cref{ihs:dz:sec:volterra} distinguishes them.

\begin{corollary}[Recovery of the normal case]\label{ihs:dz:cor:normal}
For a bounded normal operator $T$ on a nonzero ordinary complex Hilbert space
$H$, with $\Alg=\B$,
\[
 \Sigma^{\B}_\Gamma(T)=\sigma_\C(T)\cup\bigl(\sigma_\C(T)'+\mK\bigr).
\]
Every isolated ordinary spectral point has Drazin index one, regardless of
the dimension of its eigenspace.
\end{corollary}
\begin{proof}
The ordinary spectral theorem decomposes the space at an isolated normal
spectral point into its eigenspace and orthogonal complement; on the former
$T-aI$ is zero, on the latter boundedly invertible
(\cref{ihs:hh:lem:normalgap}). Apply \eqref{ihs:dz:eq:dictionary} and
\cref{ihs:dz:thm:halo}.
\end{proof}

\begin{remark}[What is recovered from Part~I, and which hypotheses are dropped]\label{ihs:dz:rem:recover}
Since $\B((t^\Gamma))=\Ahh$, the left side of \cref{ihs:dz:cor:normal} is
exactly $\sigma^{\mathrm{adj}}_\Gamma(T)$ of \cref{ihs:hh:sec:twodefs}.
The corollary is therefore the $\sigma^{\mathrm{adj}}$ half of
\cref{ihs:hh:thm:spectrum}, obtained by a different route: straightening and
Laurent descent instead of the defect splitting. At a Drazin point of index
one, \eqref{ihs:dz:eq:epsinverse} with $S=T-aI$ reads
$-\eps^{-1}P_{\mathrm{nil}}+\sum_{n\ge0}\eps^n(S^{\mathrm D})^{n+1}$; with
the notation of \cref{ihs:hh:lem:normalgap}, $P_{\mathrm{nil}}=P$ and
$S^{\mathrm D}$ is the inverse of $(T-aI)|_M$ composed with $Q$, and this is
the isolated-point inverse \eqref{ihs:hh:eq:isolatedres}. Relative to that theorem,
\cref{ihs:dz:thm:halo} drops \emph{normality} (the accumulation set is
replaced by the Drazin spectrum) and the \emph{Hilbert and Banach structure}
(any unital complex algebra; a norm enters only to identify $\sigD$ through
\eqref{ihs:dz:eq:dictionary}). It does \emph{not} drop divisibility relative
to the present Part~I, which has not assumed divisibility since the revision
recorded in \cref{ihs:hh:sec:field}; the source manuscript's claim to drop it
refers to an earlier state of Part~I (\cref{ihs:app:staleC}).

Two parts of \cref{ihs:hh:thm:spectrum} are \emph{not} recovered. The first
is the failure-of-bijectivity spectrum $\sigma^{\mathrm{alg}}_\Gamma(T)$ on
$\Hil$: Part~III is a statement about invertibility in an algebra and says
nothing about bijectivity on a vector space, except in its explicitly worked
examples (the Volterra operator of \cref{ihs:dz:sec:volterra}, and the block
operator of \cref{ihs:dz:sec:blocks}, both at every infinitesimal
displacement). The second is point-spectrum rigidity,
\cref{ihs:hh:thm:point}. \Cref{ihs:hh:thm:spectrum} proves
$\sigma^{\mathrm{alg}}_\Gamma(T)=\sigma^{\mathrm{adj}}_\Gamma(T)$ for
constant normal $T$; together with \cref{ihs:dz:cor:normal} this gives the
full normal formula by two proofs of its algebra half. For nonnormal $T$ one
always has $\sigma^{\mathrm{alg}}_\Gamma(T)\subseteq
\Sigma^{\B}_\Gamma(T)$, since an inverse in $\Ahh$ is a bijection; equality is
not asserted in general (\cref{ihs:dz:sec:questions}).
\end{remark}

\begin{corollary}[Algebraicity criterion]\label{ihs:dz:cor:algebraic}
For an element $T$ of a Banach algebra $\Alg$ the following are equivalent:
\begin{enumerate}[label=\textup{(\roman*)}]
\item $\SigA(T)\subseteq\C$;
\item $\sigD(T)=\varnothing$;
\item $T$ is algebraic over $\C$.
\end{enumerate}
In that case $\SigA(T)=\sigma_\Alg(T)$ is finite.
\end{corollary}
\begin{proof}
The first equivalence follows from \cref{ihs:dz:thm:halo} and the existence
of a positive element of $\Gamma\ne0$. For completeness, the classical second
equivalence can be proved as follows. If $\sigD(T)$ is empty, all points of
the ordinary spectrum are isolated poles, so the compact spectrum is finite,
say $\{a_1,\ldots,a_s\}$. Its finite Riesz decomposition gives commuting
complementary projections $P_j$ with sum one and
$(T-a_j1)^{\nu_j}P_j=0$ for the finite pole orders $\nu_j$; hence
$\prod_j(T-a_j1)^{\nu_j}=0$. Conversely, factor an annihilating polynomial
into powers of distinct linear factors. Polynomial Chinese remaindering gives
complementary idempotents polynomial in $T$; on each part $T$ is a scalar
plus a finite-index nilpotent, so every scalar shift is a direct sum of an
invertible and a finite-index nilpotent part, hence Drazin invertible. The
empty-Drazin-spectrum characterization itself is classical \cite{Boasso};
its exact Hahn spectral interpretation is the consequence asserted here.
\end{proof}

The Banach assumption is essential for the algebraicity implication. For
$\Alg=\C(X)$ and $T=X$, every $T-a$ is invertible, both ordinary and Drazin
spectra are empty, yet $T$ is transcendental over $\C$. No algebraicity claim
is made for arbitrary unital complex algebras. For an ordinary bounded
self-adjoint $T$, being algebraic is the same as having finite spectrum, so
\cref{ihs:dz:cor:algebraic} is consistent with the spectral-reality criterion
\cref{ihs:hh:cor:reality} of Part~I; the two statements concern different
spectra ($\sigma^{\mathrm{adj}}$ here, $\sigma^{\mathrm{alg}}$ there), which
agree for constant normal operators by \cref{ihs:hh:thm:spectrum}.

\section{Ramified polynomial spectral mapping}\label{ihs:dz:sec:poly}

We return to an arbitrary nonzero unital complex algebra $\Alg$ and a
nonconstant $p\in\C[X]$. A nondivisible value group obstructs scalar roots;
the following lemmas turn that obstruction into an exact spectral-mapping
theorem.

\begin{lemma}[Power image]\label{ihs:dz:lem:powers}
For every integer $e\ge1$,
\begin{equation}\label{ihs:dz:eq:powers}
 \{x^e:x\in\mK\}=\{0\}\cup\{\zeta\in\mK\setminus\{0\}:v(\zeta)\in e\Gamma\}.
\end{equation}
\end{lemma}
\begin{proof}
Necessity follows from $v(x^e)=ev(x)$. Conversely write
$\zeta=d\,t^{e\gamma}(1+k)$ with $d\in\C^*$ and $v(k)>0$. Since
$e\gamma>0$, also $\gamma>0$. Choose a complex root $d^{1/e}$ and put
$x=d^{1/e}t^\gamma(1+k)^{1/e}$, by the strong binomial series. Then
$x\in\mK$ and $x^e=\zeta$.
\end{proof}

\begin{lemma}[Exact local polynomial image]\label{ihs:dz:lem:local}
Suppose $p(a)=b$ with $a,b\in\C$, and put $e=\operatorname{ord}_{X=a}(p(X)-b)$.
There is a bijection $\varphi_a:\mK\to\mK$, induced by a formal complex power
series with nonzero linear coefficient, such that
$p(a+x)-b=\varphi_a(x)^e$ for $x\in\mK$. Consequently
$p(a+\mK)-b=\{x^e:x\in\mK\}$.
\end{lemma}
\begin{proof}
Factor $p(a+X)-b=cX^e(1+Xg(X))$ with $c\in\C^*$ and $g\in\C[X]$. Choose
$c^{1/e}$ and put $\varphi_a(X)=c^{1/e}X(1+Xg(X))^{1/e}\in X\C[[X]]$. Its
nonzero linear coefficient gives a unique formal compositional inverse
$\psi_a\in X\C[[X]]$: at each degree the next unknown coefficient is
multiplied by that nonzero linear coefficient and all other terms are
already determined. Every series and composition in these identities can be
evaluated at any element of $\mK$ by \cref{ihs:cor:binomial}; the values stay
in $\mK$, and the identities make the resulting maps mutually inverse on the
whole of $\mK$. The defining $e$th-power identity proves the claim. No
ordinary convergence radius was used.
For $x\ne0$ the nonzero linear coefficient is the unique leading term in
each substitution, so $v(\varphi_a(x))=v(x)=v(\psi_a(x))$. In particular,
both substitutions preserve nonzeroness as well as positive valuation.
\end{proof}

A ramification index is thus an exact condition on the value group, not an
approximate estimate along one selected scale; the leading coefficient causes
no additional obstruction because $\C$ contains its roots.

\begin{lemma}[Commuting factors]\label{ihs:dz:lem:commuting}
A finite product of commuting elements in a unital algebra is a unit if and
only if every factor is a unit.
\end{lemma}
\begin{proof}
If the product $q$ is a unit, each factor commutes with $q$ and hence with
$q^{-1}$, and the product of all other factors times $q^{-1}$ is both a left
and a right inverse of that factor. The converse is immediate.
\end{proof}

\begin{proposition}[Ordinary polynomial mapping]\label{ihs:dz:prop:ordinarymapping}
For every nonconstant $p\in\C[X]$,
\begin{equation}\label{ihs:dz:eq:ordinarymapping}
 \sigma_\Alg(p(T))=p(\sigma_\Alg(T)),\qquad \sigD(p(T))=p(\sigD(T)).
\end{equation}
\end{proposition}
\begin{proof}
Factor $p(X)-b$ over $\C$ and apply \cref{ihs:dz:lem:commuting} at $T$ for
the first equality. For the second, work temporarily in the auxiliary field
$\C((u^\Q))$; this imposes no divisibility assumption on the workspace of the
later theorems. Let $a_1,\ldots,a_s$ be the distinct roots of $p(X)-b$, with
multiplicities $e_1,\ldots,e_s$. For every $e_j$th root of unity $\xi$,
\cref{ihs:dz:lem:local} gives a root $\alpha_{j,\xi}=a_j+\psi_{a_j}(\xi
u^{1/e_j})$ of $p(X)-b-u$. These roots are distinct: different $a_j$ have
different residues, and for the same $a_j$ the series $\psi_{a_j}$ is
injective on infinitesimals. Their number is $\sum_je_j=\deg p$, so with
$c_p$ the leading coefficient of $p$,
\[
 p(X)-b-u=c_p\prod_{j=1}^s\prod_{\xi^{e_j}=1}(X-\alpha_{j,\xi}).
\]
After evaluation at $T$ all factors commute. By \cref{ihs:dz:thm:descent}, a
factor is a unit exactly when $T-a_j1$ is Drazin invertible, and the product
is a unit exactly when $p(T)-b1$ is Drazin invertible. The commuting-factor
lemma proves that this happens exactly when every $T-a_j1$ is Drazin
invertible, which is the second equality.
\end{proof}

\begin{theorem}[Exact ramified spectral image]\label{ihs:dz:thm:image}
Let $Z=\sigD(T)$ and, for $b\in p(Z)$, let $Z_b$ and $e_a$ be as in
\cref{ihs:dz:sec:glance}; these multiplicities refer to the polynomial germ,
not to the multiplicity of an operator eigenvalue. Then for $b\in p(Z)$,
\begin{equation}\label{ihs:dz:eq:image}
 p(\SigA(T))\cap(b+\mK)=\{b\}\cup
 \Bigl\{b+\zeta:0\ne\zeta\in\mK,\ v(\zeta)\in\bigcup_{a\in Z_b}e_a\Gamma\Bigr\}.
\end{equation}
For $b\in p(\sigma_\Alg(T))\setminus p(Z)$ this intersection is $\{b\}$; for
$b\notin p(\sigma_\Alg(T))$ it is empty. In contrast,
\begin{equation}\label{ihs:dz:eq:polyspectrum}
 \SigA(p(T))=p(\sigma_\Alg(T))\cup\bigl(p(Z)+\mK\bigr).
\end{equation}
In particular $p(\SigA(T))\subseteq\SigA(p(T))$, and these formulas describe
every possible omission from the polynomial image.
\end{theorem}
\begin{proof}
By \cref{ihs:dz:thm:halo}, a spectral point is either an ordinary point of
$\sigma_\Alg(T)$ or a member of a halo $a+\mK$ with $a\in Z$. The standard
part of $p(a+\eps)$ is $p(a)$, so contributions to a halo $b+\mK$ come only
from its ordinary fibre. For each $a\in Z_b$,
\cref{ihs:dz:lem:local,ihs:dz:lem:powers} give exactly the condition
$v(\zeta)\in e_a\Gamma$. Their finite union proves \eqref{ihs:dz:eq:image},
including the cases with no Drazin-spectral preimage. Finally apply
\cref{ihs:dz:thm:halo} to $p(T)$ and substitute
\cref{ihs:dz:prop:ordinarymapping}.
\end{proof}

\begin{lemma}[Finite unions of integer-multiple subgroups]\label{ihs:dz:lem:group}
For any abelian group $G$ and positive integers $e_1,\ldots,e_s$ with
$s\ge1$, one has $G=\bigcup_je_jG$ if and only if $e_jG=G$ for some $j$. If
$G$ is ordered and the union is proper, it omits a positive element.
\end{lemma}
\begin{proof}
Assume every $e_jG$ is proper. Choose a prime divisor $p_j$ of $e_j$ with
$p_jG\ne G$; one exists, since surjectivity of multiplication by all prime
factors would give surjectivity of multiplication by $e_j$. Let $\mathcal P$
be the finite set of chosen primes and choose $x_\ell\notin \ell G$ for each
$\ell\in\mathcal P$. By the integer Chinese remainder theorem choose
$c_\ell\equiv1\pmod \ell$ and $c_\ell\equiv0\pmod q$ for the other
$q\in\mathcal P$. Then $x=\sum_{\ell\in\mathcal P}c_\ell x_\ell$ has the same
class as $x_\ell$ modulo $\ell G$, so lies in none of the $\ell G$ and hence
in none of the $e_jG$. The converse is immediate. In an ordered group
$x\ne0$, and one of $x,-x$ is positive and still outside every subgroup in
the union.
\end{proof}

\begin{theorem}[Sharp spectral-mapping criterion]\label{ihs:dz:thm:criterion}
For a nonconstant $p\in\C[X]$, the equality
\begin{equation}\label{ihs:dz:eq:criterion}
 \SigA(p(T))=p(\SigA(T))
\end{equation}
holds if and only if, for every $b\in p(\sigD(T))$, at least one $a\in Z_b$
has $e_a\Gamma=\Gamma$. If the condition fails at $b$, there is a positive
$\delta\in\Gamma$ with
\begin{equation}\label{ihs:dz:eq:witness}
 b+t^\delta\in\SigA(p(T))\setminus p(\SigA(T)).
\end{equation}
\end{theorem}
\begin{proof}
\Cref{ihs:dz:thm:image} says equality holds exactly when
$\bigcup_{a\in Z_b}e_a\Gamma$ contains every positive element for every
relevant $b$; as a union of subgroups it then contains all of $\Gamma$.
\Cref{ihs:dz:lem:group} gives exactly the stated condition, and if the union
is proper its positive omitted element supplies \eqref{ihs:dz:eq:witness}.
\end{proof}

For divisible $\Gamma$ the criterion always holds. More generally, an
unramified Drazin-spectral preimage, $e_a=1$, suffices for its fibre whatever
the group. For $\Gamma=\Z$ every relevant fibre must contain such an
unramified preimage; for $\Gamma=\Z[1/2]$, ramification of order a power of
two creates no obstruction, but an exclusively order-three fibre does. If $T$
is algebraic, its Drazin spectrum is empty and the criterion is vacuous, so
finite matrices alone cannot exhibit the nontrivial halo obstruction. The
obstruction couples a non-Drazin singularity with insufficient scalar roots
in the chosen workspace. It is not a failure of ordinary complex polynomial
spectral mapping, and it is not a claimed failure over the full surcomplex
field (\cref{ihs:dz:cor:fullmapping}).

\begin{proposition}[Pole-order profile]\label{ihs:dz:prop:polemap}
Let $\Alg$ be a Banach algebra and $b\in p(\sigma_\Alg(T))\setminus
p(\sigD(T))$. For each spectral $a$ with $p(a)=b$ let $\nu_a$ be its finite
pole order and $e_a=\operatorname{ord}_{X=a}(p(X)-b)$. The Drazin index of
$p(T)-b1$ is
\begin{equation}\label{ihs:dz:eq:polemap}
 \nu_b=\max_{a\in\sigma_\Alg(T),\ p(a)=b}\left\lceil\frac{\nu_a}{e_a}\right\rceil,
\end{equation}
and its inverse at a displacement $0\ne\eps\in\mK$ has exact valuation
$-\nu_bv(\eps)$.
\end{proposition}
\begin{proof}
The indicated spectral preimages are finite poles. Their finite Riesz
decomposition separates them from a part on which $p(T)-b1$ is invertible.
More precisely, let $P_a$ be the corresponding Riesz idempotents and work
in the corner $P_a\Alg P_a$, whose identity is $P_a$. There
$TP_a=aP_a+N_a$ with $N_a=(T-a1)P_a$ of nilpotency index $\nu_a$.
Thus $(p(T)-b1)P_a=N_a^{e_a}w_a(N_a)$, where $w_a(0)\ne0$ and the
polynomial $w_a$ is evaluated with corner identity $P_a$; the factor
$w_a(N_a)$ is invertible and commutes with $N_a$, so the $k$th power vanishes
exactly when $N_a^{e_ak}$ does, giving index $\lceil\nu_a/e_a\rceil$. The
index of a finite direct sum is the maximum. The complementary corner, if
nonzero, has index zero; the nonempty spectral fibre ensures $\nu_b\ge1$.
Apply \cref{ihs:dz:thm:scale}
for the valuation.
\end{proof}

The ceiling formula is an elementary nilpotent-algebra fact, not claimed as
new by itself; here it gives the exact resolvent exponent.

\section{Examples separating poles, halos and scalar-root defects}\label{ihs:dz:sec:examples}

\subsection{Finite Jordan blocks and infinite-rank nilpotents}\label{ihs:dz:sec:jordan}
Let $J_\nu$ be the nilpotent $\nu\times\nu$ Jordan matrix with ones on the
superdiagonal. For $T=aI+J_\nu$,
\[
 \SigA(T)=\{a\},\qquad
 (T-aI-\eps I)^{-1}=-\sum_{j=0}^{\nu-1}\eps^{-j-1}J_\nu^j.
\]
Every nonzero infinitesimal displacement is invertible, although the inverse
has valuation $-\nu v(\eps)$: a large inverse and a spectral value are
different notions. For example $J_5^2$ has nilpotency index three, and its
inverse at displacement $\eps$ has valuation $-3v(\eps)$, not $-5v(\eps)$ or
$-\tfrac52v(\eps)$.

Infinite rank alone does not force a halo. On $H\oplus H$ let
$N=\begin{psmallmatrix}0&X\\0&0\end{psmallmatrix}$ with $X$ any bounded
operator, possibly compact of infinite rank. Then $N^2=0$, so
$\Sigma^{\B}_\Gamma(N)=\{0\}$ and its shifted inverse has only finitely many
negative powers.

\subsection{Volterra: an isolated point that acquires a whole halo}\label{ihs:dz:sec:volterra}
Let $H=L^2([0,1],\C)$ and $(\mathcal Vf)(x)=\int_0^xf(s)\,ds$. The operator is
injective; its range consists exactly of the $L^2$ classes with an absolutely
continuous representative $g$ satisfying $g(0)=0$ and $g'\in L^2([0,1])$.
Indeed, an indefinite integral has those properties, and conversely
$g=\mathcal Vg'$ almost everywhere. The constant function $1$ has no such
representative and is not in the range. Repeated integration gives
\begin{equation}\label{ihs:dz:eq:volterra}
 (\mathcal V^nf)(x)=\frac1{(n-1)!}\int_0^x(x-s)^{n-1}f(s)\,ds\qquad(n\ge1),
\end{equation}
and the squared Hilbert--Schmidt norm of this kernel is
$1/\bigl(((n-1)!)^2(2n-1)(2n)\bigr)\le1/(n!)^2$. Thus
$\norm{\mathcal V^n}\le1/n!$, and for every ordinary $z\ne0$ the series
$-\sum_{n\ge0}z^{-n-1}\mathcal V^n$ converges in operator norm and inverts
$\mathcal V-zI$. Since $\mathcal V$ is not onto, its ordinary spectrum is
exactly $\{0\}$.

Nevertheless $\mathcal V^n1=x^n/n!\ne0$ for every $n$, so $\mathcal V$ is not
nilpotent. It is not Drazin invertible either: if $S^{\mathrm D}$ were a
Drazin inverse, then $P_{\mathrm{inv}}=\mathcal V^n(S^{\mathrm D})^n$ for
every $n\ge1$, of norm at most $\norm{S^{\mathrm D}}^n/n!\to0$; hence
$P_{\mathrm{inv}}=0$ and $\mathcal V$ would be nilpotent. \Cref{ihs:dz:thm:halo}
gives
\begin{equation}\label{ihs:dz:eq:volterrahalo}
 \Sigma^{\B}_\Gamma(\mathcal V)=\mK.
\end{equation}
The ordinary spectrum has no accumulation points at all, so an
accumulation-only formula is false without normality; this is the
nonnormal counterpart of \cref{ihs:hh:ex:backshift}, which exhibits the
failure of point-spectrum rigidity rather than of the spectral formula.

In this example the spectral obstruction is also genuine vector
nonsurjectivity on $H((t^\Gamma))$. If $x\ne0$ has leading exponent $\delta$,
the coefficient of $(\mathcal V-\eps I)x$ at $\delta$ is $\mathcal Vx_\delta\ne0$,
because the infinitesimal term has strictly larger order. The operator is
therefore injective and preserves the leading exponent; if its image
contained the constant vector $1$, then $\delta=0$ and $\mathcal Vx_0=1$,
which is impossible. This is \cref{ihs:hh:thm:defect} with $S=\mathcal V$ and
$E=-\eps I$, and it does not rely on any equivalence of categories. Thus
$\sigma^{\mathrm{alg}}_\Gamma(\mathcal V)=\Sigma^{\B}_\Gamma(\mathcal V)=\mK$:
away from $\mK$ the algebra inverses are bijections.

\subsection{The missing odd valuation levels}\label{ihs:dz:sec:odd}
Choose $\Gamma=\Z$ and $p(X)=X^e$ with $e\ge2$. For Volterra, $Z=\{0\}$ and
the relevant ramification index is $e$. \Cref{ihs:dz:thm:image} gives
\begin{equation}\label{ihs:dz:eq:volterrapower}
 \Sigma^{\B}_\Gamma(\mathcal V^e)=\mK,\qquad
 \{\lambda^e:\lambda\in\Sigma^{\B}_\Gamma(\mathcal V)\}
 =\{0\}\cup\{\zeta\in\mK\setminus\{0\}:v(\zeta)\in e\Z\}.
\end{equation}
For $e=2$ every odd positive valuation is missing from the image:
$t\in\Sigma^{\B}_\Gamma(\mathcal V^2)$ but $t$ is not the square of a point
of $\Sigma^{\B}_\Gamma(\mathcal V)$. After passage to $\Gamma=\Q$, $t^{1/2}$
exists and belongs to the spectrum of $\mathcal V$, so this image defect
disappears; yet $\mathcal V^2-tI$ is still not invertible there, by
\cref{ihs:dz:cor:workspace}. What changed was the existence of a spectral
preimage, not the invertibility of the target pencil.

\subsection{Norm-convergent finite sections can miss an entire halo}\label{ihs:dz:sec:blocks}
On $H=\bigoplus_{n\ge1}\C^n$ consider
\begin{equation}\label{ihs:dz:eq:blocks}
 \mathcal N=\bigoplus_{n\ge1}2^{-n}J_n.
\end{equation}
Let $\mathcal N_M$ retain the first $M$ blocks and vanish elsewhere. It has
finite rank and is nilpotent, with $\norm{\mathcal N-\mathcal N_M}\le
2^{-(M+1)}\to0$; hence $\mathcal N$ is compact. It is not nilpotent, because
arbitrarily long blocks occur. For $z\in\C^*$, finitely many initial blocks
have finite-dimensional inverses, and on the remaining blocks
$\norm{2^{-n}J_n}<\abs z/2$, so ordinary Neumann series bound their inverse
norms by $2/\abs z$; the inverses form a bounded direct sum. Thus
$\sigma_\C(\mathcal N)=\{0\}$, as $\mathcal N$ has nonzero kernel. It is not
Drazin invertible: for $k\ge1$ only the blocks with $n>k$ contribute to
$\mathcal N^k$, so $\norm{\mathcal N^k}=2^{-k(k+1)}$, and Drazin data would
give $\norm{P_{\mathrm{inv}}}\le2^{-k(k+1)}\norm{S^{\mathrm D}}^k\to0$, making
$\mathcal N$ nilpotent. Therefore
\begin{equation}\label{ihs:dz:eq:finitefailure}
 \Sigma^{\B}_\Gamma(\mathcal N_M)=\{0\}\ \text{for every }M,\qquad
 \Sigma^{\B}_\Gamma(\mathcal N)=\mK.
\end{equation}
There is a direct support witness for every $0\ne\eps\in\mK$. Put
$\eta=v(\eps)>0$ and $c=\lc(\eps)$. The constant vector $y$ whose $n$th
block is $2^{-n}e_n$ lies in $H$, since $\sum_n4^{-n}<\infty$. A putative
solution $x$ of $(\mathcal N-\eps I)x=y$ must, in the $n$th block, equal
\[
 x^{(n)}=-\sum_{j=0}^{n-1}\eps^{-j-1}(2^{-n}J_n)^j2^{-n}e_n.
\]
Its final term has valuation $-n\eta$ and leading vector
$-c^{-n}2^{-n^2}e_1\ne0$; every earlier term has strictly larger valuation.
The coefficient of the global vector $x$ at $-n\eta$ is therefore nonzero,
since its projection onto the $n$th block is nonzero. Its support would contain
$-\eta,-2\eta,-3\eta,\ldots$, which is not well ordered, so no Hahn-vector
solution exists. A finite number of block checks can never establish the
uniform bound on nilpotent chain lengths that would remove this obstruction.
At $\eps=0$ the operator already has nonzero kernel,
whereas outside $\mK$ its algebra inverse gives a bijection. Consequently
\[
 \sigma^{\mathrm{alg}}_\Gamma(\mathcal N)
 =\Sigma^{\B}_\Gamma(\mathcal N)=\mK.
\]
This extension to nonmonomial displacements follows from the same support
argument; the source's explicit block calculation used a monomial.
This is the Drazin counterpart of \cref{ihs:sec:finitesections}, where every
finite section of a normal operator is invertible but the infinite one is
not.

\subsection{Two ranks and a nonmonomial displacement}\label{ihs:dz:sec:tworank}
Let $\Gamma=\Z\times\Z$ in lexicographic order and $\eta=(0,1)>0$. The
sequence $n\eta$ never exceeds $(1,0)$, yet all the preceding inverse
formulas remain strongly valid. For example take
$\eps=t^{(0,1)}\bigl(1+t^{(0,2)}+t^{(1,0)}\bigr)$. At a Drazin point of index
three, its shifted inverse has exact valuation $(0,-3)$, irrespective of the
nonmonomial correction. The straightening lemma applies within this very
group and adjoins no fractions. In contrast, the image of the squaring map
on infinitesimals consists of zero and the nonzero infinitesimals whose
valuation has both coordinates even; inverse descent and the scalar-power criterion answer
different questions. Under the ordered embedding $(m,n)\mapsto m\omega+n$
and $t^\gamma=\omega^{-\gamma}$, the displacement is the actual positive
surreal number $\eps=\omega^{-1}(1+\omega^{-2}+\omega^{-\omega})$, and the
exact inverse valuation at index three corresponds to the leading monomial
$\omega^3$, with operator coefficient given by \cref{ihs:dz:thm:scale}.

\section{Actual surreal and surcomplex specialization}\label{ihs:dz:sec:surreal}

The normal-form bridge is the imported classical theorem of
\cref{ihs:sec:surcomplex}, with the sign convention $t^\gamma=\omega^{-\gamma}$
of \eqref{ihs:eq:surreal-embedding} \cite{Gonshor,DriesEhrlich}; it is not a
newly proposed construction of the surreal field. Positive valuation is
exactly infinitesimal size in the embedded fields. Any \emph{set} of
surcomplex scalars fits in a common workspace: take the union of the supports
of their real and imaginary normal forms, negate the exponents, and form the
additive subgroup they generate, adjoining $1$ to make it nonzero. This is
set-sized, and its rational span is a set-sized divisible workspace when
scalar roots are wanted. No single set-sized workspace is asserted to contain
all surreal numbers.

For a fixed set-sized complex algebra $\Alg$, let $\Alg_{\No}$ denote the
class of expressions $\sum_{\gamma\in S}a_\gamma\omega^{-\gamma}$ with
$a_\gamma\in\Alg$ and $S\subset\No$ a well-ordered set. These are
\emph{operator-valued normal forms}, not scalar surreal numbers when $\Alg$ is
noncommutative. Any finite collection of them lies in a common $\AlgG$, so
additions and products are defined there and independent of enlargement.
This is a locality construction, not a theorem with a proper class passed as
its set-sized group parameter. For constant $T\in\Alg$ define
$\Sigma_{\No,\Alg}(T)$ by failure of a two-sided inverse of $T-\lambda1$ in
$\Alg_{\No}$, where $\lambda\in\No[i]$, and write $\mathfrak m_{\No[i]}$ for
the surcomplex infinitesimals, including zero.

\begin{theorem}[Surcomplex halo classification]\label{ihs:dz:thm:surcomplex}
For every constant $T\in\Alg$,
\begin{equation}\label{ihs:dz:eq:surcomplex}
 \Sigma_{\No,\Alg}(T)=\sigma_\Alg(T)\cup
 \bigl(\sigD(T)+\mathfrak m_{\No[i]}\bigr).
\end{equation}
The exact resolvent formulas and leading-scale assertions hold at every
actual surcomplex parameter in their indicated cases.
\end{theorem}
\begin{proof}
Place the normal-form support of $\lambda$ in a set-sized nonzero ordered
subgroup $\Gamma$. If \cref{ihs:dz:thm:halo} gives an inverse, its formula
already lies in $\AlgG\subseteq\Alg_{\No}$. Conversely a putative inverse in
$\Alg_{\No}$ has set-sized support; enlarge $\Gamma$ to a set-sized subgroup
containing both supports. The set-sized theorem then forces the same
ordinary or Drazin invertibility condition, which is independent of the
enlargement by \cref{ihs:dz:cor:workspace}. The exact leading terms follow
from \cref{ihs:dz:thm:scale}.
\end{proof}

\begin{corollary}[Polynomial mapping for the full surcomplex scalars]\label{ihs:dz:cor:fullmapping}
For nonconstant $p\in\C[X]$,
$\Sigma_{\No,\Alg}(p(T))=p(\Sigma_{\No,\Alg}(T))$.
\end{corollary}
\begin{proof}
For any particular parameter and its support choose a set-sized divisible
workspace as above. The criterion of \cref{ihs:dz:thm:criterion} holds there,
and inverses and spectral membership are unchanged after further
enlargement. Alternatively, every positive surreal valuation is divisible by
each positive integer, and the local germ construction supplies all the
required spectral preimages in the full field.
\end{proof}

The nondivisible counterexamples are therefore genuine phenomena of exact
Hahn subfields of the surcomplexes, while the non-Drazin halo itself persists
over the full surcomplex scalar field. Passing from the integer workspace to
the full field fills missing polynomial preimages; it does not remove the
target spectral obstruction.

For a real unital algebra $\Alg_\R$, complexify to $\Alg=\Alg_\R\otimes_\R\C$.
At real Hahn parameters, invertibility of a real pencil in the
complexification is equivalent to invertibility in the real coefficient
algebra, because the unique complex inverse is fixed by conjugation and so
has real coefficients; likewise uniqueness of the Drazin inverse shows that
Drazin data of a real element, when they exist in the complexification, are
real. The real-parameter version of \cref{ihs:dz:thm:surcomplex} follows by
intersection with $\No$. This avoids choosing rational powers of negative
real leading coefficients in the straightening proof: real coefficients
descend by conjugation only after the complex theorem is established. For
the real Volterra operator, every nonzero real surreal infinitesimal remains
a spectral value in the corresponding real operator-normal-form algebra.

\section{Part III: interpretation, boundaries and further questions}\label{ihs:dz:sec:limits}

\subsection{What the results establish}\label{ihs:dz:sec:establish}
There are two independent mechanisms. A constant pencil admits a Hahn
inverse at a nonzero infinitesimal displacement exactly when it has a Drazin
decomposition, with an invertible corner and a corner nilpotent of finite
index, and that condition survives all
exponent-workspace enlargements. A polynomial can then map the resulting
halos onto only selected valuation levels when the workspace lacks enough
roots, and the defect is determined exactly by the ramification indices over
the Drazin-spectral fibres. The argument is not a transfer of the analytic
theory of complex operators into a non-Archimedean setting: it uses finite
algebra on the singular part and exact support calculus on the regular part.
The surreal extension localizes all supports in set-sized subfields rather
than declaring a proper class to be a Banach field. The Volterra example
explains why an ordinary inverse series that converges for every nonzero
complex parameter can still be inadmissible at every nonzero infinitesimal:
its negative exponent tail descends without a least element. The
increasing-block example explains why all finite nilpotent approximants may
pass while their norm limit fails. These are structural obstructions, not
floating-point problems.

\subsection{Non-claims that matter mathematically}\label{ihs:dz:sec:nonclaims}
The spectra of this part are relative to the named algebras $\AlgG$ and
$\Alg_{\No}$, and concern \emph{constant} elements only. No classification
of all linear endomorphisms, of all possible non-Archimedean Hilbert
completions, or of all Hahn operator series is asserted. Actual vector
nonsurjectivity was proved for the two explicit examples, not inferred
universally from lack of an inverse in a specified subalgebra. The Banach
algebraicity corollary depends on compactness of the ordinary spectrum and
the finite Riesz decomposition and is false for arbitrary complex algebras.
The pole-order formula concerns a finite integer index; generalized Drazin
invertibility permits an infinite-order obstruction and cannot replace that
hypothesis. Every infinite operation uses strong support conditions, and no
conclusion about convergence of finite subsums in a higher-rank valuation
topology is implicit. No least field-valued operator norm, trace-class
determinant, spectral measure or physical application is supplied by this
part. The main theorems are proved in the text; they have not been
independently refereed or formalized in Lean, and the finite program of
\cref{ihs:sec:checksthree} cannot validate arbitrary support well ordering,
proper-class foundations, or the infinite-dimensional spectral formula.

\subsection{Natural next questions}\label{ihs:dz:sec:questions}
The first genuine extension is from a constant pencil to $T_0+E-\lambda1$,
with $E$ of positive Hahn order and not necessarily commuting with $T_0$.
The straightening argument no longer removes $E$, and a singular Drazin
corner can interact with its complement. A useful goal is an exact
Schur-reduction criterion with a support certificate, particularly when the
singular corner is infinite dimensional; no such general classification is
claimed. \emph{Status in this report:} Part~IV answers one case of this
question, for $\Alg=\B$, $T_0$ injective compact self-adjoint and $E$
self-adjoint in $\Ahh$, over a divisible $\Gamma$
(\cref{ihs:fr:thm:main}); there the halo of the Drazin-spectral point $0$
stays the whole of $\mK$, and each isolated simple pole $c\ne0$ is replaced
by finitely many Hermitian cluster eigenvalues in $c+\mK$. The general
question remains open.

A second question concerns the operator category. For which nondivisible
value groups and vector modules does bijectivity of an everywhere-defined
coefficientwise bounded Hahn operator force its inverse to have a common
well-ordered operator support? An answer would relate the present
algebra-relative spectrum to larger endomorphism spectra, such as
$\sigma^{\mathrm{alg}}_\Gamma$ of Part~I; the examples show why assuming that
these categories coincide is unsafe. A third direction is an analytic
functional calculus beyond polynomials: ordinary holomorphic germs with
nonzero finite ramification can be handled locally by the same formal
inverse construction, but a global Hahn-valued calculus must specify common
domains and admissible coefficient supports. For formal verification, the
clean first target is the finite-Laurent Drazin criterion, which needs only a
unital algebra and coefficient recurrences; the next are cyclic-support
projection and strong straightening. No checked implementation is included.

\clearpage
\part*{Part IV. Fredholm windows: compact self-adjoint residues under noncommuting Hahn perturbations}
\addcontentsline{toc}{part}{Part IV. Fredholm windows}

\emph{Standing hypotheses for the spectral classification in Part IV.} $H$ is an ordinary
\emph{infinite-dimensional separable} complex Hilbert space, $\Gamma$ is a
nonzero set-sized ordered abelian group that is \textbf{divisible}, and the
vector space, inner product and operator algebra are those of Part~I:
$\Hil=H((t^\Gamma))$ with the inner product \eqref{ihs:hh:eq:inner}, linear in
the first variable, and $\Ahh=\B((t^\Gamma))$. Divisibility is a hypothesis
of this classification; it is used to make $\K$ algebraically closed and $\Fld$
real closed, so that finite Hermitian clusters diagonalize inside the field
(\cref{ihs:fr:lem:hermitian}). The trace and determinant results in
\cref{ihs:fr:sec:determinant} separately allow arbitrary $H$ and $\Gamma$;
the spectral statements keep the hypotheses above. ``Spectrum'' in this part means both Part~I
spectra: $\sigma^{\mathrm{alg}}_\Gamma$, failure of bijectivity on $\Hil$,
and $\sigma^{\mathrm{adj}}_\Gamma$, failure of invertibility in $\Ahh$; the
main theorem proves that they agree for the operators studied. The source
manuscript wrote $\mathscr B$, $\sigma_{\mathscr B}$ and $\sigma_{\mathrm{bij}}$
for $\Ahh$, $\sigma^{\mathrm{adj}}_\Gamma$ and $\sigma^{\mathrm{alg}}_\Gamma$,
$\operatorname{res}$ for the standard part $\st$, and $W$, $M$, $U$, $V$, $D$, $B$ for
objects that are renamed below to avoid Part~I and Part~II meanings.

\section{The question and the answer}\label{ihs:fr:sec:intro}

A useful surcomplex spectral theory must distinguish three operations:
solving each finite spectral cluster, assembling all clusters into one
operator, and detecting the spectrum by a scalar analytic determinant. In
the coefficientwise Hahn setting these operations have different answers.
This part gives a class in which all three answers are proved exactly. The
principal example starts with a compact self-adjoint operator whose ordinary
eigenvalues accumulate only at zero. An infinitesimal noncommuting
perturbation moves the individual eigenvalues. Nevertheless zero does not
simply acquire a sequence of displaced eigenvalues: it retains an entire
complex infinitesimal neighbourhood of non-eigenvalue spectral points. Every
nonzero ordinary spectral window is finite-dimensional, but the accumulation
window is not Fredholm in the algebraic sense of \cref{ihs:fr:sec:essential}.

Let $T$ be an ordinary injective compact self-adjoint operator on $H$ and
\begin{equation}\label{ihs:fr:eq:perturbation}
 A=T+E,\qquad E=E^*\in\Ahh,\qquad v(E)>0.
\end{equation}
The coefficient operators of $E$ need only be bounded; they need not be
compact or commute with each other or with $T$. For an ordinary eigenvalue
$c\ne0$ of $T$ write $m_c=\dim_\C\ker(T-cI)$. \Cref{ihs:fr:thm:main}
constructs $m_c$ real Hahn eigenvalues $\lambda_{c,1},\ldots,\lambda_{c,m_c}$
with standard part $c$, counted with multiplicity, and proves
\begin{equation}\label{ihs:fr:eq:headline}
 \sigma^{\mathrm{adj}}_\Gamma(A)=\sigma^{\mathrm{alg}}_\Gamma(A)
 =\mK\cup\bigcup_{c\in\sigma_\C(T)\setminus\{0\}}
 \{\lambda_{c,1},\ldots,\lambda_{c,m_c}\}.
\end{equation}
The second set is exactly the point spectrum. For $z\in\mK$, $A-zI$ is
injective but not surjective, with cokernel isomorphic to the Hahn extension
of an algebraic complement of $\Ran T$. Thus even nonreal infinitesimals
belong to the spectrum of this self-adjoint operator. This is not a claim
about the spectrum of an ordinary complex Hilbert-space operator.

\emph{Relation to Part~I.} For $E=0$, \eqref{ihs:fr:eq:headline} is the
injective compact self-adjoint case of \cref{ihs:hh:cor:compact}, and
Part~I is more general there (normal, not necessarily injective, any
$\Gamma$, any $H$). The new content is the noncommuting perturbation. This is
the case that \cref{ihs:hh:sec:different} describes as requiring a different
theorem. The local answer, however, does not assemble into a global
bounded-coefficient spectral theorem: \cref{ihs:fr:thm:nounitary} gives an
explicit tridiagonal perturbation with trace-class coefficients and
individually defined eigenbranches but no Hahn-unitary diagonalizer in
$\Ahh$, and \cref{ihs:fr:thm:nolidskii,ihs:fr:thm:noproduct} show that even
ordinary coefficientwise eigenvalue sums and products fail for it. A
coefficientwise lift of the trace-class Fredholm determinant nevertheless
exists and has an exact Fredholm alternative
(\cref{ihs:fr:thm:fredholm}); it detects every finite-coupling spectral
window, but no nonzero coherent Hahn analytic function can detect the whole
spectral monad at the accumulation point
(\cref{ihs:fr:thm:nocharacteristic}).

\emph{Relation to Part~III.} In the language of Part~III,
$\sigD(T)=\{0\}$ for $\Alg=\B$: zero is an accumulation point, and each
nonzero eigenvalue is an isolated simple pole. \Cref{ihs:dz:thm:halo} then
gives $\Sigma^{\B}_\Gamma(T)=\mK\cup(\sigma_\C(T)\setminus\{0\})$ at $E=0$.
Part~IV answers the first question of \cref{ihs:dz:sec:questions} for this
class: under a noncommuting $E$ the halo of the Drazin-spectral point stays
all of $\mK$, while each pole $c$ is replaced by a finite Hermitian cluster.

\emph{What is inherited and what is added.} The range-defect factorization
used below is \cref{ihs:hh:thm:defect} of Part~I, re-proved in the source by
a different factorization (\cref{ihs:fr:rem:secondproof}); the Hahn support
lemma, ordinary compact spectral theory, the Riesz projection construction,
finite Hermitian diagonalization, and the classical trace-class determinant
identities are likewise inherited or classical, and their transfer alone is
not presented as new. Ordinary analytic perturbation theory supplies isolated
spectral projections and finite-dimensional reduction \cite{Kato}; the
direct rotation of two nearby orthogonal projections is the classical
Sz.-Nagy--Kato construction, with the explicit formula in Simon's account
\cite{SimonProjections}; ordinary Fredholm determinant identities are
standard trace-ideal theory, used in Bornemann's formulation
\cite{Bornemann}; non-Archimedean Fredholm theory has a substantial
precedent in Serre's work \cite{Serre}, whose category of completely
continuous operators on a non-Archimedean Banach space is not the
coefficient algebra used here. The proposed additions are the full formula
\eqref{ihs:fr:eq:headline} for arbitrary noncommuting self-adjoint
positive-order perturbations, the trace-class counterexample to global
bounded-coefficient diagonalization and to spectral sums and products, and
the exact interface between Hahn Fredholm determinants and the unavoidable
accumulation monad. A targeted search did not locate these combined
statements; it is not an exhaustive priority determination.

\section{The operator category and the identity-lifting device}\label{ihs:fr:sec:framework}

\subsection{The algebra is specified in advance}\label{ihs:fr:sec:category}
All objects are those of \cref{ihs:sec:construction,ihs:sec:adjoints}. Part~I
proves that $\Ahh$ is exactly the algebra of everywhere-defined adjointable
operators (\cref{ihs:hh:thm:adjoint}); this part does not need that theorem,
because it specifies $\Ahh$ at the outset. Membership in $\Ahh$ requires every
coefficient to be an ordinary bounded operator and one well-ordered operator
support; it does not require $\sup_\gamma\norm{A_\gamma}<\infty$. Conversely,
a matrix whose entries are individually valid Hahn series need not define an
element of $\Ahh$: its coefficient at a single exponent can be an unbounded
ordinary matrix. This is the obstruction used in
\cref{ihs:fr:sec:obstruction}. For $X\in\Ahh$ also write
$\sigma_{p,\Gamma}(X)=\{\lambda:\ker(X-\lambda I)\ne0\}$. The inclusion
$\sigma^{\mathrm{alg}}_\Gamma(X)\subseteq\sigma^{\mathrm{adj}}_\Gamma(X)$
always holds; equality is proved below for the operators studied, not
inferred for every $X\in\Ahh$.

Divisibility makes $\K$ algebraically closed and $\Fld$ real closed. The
complex statement is the standard Hahn-field closedness theorem
\cite[Corollary~4]{Poonen}; the real statement follows equivalently from
positive square roots and complexification, and a positive square root can
be built directly from half the leading exponent and a near-one binomial
series (\cref{ihs:hh:sec:field}). The results apply to genuine surcomplex
coefficients through \eqref{ihs:eq:surreal-embedding}; they posit no set
containing the proper class $\No[i]$ and do not identify the intrinsic
valuation topology with the fine topology of the surreal class.

The support tools are \cref{ihs:lem:support,ihs:lem:neumann}; the correction
exponents of $(I+F)^{-1}$ lie in $(\supp F)^*\setminus\{0\}$. If
$X=X_0+X_+$ with $X_0$ an invertible ordinary bounded operator and
$v(X_+)>0$, factoring out $X_0$ gives an inverse in $\Ahh$ with nonnegative
support. This criterion is sufficient, not necessary: $t^\eta I$ is
invertible although its standard part is zero.

\subsection{Lifting analytic identities coefficientwise}\label{ihs:fr:sec:lift}
Let $X,Y$ be ordinary complex Banach spaces and let $f:U\to Y$ be
holomorphic on an open neighbourhood $U\subseteq X$ of $x_0$. For a positive-support Hahn perturbation
$h=\sum_{s\in S}h_st^s$ define
\begin{equation}\label{ihs:fr:eq:taylorlift}
 f^{\mathrm H}(x_0+h)=\sum_{n\ge0}\frac1{n!}\,\mathrm d^nf(x_0)[h,\ldots,h],
\end{equation}
expanded into the finite coefficient words of \cref{ihs:lem:support}(ii). It
is a formal Hahn sum; no inequality involving ordinary norms of the $h_s$ is
required.

\begin{lemma}[Identity lifting]\label{ihs:fr:lem:lift}
Ordinary holomorphic identities between maps as above lift to identities of
their Hahn Taylor expansions whenever the ordinary base points lie in the
domains of the identities. Finite compositions, products and continuous
multilinear operations commute with this lifting.
\end{lemma}
\begin{proof}
Fix an exponent $\gamma$. Only finitely many words in $S$ contribute to that
coefficient, and only finitely many directions $h_s$ occur in them.
Introduce one ordinary complex variable for each such direction and restrict
the ordinary identity to the resulting finite-dimensional parameter space
near zero. Equality of its Taylor coefficients gives the coefficient
equality at $\gamma$. For composition, zero-constant inner increments and the
same finite-word condition justify all regrouping.
For several inputs use the finite union of their positive supports; each
output coefficient still uses finitely many directions and word lengths.
The argument bounds the number of contributors to one coefficient, not the
ordinary norms of all Hahn coefficients simultaneously.
\end{proof}

A real-analytic identity involving adjoints is treated by real auxiliary
parameters and compared coefficientwise in the same way; this is used for
orthogonal projections of self-adjoint families. An infinite ordinary sum at
one exponent must first be evaluated in its ordinary coefficient space:
infinitely many nonzero constants are not a strongly Hahn-summable family.
Ordinary traces, contour integrals and determinants are taken inside the
coefficient space before the Hahn support argument is applied.

\section{The inherited range defect and the infinitesimal reservoir}\label{ihs:fr:sec:defect}

The range-defect mechanism is \cref{ihs:hh:thm:defect}: for $S:V\to V$
injective, $F\in\End_\C(V)((t^\Gamma))$ of positive order and an algebraic
complement $V=\Ran S\oplus W$, the operator $S+F$ is injective and
$\coker(S+F)\cong W((t^\Gamma))$. Its elementary factorization allows
completely noncommuting perturbations, and it is printed once, in Part~I.

\begin{remark}[A second proof, from the Part~IV source]\label{ihs:fr:rem:secondproof}
Define $L:V\to V$ by $L(Sx+w)=x$, so that $LS=I$ and
\[
 S+F=(I+FL)S.
\]
Since $FL$ has positive support, \cref{ihs:lem:neumann} inverts $I+FL$ in
the Hahn algebra of all endomorphisms, so multiplication by $I+FL$ is an
automorphism of $V((t^\Gamma))$ carrying $S(V((t^\Gamma)))$ onto
$\Ran(S+F)$. Applying the two ordinary projections of $V=\Ran S\oplus W$ to
every coefficient does not enlarge supports, so
$V((t^\Gamma))=(\Ran S)((t^\Gamma))\oplus W((t^\Gamma))$, and the
coefficientwise inverse of $S$ on its range gives
$S(V((t^\Gamma)))=(\Ran S)((t^\Gamma))$. Both conclusions follow. This
factorizes on the left of $S$, whereas the Part~I proof builds a left
inverse and a complementary projection; both use only the Neumann lemma. For
compact injective $S$ with nonclosed range, $L$ is generally unbounded; it is
used only in the auxiliary algebra of all endomorphisms, and the
factorization produces no bounded inverse, no bounded projection onto $W$,
and no orthogonal decomposition of $\Hil$.
\end{remark}

\begin{proposition}[The size of a compact range defect]\label{ihs:fr:prop:continuum}
If $T$ is injective, compact and self-adjoint on an ordinary
infinite-dimensional separable $H$, then $\dim_\C(H/\Ran T)=2^{\aleph_0}$.
Consequently $W((t^\Gamma))$ contains at least $2^{\aleph_0}$ vectors that
are linearly independent over $\K$.
\end{proposition}
\begin{proof}
The compact spectral theorem gives an orthonormal eigenbasis with nonzero
real eigenvalues $d_n\to0$. Choose a subsequence with
$0<\abs{d_{n_k}}\le2^{-k^2}$ and, for real $a\in(0,1)$, put
$x^{(a)}_{n_k}=\abs{d_{n_k}}^a$ and $x^{(a)}_j=0$ off the subsequence; the
vector is square summable. For a nonzero combination
$x=\sum_{j=1}^rc_jx^{(a_j)}$ with $a_1<\cdots<a_r$ and $c_1\ne0$,
\[
 \frac{x_{n_k}}{d_{n_k}}=\frac{\abs{d_{n_k}}^{a_1}}{d_{n_k}}
 \Bigl(c_1+\sum_{j=2}^rc_j\abs{d_{n_k}}^{a_j-a_1}\Bigr);
\]
the bracket tends to $c_1$ and the first factor tends to infinity in
absolute value, so these coordinates are not square summable and
$x\notin\Ran T$. This gives continuum many independent cosets; the
underlying set of a separable Hilbert space has the cardinality of the
continuum, which gives the other bound. Independent cosets are represented by
independent vectors of $W$. A finite $\K$-linear dependence among their
constant embeddings would, at any exponent with a nonzero scalar coefficient,
give an ordinary complex dependence, so they remain independent over $\K$
(this is the argument of \cref{ihs:hh:cor:independent}).
\end{proof}

The exact $\K$-dimension of $W((t^\Gamma))$ is not asserted for every large
$\Gamma$; only the lower bound is used (compare \cref{ihs:hh:warn:dimension}).
For the particular operator $De_n=n^{-1}e_n$, \cref{ihs:hh:lem:continuum}
gives the same lower bound by explicit power vectors.

\begin{corollary}[The infinitesimal reservoir]\label{ihs:fr:cor:reservoir}
Under \eqref{ihs:fr:eq:perturbation}, every $z\in\mK$ satisfies
\[
 \ker(A-zI)=0,\qquad \coker(A-zI)\cong W((t^\Gamma)),\qquad
 H=\Ran T\oplus W.
\]
In particular $\mK\subseteq\sigma^{\mathrm{alg}}_\Gamma(A)$, and no point of
$\mK$ is an eigenvalue.
\end{corollary}
\begin{proof}
Apply \cref{ihs:hh:thm:defect} with $S=T$ and positive-order perturbation
$E-zI$ (zero allowed); \cref{ihs:fr:prop:continuum} makes the cokernel
nonzero.
\end{proof}

\section{Isolated spectral windows are exactly finite-dimensional}\label{ihs:fr:sec:clusters}

Fix a nonzero ordinary eigenvalue $c$ of $T$. Let $P^\circ_c$ be its ordinary
orthogonal spectral projection and $H_c=P^\circ_cH$. Compactness makes $c$
isolated and $m_c=\dim H_c<\infty$. Choose a positively oriented ordinary
circle $\mathcal C_c$ around $c$ containing no other point of $\sigma_\C(T)$
in its closed interior and none on its boundary; its radius may shrink as
$c\to0$, and no common radius is required.

\subsection{Coefficientwise Riesz projections and direct rotation}\label{ihs:fr:sec:riesz}
Put $R_0(\zeta)=(\zeta I-T)^{-1}$ for $\zeta\in\mathcal C_c$ and define
\begin{equation}\label{ihs:fr:eq:riesz}
 R(\zeta)=\sum_{n\ge0}(R_0(\zeta)E)^nR_0(\zeta),\qquad
 P_c=\frac1{2\pi i}\int_{\mathcal C_c}R(\zeta)\,d\zeta.
\end{equation}
The integral is an ordinary bounded-operator-valued contour integral of each
Hahn coefficient, not an integral over a surcomplex contour.

\begin{proposition}\label{ihs:fr:prop:riesz}
Let $S=\supp E\subseteq\Gamma_{>0}$. Formula \eqref{ihs:fr:eq:riesz} defines
$P_c\in\Ahh$ with
\[
 P_c^2=P_c=P_c^*,\qquad AP_c=P_cA,\qquad
 P_c=P^\circ_c+\text{positive-order terms},
\]
and $\supp(P_c-P^\circ_c)\subseteq S^*\setminus\{0\}$.
\end{proposition}
\begin{proof}
At a fixed exponent the resolvent expansion has finitely many ordered
coefficient words, so each coefficient is a finite sum of products of
ordinary resolvents and bounded operators; on the compact circle these are
norm-continuous, and their contour integrals exist. The constant coefficient
is $P^\circ_c$. For any finite collection of directions $E_s$, the ordinary
Riesz projection of $T+\sum_sz_sE_s$ on the fixed circle, for complex
parameters near zero, is holomorphic, idempotent and commutes with the
perturbed operator; its Taylor coefficients are the contour coefficients
above. To keep this contour in the resolvent set, put
$M_c=\max_{\zeta\in\mathcal C_c}\norm{R_0(\zeta)}<\infty$ and restrict the
finitely many parameters by $\sum_s|z_s|\norm{E_s}<1/M_c$. The resulting
uniform geometric expansion justifies integrating its Taylor coefficients.
No common parameter neighbourhood for all directions or all $c$ is required.
Now \cref{ihs:fr:lem:lift} proves idempotence and commutation. Each
$E_s$ is self-adjoint because $E=E^*$, so for real parameters near zero the
projection is orthogonal; comparing real Taylor coefficients gives
$P_c^*=P_c$. \Cref{ihs:lem:support} gives the support containment.
\end{proof}

\begin{lemma}[Support-controlled direct rotation]\label{ihs:fr:lem:rotation}
Let $P=P^2=P^*\in\Ahh$ and let $P^\circ$ be an ordinary orthogonal projection
with $v(P-P^\circ)>0$. Put
\begin{equation}\label{ihs:fr:eq:rotation}
 \Omega=PP^\circ+(I-P)(I-P^\circ),\qquad
 U=\Omega\bigl(I-(P-P^\circ)^2\bigr)^{-1/2},
\end{equation}
the inverse square root being its binomial expansion at $I$. Then
$U^*U=UU^*=I$, $UP^\circ=PU$ and $v(U-I)>0$. If
$\supp(P-P^\circ)\subseteq S^*\setminus\{0\}$, then
$\supp(U-I)\subseteq S^*\setminus\{0\}$.
\end{lemma}
\begin{proof}
Direct multiplication using $P^2=P$ and $(P^\circ)^2=P^\circ$ gives
$\Omega^*\Omega=\Omega\Omega^*=I-(P-P^\circ)^2$ and $P\Omega=\Omega P^\circ$.
The square $(P-P^\circ)^2$ commutes with both projections and hence with
$\Omega$. The binomial series for $(I-(P-P^\circ)^2)^{-1/2}$ exists by
\cref{ihs:cor:binomial}, is self-adjoint, commutes with $\Omega,P,P^\circ$,
and squares to the inverse. The unitary and intertwining identities follow,
and all nonconstant terms have support in the stated monoid. In particular,
$\Omega-I=(P-P^\circ)(2P^\circ-I)$ has positive support, and the binomial
factor has constant coefficient $I$, so $U-I$ has positive support as claimed.
\end{proof}

This is the classical Kato two-projection formula, as recalled by Simon
\cite[Equation~(4)]{SimonProjections},
recorded in the present algebra. Apply it to $P=P_c$ and $P^\circ=P^\circ_c$,
and call the unitary $U_c$. Then $U_c^*AU_c$ commutes with $P^\circ_c$, so
relative to $H=H_c\oplus H_c^\perp$ it has an exact block form
\begin{equation}\label{ihs:fr:eq:block}
 U_c^*AU_c=A_c\oplus A_c^\perp,\qquad A_c\in M_{m_c}(\K),\quad A_c=A_c^*.
\end{equation}
The standard part of $A_c$ is $cI_{m_c}$ and that of $A_c^\perp$ is
$T|_{H_c^\perp}$; both blocks have nonnegative support, and their positive
exponents lie in $S^*$.

\subsection{Finite Hermitian algebra over the Hahn field}\label{ihs:fr:sec:hermitian}
\begin{lemma}[Finite Hermitian diagonalization]\label{ihs:fr:lem:hermitian}
Every Hermitian matrix over $\K=\Fld[i]$ is unitarily diagonalizable, with
eigenvalues in $\Fld$. If its entries are finite and its standard part is
$cI_m$, all its eigenvalues are finite with standard part $c$.
\end{lemma}
\begin{proof}
The zero-dimensional case is empty. In positive dimension, since $\K$ is
algebraically closed, the characteristic polynomial has a root
and an eigenvector $x\ne0$. For a Hermitian $X$,
$\lambda\ip xx=\ip{Xx}x=\overline{\ip{Xx}x}$ with $\ip xx>0$ in $\Fld$, so
$\lambda$ is real. The positive square root of $\ip xx$ normalizes $x$; its
orthogonal complement is invariant, and induction gives an orthonormal
eigenbasis. For the standard-part assertion, write the monic characteristic
polynomial as $X^m+\sum_{j<m}a_jX^j$ with $v(a_j)\ge0$. A root of negative
valuation would make $X^m$ strictly smaller in valuation than every other
term, impossible in a zero sum; so every root is finite, and reducing the
characteristic equation gives $(\st\lambda-c)^m=0$.
\end{proof}

\begin{theorem}[Exact isolated-cluster reduction]\label{ihs:fr:thm:cluster}
For every nonzero ordinary eigenvalue $c$ of $T$, the construction above
produces a finite Hermitian matrix $A_c$ with exactly $m_c$ real Hahn
eigenvalues $\lambda_{c,1},\ldots,\lambda_{c,m_c}$ counted with
multiplicity, all with standard part $c$. For every $\lambda\in c+\mK$,
\begin{equation}\label{ihs:fr:eq:clusterkernel}
 \ker(A-\lambda I)\cong\ker(A_c-\lambda I),\qquad
 A-\lambda I\text{ invertible in }\Ahh\iff\det(\lambda I_{m_c}-A_c)\ne0.
\end{equation}
The kernel is finite-dimensional, of dimension the multiplicity of
$\lambda$ in the Hermitian eigenvalue list.
\end{theorem}
\begin{proof}
\Cref{ihs:fr:lem:hermitian} applies to $A_c$. On the complementary block the
standard part of $A_c^\perp-\lambda I$ is $T|_{H_c^\perp}-cI$, an ordinary
bounded invertible operator because $c$ has been removed from the compact
spectrum; the remaining terms have positive valuation, and
\cref{ihs:lem:neumann} inverts the block in its bounded-coefficient Hahn
algebra. If the finite determinant is nonzero, invert $A_c-\lambda I$ over
$\K$, combine with the complementary inverse and conjugate by $U_c$; a
finite matrix of scalar Hahn entries acts by a bounded coefficient at each
exponent on the finite-dimensional block, so the inverse lies in $\Ahh$,
even if some exponents are negative. If the determinant is zero, a nonzero
finite block eigenvector gives a nonzero vector of $\Hil$ after conjugation.
The complementary inverse also proves the kernel isomorphism, explicitly
$y\mapsto U_c(y,0)$. On quotients the map induced by
$x\mapsto P^\circ_cU_c^*x$ gives
$\coker(A-\lambda I)\cong\coker(A_c-\lambda I)$: the complementary
coordinate is in the image because its block is invertible. Both finite
defect dimensions therefore equal the multiplicity in the Hermitian
eigenvalue list, including zero when $\lambda$ is absent from it.
\end{proof}

\begin{remark}[Precisely what is support controlled]\label{ihs:fr:rem:support}
The entries of $A_c-cI$ and the projection and rotation corrections have
support in the original positive monoid $S^*$. For a \emph{multiple}
residue eigenvalue it is not asserted that each root of the finite
characteristic polynomial has support in that monoid; the roots belong to
the original divisible Hahn field, and \cref{ihs:fr:sec:monoid} shows that
their supports can leave $S^*$. For a \emph{simple} residue eigenvalue,
ordinary analytic eigenvalue perturbation and \cref{ihs:fr:lem:lift} give
the eigenbranch directly with support in $S^*$. This matches the Part~II
boundary \cref{ihs:rf:rem:monoid}: monoid support is a consequence of simple
residues, not an unconditional invariant.
\end{remark}

\section{The complete spectrum of a noncommuting perturbation}\label{ihs:fr:sec:main}

\begin{theorem}[Compact-residue spectral classification]\label{ihs:fr:thm:main}
Under the standing hypotheses of Part~IV and \eqref{ihs:fr:eq:perturbation},
with $T$ injective, compact and self-adjoint,
\begin{align}
 \sigma^{\mathrm{adj}}_\Gamma(A)=\sigma^{\mathrm{alg}}_\Gamma(A)
 &=\mK\cup\bigcup_{c\in\sigma_\C(T)\setminus\{0\}}
   \{\lambda_{c,1},\ldots,\lambda_{c,m_c}\},\label{ihs:fr:eq:mainspectrum}\\
 \sigma_{p,\Gamma}(A)
 &=\bigcup_{c\in\sigma_\C(T)\setminus\{0\}}
   \{\lambda_{c,1},\ldots,\lambda_{c,m_c}\}.\label{ihs:fr:eq:mainpoint}
\end{align}
All eigenvalues are real with finite-dimensional eigenspaces. At every
$z\in\mK$, $A-zI$ is injective with cokernel $W((t^\Gamma))$, where
$H=\Ran T\oplus W$ is any algebraic complement; this cokernel has
$\K$-dimension at least $2^{\aleph_0}$.
\end{theorem}
\begin{proof}
Every $\lambda\in\K$ falls in one of four disjoint cases. If
$v(\lambda)<0$, then $A-\lambda I=-\lambda(I-\lambda^{-1}A)$ with
$v(\lambda^{-1}A)>0$, so the second factor has a Neumann inverse in $\Ahh$.
If $\lambda$ is finite and $c=\st\lambda\notin\sigma_\C(T)$, then
$A-\lambda I=(T-cI)+(E-(\lambda-c)I)$ is an ordinary invertible operator plus
a positive-order term, again invertible in $\Ahh$. If $c\ne0$ belongs to
$\sigma_\C(T)$, it is an isolated eigenvalue of finite multiplicity, and
\cref{ihs:fr:thm:cluster} gives exactly the listed eigenvalues in its halo;
every other point of the halo has an inverse in $\Ahh$. Finally, if
$\st\lambda=0$, including $\lambda=0$, then $\lambda\in\mK$ and
\cref{ihs:fr:cor:reservoir} gives injectivity and non-surjectivity with the
stated cokernel. Every excluded point has a constructed algebra inverse and
every included point fails bijectivity, proving equality of the two spectra.
The point-spectrum and dimension statements follow from the finite Hermitian
blocks and \cref{ihs:fr:prop:continuum}.
\end{proof}

Solving formal characteristic equations near each nonzero $c$ would establish
only that the displayed eigenvalues belong to the point spectrum; it would
say nothing about infinite scalars, ordinary resolvent halos, or the monad at
zero. The four-case proof controls all of these. The hypotheses have
identifiable roles: compactness makes every nonzero ordinary spectral point
an isolated finite cluster; injectivity makes the whole zero monad a
non-eigenvalue reservoir through \cref{ihs:hh:thm:defect}; self-adjointness
makes the clusters Hermitian and their roots real; divisibility keeps those
roots inside the selected field; and infinite dimension ensures that an
injective compact operator is not onto. None of these roles is replaced by a
claim of completeness of a surreal Hilbert space.

\subsection{The algebraic Fredholm spectrum}\label{ihs:fr:sec:essential}
Call a $\K$-linear endomorphism of $\Hil$ \emph{algebraically Fredholm} when
its kernel and cokernel are finite-dimensional over $\K$, with algebraic
index the difference of these dimensions; no topological closed-range
condition is included. Put
$\sigma_{\mathrm{ess,alg}}(X)=\{\lambda\in\K:X-\lambda I\text{ is not
algebraically Fredholm}\}$.

\begin{corollary}\label{ihs:fr:cor:essential}
Under the hypotheses of \cref{ihs:fr:thm:main},
$\sigma_{\mathrm{ess,alg}}(A)=\mK$, and at every point outside $\mK$ the
algebraic Fredholm index of $A-\lambda I$ is zero.
\end{corollary}
\begin{proof}
At a point of $\mK$ the cokernel has dimension at least the continuum.
Outside that monad the operator is either invertible or conjugate to a
finite square matrix plus an invertible complementary block, and a finite
square matrix has kernel and cokernel of the same finite dimension.
\end{proof}

Arbitrary positive-order bounded self-adjoint coefficient perturbations
therefore preserve the full algebraic Fredholm spectral monad of this
compact residue. The terminology is explicitly algebraic; no identification
with a Banach-space or compactoid essential spectrum is made, and the prime
of Part~I is not an essential spectrum either (\cref{ihs:hh:sec:reality}).

\subsection{Extension of the exponent group}\label{ihs:fr:sec:extension}
\begin{corollary}[Fixed point spectrum, growing spectral monad]\label{ihs:fr:cor:extension}
Let $\iota:\Gamma\hookrightarrow\Gamma'$ be an ordered embedding into another
nonzero divisible set-sized ordered abelian group, and extend all Hahn series
coefficientwise. The point spectrum and the finite cluster multiplicities of
$A$ are unchanged under the scalar embedding, whereas
\begin{equation}\label{ihs:fr:eq:extension}
 \sigma^{\mathrm{adj}}_{\Gamma'}(A)=\mathfrak m_{K_{\Gamma'}}\cup
 \iota\bigl(\sigma_{p,\Gamma}(A)\bigr).
\end{equation}
Enlarging the workspace can add spectral points without adding eigenvalues.
\end{corollary}
\begin{proof}
The constructions of $P_c,U_c,A_c$ commute with an ordered exponent
embedding: at every exponent the same finite algebraic and ordinary contour
operations are performed. The characteristic polynomial of $A_c$ already
splits over $\K$, so it has no new roots in the extension and its
multiplicities are unchanged. Apply \cref{ihs:fr:thm:main} over $\Gamma'$;
its non-eigenvalue part is exactly the enlarged infinitesimal ideal.
\end{proof}

For example, extend $\Gamma=\Q$ into $\Q^2$ with lexicographic order by
$q\mapsto(0,q)$; then $it^{(1,0)}$ is a new spectral point but not a new
eigenvalue. The formula concerns spectra in their respective modules; it
does not assert that scalar extension leaves the module unchanged. It is
consistent with \cref{ihs:hh:cor:extension,ihs:dz:cor:workspace}, which say
that each old scalar keeps its spectral status.

\section{A trace-class counterexample to simultaneous diagonalization}\label{ihs:fr:sec:obstruction}

The local finite-dimensional theorem does not imply a global
bounded-coefficient diagonalization, even for a one-scale perturbation with
exceptionally small ordinary coefficient ideals.

\subsection{The operator and its individual eigenvalues}\label{ihs:fr:sec:operator}
Let $H=\ell^2(\N_{\ge1})$, fix $\eta>0$ in $\Gamma$, and put
\begin{equation}\label{ihs:fr:eq:counterexample}
 D_\star e_n=2^{-n}e_n,\qquad
 Y=\sum_{n\ge1}\frac1{n^2}\bigl(\lvert e_n\rangle\langle e_{n+1}\rvert
       +\lvert e_{n+1}\rangle\langle e_n\rvert\bigr),\qquad
 A_\star=D_\star+t^\eta Y.
\end{equation}
The sum defining $Y$ is an ordinary trace-norm convergent series of rank-one
operators, evaluated before it is used as a Hahn coefficient:
$\norm{D_\star}_1=1$ and $\norm Y_1\le2\sum_n n^{-2}<\infty$. Both
coefficients are self-adjoint and trace class, and they do not commute.
$A_\star$ is positive definite: if $v(x)=\alpha$, then $\ip{A_\star x}x$ has
leading term $\ip{D_\star x_\alpha}{x_\alpha}_Ht^{2\alpha}>0$. For each $n$,
\cref{ihs:fr:thm:cluster} supplies one eigenvalue $\lambda_n$ with standard
part $2^{-n}$; since that residue is simple, $\lambda_n$ is the unique
ordinary analytic perturbation series in $s=t^\eta$, so
$\lambda_n\in\C[[s]]\subseteq\K$. The existence of these formal series does
not give one common ordinary analytic radius as $n\to\infty$.

\begin{theorem}[No global bounded-coefficient Hahn unitary]\label{ihs:fr:thm:nounitary}
For $A_\star$ of \eqref{ihs:fr:eq:counterexample}, no $U\in\Ahh$ with
$U^*U=UU^*=I$ makes $U^*A_\star U$ diagonal relative to $(e_n)$. This holds
even when the support of $U$ may use arbitrary exponents of $\Gamma$, not
only integer multiples of $\eta$.
\end{theorem}
\begin{proof}
Suppose such $U$ exists and put $\Lambda=U^*A_\star U$. If $\alpha=v(U)$,
the coefficient of $U^*U$ at $2\alpha$ is $U_\alpha^*U_\alpha\ne0$, so
$2\alpha=0$ and $\alpha=0$; the two unitary identities make $U_0$ an ordinary
unitary. Since $U_0^*D_\star U_0$ is diagonal and $D_\star$ has distinct
eigenvalues, every column of $U_0$ is a phase multiple of one basis vector,
and unitarity makes the assignment a permutation. Right multiplication by
the inverse ordinary monomial unitary reduces to $U_0=I$, $\Lambda_0=D_\star$,
with both $U$ and $\Lambda$ of nonnegative support.

Compare coefficients in $A_\star U=U\Lambda$. For every $0<\gamma<\eta$,
$\Lambda_\gamma=0$ and $U_\gamma$ is diagonal, by well-ordered induction over
the union of the positive supports: at such $\gamma$,
\[
 D_\star U_\gamma-U_\gamma D_\star=\Lambda_\gamma+
 \sum_{\alpha+\beta=\gamma,\ \alpha,\beta>0}U_\alpha\Lambda_\beta,
\]
every $\beta$ on the right is smaller than $\gamma$, so the sum vanishes; the
left side has zero diagonal while $\Lambda_\gamma$ is diagonal, so
$\Lambda_\gamma=0$, and distinct diagonal entries make $U_\gamma$ diagonal.
This covers supports with infinitely many exponents below $\eta$. At
exponent $\eta$ the same equation acquires the term $Y$, and the lower-order
products still vanish, so $[D_\star,U_\eta]+Y=\Lambda_\eta$. The $(n,n+1)$
entry gives
\begin{equation}\label{ihs:fr:eq:unboundedgenerator}
 (U_\eta)_{n,n+1}=\frac{Y_{n,n+1}}{2^{-(n+1)}-2^{-n}}=-\frac{2^{n+1}}{n^2}.
\end{equation}
These entries are unbounded, and every entry of an ordinary bounded operator
is bounded by its norm, so $U_\eta\notin\B$, contradicting $U\in\Ahh$.
\end{proof}

This is a global, arbitrary-support strengthening of the first-order
obstruction \cref{ihs:rf:prop:bounded} of Part~II, which rules out a bounded
first coefficient $X$ in a formal expansion $U(t)=\Id+tX+\cdots$ for
$\diag(1/n)$ and the unweighted adjacency operator. Here the perturbation is
trace class, the exponents of $U$ are arbitrary, and the conclusion is about
any Hahn unitary in $\Ahh$.

\begin{proposition}[The eigenvalue list is not a bounded coefficient]\label{ihs:fr:prop:eigencoeff}
For $A_\star$, all odd powers of $s=t^\eta$ in each $\lambda_n$ vanish, and
\begin{equation}\label{ihs:fr:eq:eigensecond}
 \lambda_1=\tfrac12+4s^2+O(s^4),\qquad
 \lambda_n=2^{-n}+2^n\Bigl(\frac2{n^4}-\frac1{(n-1)^4}\Bigr)s^2+O(s^4)
 \quad(n\ge2).
\end{equation}
The coefficient of $s^2$ is at least $2^{n-1}/n^4$ for $n\ge16$.
Consequently $\diag(\lambda_n)$ is not an element of $\Ahh$, nor is any
permutation of this diagonal list.
\end{proposition}
\begin{proof}
Write an eigenvector with $n$th coordinate normalized to one as
$e_n+su^{(1)}+s^2u^{(2)}+\cdots$. The first coefficient equation gives
$(u^{(1)})_j=Y_{jn}/(d_n-d_j)$ for $j\ne n$, with $d_n=2^{-n}$, and
$[s]\lambda_n=Y_{nn}=0$. The $n$th coordinate of the second equation gives
$[s^2]\lambda_n=\sum_{j\ne n}Y_{nj}Y_{jn}/(d_n-d_j)$. Only adjacent indices
occur: for $n=1$ this is $4$, and for $n\ge2$ the two terms are
$-2^n/(n-1)^4$ and $2^{n+1}/n^4$. With $Je_n=(-1)^ne_n$ one has
$JD_\star J=D_\star$ and $JYJ=-Y$, so $A_\star(s)$ and $A_\star(-s)$ have the
same uniquely labelled simple eigenbranch at residue $d_n$, and
$\lambda_n(s)=\lambda_n(-s)$; this proves the odd-coefficient assertion
without any convergence argument in $\Gamma$. For $n\ge16$,
$(n/(n-1))^4\le(16/15)^4<3/2$, so the bracket in \eqref{ihs:fr:eq:eigensecond}
is at least $1/(2n^4)$, and the resulting positive sequence is unbounded. A
diagonal coefficient operator is bounded exactly when its entries are, and
permutation does not change that.
\end{proof}

These coefficients also rule out a common ordinary complex disc of analytic
eigenbranches. If all $\lambda_n(s)$ were analytic eigenvalue functions on
$|s|<r$ for some $r>0$, fix $0<\rho<r$. On $|s|=\rho$ every eigenvalue has
$|\lambda_n(s)|\le\norm{D_\star+sY}\le\norm{D_\star}+\rho\norm Y$.
Cauchy's coefficient estimate would then bound all $[s^2]\lambda_n$ by
$(\norm{D_\star}+\rho\norm Y)/\rho^2$, contradicting their unboundedness.
Each branch still has its own analytic neighbourhood at zero.

\subsection{Why Part~II has a different answer}\label{ihs:fr:sec:rowfinite}
For the locally finite path matrix in \eqref{ihs:fr:eq:counterexample},
each fixed formal perturbation coefficient of an eigenvector has finite
coordinate support: formal recursion propagates a bounded number of edges at
a fixed order. The coefficients can therefore be assembled as row- and
column-finite matrices, and the entries \eqref{ihs:fr:eq:unboundedgenerator}
are allowed there, since row- and column-finiteness does not bound their
sizes. This is why \cref{ihs:fr:thm:nounitary} does not contradict
\cref{ihs:rf:thm:unitary}: that theorem acts on $\C^{(\N)}((t^\Gamma))$, not
on $\ell^2(\N)((t^\Gamma))$, in a coefficient algebra without ordinary
boundedness, and \cref{ihs:fr:thm:nounitary} proves that its conclusion
cannot simply be transferred to $\Ahh$. Nor does the result deny all
spectral representations in other categories: it rules out the specified
global Hahn unitary and diagonal coefficient operator. A theory allowing
unbounded coefficient operators, different vector spaces or a different
summation notion would have to state its own domain and action requirements.

\section{A trace-class determinant with an exact Fredholm alternative}\label{ihs:fr:sec:determinant}

\emph{Generality of this section.} Here $H$ may be any ordinary complex
Hilbert space and $\Gamma$ any set-sized ordered abelian group, possibly
zero and not necessarily divisible. Set $\K=\C((t^\Gamma))$,
$\Hil=H((t^\Gamma))$ and $\Ahh=\B((t^\Gamma))$ as before. None of the
arguments in this section uses the compact residue $T$, self-adjointness,
separability or infinite dimension. For $H=\{0\}$ the determinant is $1$
and the Fredholm alternative is trivial; for $\Gamma=\{0\}$ the construction
is the ordinary determinant. Subsequent applications to $A=T+E$ retain
their stated spectral hypotheses.

\subsection{Trace class is imposed coefficient by coefficient}\label{ihs:fr:sec:traceclass}
Let $\SH$ be the ordinary complex trace-class ideal with its trace norm, and
\[
 \TT=\SH((t^\Gamma)),\qquad \TT_+=\{C\in\TT:v(C)\ge0\}.
\]
$\TT$ is a two-sided ideal of $\Ahh$ under finite coefficient convolution. It
is a specified coefficient ideal; no identification with all nuclear or
compactoid operators for a non-Archimedean norm is assumed. Define
\begin{equation}\label{ihs:fr:eq:hahntrace}
 \Tr_\Gamma C=\sum_\gamma(\Tr C_\gamma)t^\gamma.
\end{equation}
Ordinary traces are taken before forming the Hahn series; the support can
shrink by cancellation but cannot enlarge. Ordinary cyclicity at each of the
finitely many coefficient products gives
$\Tr_\Gamma(XC)=\Tr_\Gamma(CX)$ for $X\in\Ahh$ and $C\in\TT$.

The following classical inputs are not claimed as new. The map
$C\mapsto\det_H(I+C)$ on $\SH$ is entire; it is multiplicative on products of
operators of the form $I+C$, agrees with finite-dimensional determinants,
satisfies the finite-rank Sylvester identity, and is nonzero exactly when
$I+C$ is invertible; near $C=0$ it satisfies
$\det_H(I+C)=\exp(\Tr\log(I+C))$ \cite[Section~3]{Bornemann}.
For the Banach-space analyticity needed here, use the exterior-power
expansion $\det_H(I+C)=\sum_{n\ge0}\Tr(\bigwedge^n C)$: its terms are
continuous homogeneous polynomials and satisfy
$|\Tr(\bigwedge^n C)|\le\norm C_1^n/n!$. Thus the series converges locally
uniformly on trace-norm bounded sets and defines an entire map on $\SH$,
not merely an entire function along each scalar line.

The finite-rank identity used below is
$\det_H(I+UV)=\det_m(I_m+VU)$ for bounded $U:\C^m\to H$ and
$V:H\to\C^m$. Indeed, $UV$ takes values in the finite-dimensional space
$\Ran U$, so its Fredholm determinant reduces to that space; the usual
finite-dimensional rectangular determinant identity then gives the displayed
formula. This also covers $m=0$ with determinant $1$.

\begin{definition}[Hahn Fredholm determinant]\label{ihs:fr:def:determinant}
Every $C\in\TT_+$ decomposes uniquely as $C=C_0+C_+$ with $C_0\in\SH$ and
$v(C_+)>0$. Define
\begin{equation}\label{ihs:fr:eq:hahndet}
 \Det(I+C)=\sum_{k\ge0}\frac1{k!}\,\mathrm d^k{\det}_H(I+\cdot\,)(C_0)[C_+,\ldots,C_+],
\end{equation}
each multilinear expression being expanded into ordered coefficient words as
in \eqref{ihs:fr:eq:taylorlift}.
\end{definition}

This is well defined by \cref{ihs:lem:support}, with support in
$(\supp C_+)^*$. It is independent of any basis and needs no labelling of
eigenvalues. The ordinary coefficient $C_0$ may have large trace norm and
$I+C_0$ may be singular, so this is not merely a determinant on the
near-identity group $I+\TT_{>0}$.

\begin{proposition}[Determinant identities]\label{ihs:fr:prop:detidentities}
For $C,C'\in\TT_+$,
\begin{equation}\label{ihs:fr:eq:detidentities}
 \st\Det(I+C)=\det_H(I+C_0),\qquad
 \Det\bigl((I+C)(I+C')\bigr)=\Det(I+C)\Det(I+C').
\end{equation}
The canonical adjoint identity is
$\Det(I+C^*)=\overline{\Det(I+C)}$. If $\mathcal J$ is any fixed
conjugation on $H$ (an antilinear isometric involution), then also
$\Det(I+\mathcal JC\mathcal J)=\overline{\Det(I+C)}$, with
$\mathcal J$ applied coefficientwise. The determinant agrees with a finite
matrix determinant on a finite-dimensional ordinary summand. If $v(C)>0$,
\begin{equation}\label{ihs:fr:eq:detlog}
 \Det(I+C)=\exp\Bigl(\Tr_\Gamma\sum_{k\ge1}\frac{(-1)^{k+1}}kC^k\Bigr),
\end{equation}
all powers, logarithms and exponentials being strong Hahn sums.
\end{proposition}
\begin{proof}
The standard-part assertion is the constant Taylor coefficient. For
multiplicativity apply \cref{ihs:fr:lem:lift} to the ordinary identity
$\det_H(I+C+C'+CC')=\det_H(I+C)\det_H(I+C')$ near the ordinary pair
$(C_0,C'_0)$; the nonconstant increments of $C+C'+CC'$ have positive support.
The adjoint and fixed-conjugation assertions follow by real-analytic
coefficient comparison, and the finite-matrix assertion by holomorphic
comparison. No canonical conjugation of operators on an abstract complex
Hilbert space is assumed. For
$v(C)>0$, the ordinary logarithmic identity holds near zero in trace norm;
compare Taylor coefficients at zero and lift. Positive-word finiteness
justifies the operator logarithm and the scalar exponential without any
claim of sequential convergence in the value group.
\end{proof}

\subsection{The Fredholm alternative}\label{ihs:fr:sec:alternative}
The nonzero-determinant criterion needs more than residue testing: a
determinant with zero standard part can still be a nonzero Hahn number. The
argument reduces the singular ordinary residue to one finite matrix.

\begin{lemma}[Ordinary finite-rank correction]\label{ihs:fr:lem:correction}
For $C_0\in\SH$ there is a finite-rank ordinary bounded operator
$F=\iota\kappa$, with bounded linear maps $\iota:\C^m\to H$ and
$\kappa:H\to\C^m$, such that $I+C_0+F$ is
invertible.
\end{lemma}
\begin{proof}
$I+C_0$ is Fredholm of index zero. Decompose its domain as its
finite-dimensional kernel plus the orthogonal complement, and its target as
its closed range plus a finite-dimensional orthogonal complement of equal
dimension. Choose an isomorphism from the kernel onto that complement and
extend it by zero; adding it makes both block maps bijective and boundedly
invertible. Here the restriction from the kernel's orthogonal complement
onto the closed range has a bounded inverse by the ordinary bounded inverse
theorem. Take $m=\dim\ker(I+C_0)$, let $\kappa$ be the coordinates of the
orthogonal projection onto that kernel, and let $\iota$ be the chosen
isomorphism into the complementary target space. Both are bounded. If the
kernel is zero, take $m=0$ and $F=0$.
\end{proof}

\begin{theorem}[Hahn Fredholm alternative]\label{ihs:fr:thm:fredholm}
Let $C\in\TT_+$. Then
\begin{equation}\label{ihs:fr:eq:fredholmequivalence}
 \Det(I+C)\ne0\iff(I+C)^{-1}\in\Ahh\iff I+C:\Hil\to\Hil\text{ is bijective}.
\end{equation}
Every such $I+C$ is algebraically Fredholm of index zero; if the determinant
vanishes, $I+C$ has a nonzero finite-dimensional kernel and a cokernel of
the same dimension. Explicitly, with $F=\iota\kappa$ as in
\cref{ihs:fr:lem:correction}, $M_F=I+C+F$ and $L_F=I_m-\kappa M_F^{-1}\iota$,
\begin{equation}\label{ihs:fr:eq:finitefredholm}
 \Det(I+C)=\Det(M_F)\,\det_mL_F,\qquad \st\Det(M_F)\ne0,
\end{equation}
so one finite matrix decides the exact invertibility problem.
\end{theorem}
\begin{proof}
The constant coefficient of $M_F$ is invertible, so $M_F^{-1}$ exists in the
nonnegative part of $\Ahh$ by \cref{ihs:lem:neumann}. Let
$R_F=M_F^{-1}\iota$; then $I+C=M_F(I-R_F\kappa)$. Every coefficient of
$R_F\kappa=M_F^{-1}F$ is trace class, with nonnegative support.
Both $M_F-I$ and $-R_F\kappa$ therefore lie in the determinant's domain
$\TT_+$. To lift Sylvester, regard $R_F$ as a Hahn series with coefficients
in the Banach space $\mathcal B(\C^m,H)$ and keep $\kappa$ fixed. Apply
\cref{ihs:fr:lem:lift} to
$R\mapsto\det_H(I-R\kappa)=\det_m(I_m-\kappa R)$ at the ordinary
constant coefficient of $R_F$. Its coefficients need not have their ranges
in one common finite-dimensional ordinary subspace of $H$.
Multiplicativity and the lifted finite-rank Sylvester identity give
$\Det(I+C)=\Det(M_F)\Det(I-R_F\kappa)=\Det(M_F)\det_m(I_m-\kappa R_F)$, and
$\Det(M_F)$ has the nonzero standard part $\det_H$ of the constant
coefficient. If $L_F$ is invertible, the identity
$(I-R_F\kappa)^{-1}=I+R_FL_F^{-1}\kappa$ constructs an inverse in $\Ahh$
(the entries of $L_F^{-1}$ may have negative valuation, which $\Ahh$
permits). There are only finitely many matrix entries, so their supports
have a well-ordered finite union; multiplication by the bounded coefficient
maps $R_F$ and $\kappa$ then gives an element of $\Ahh$. Combined with
$M_F^{-1}$ this proves the first implication, and an
algebra inverse is a module inverse. If $\det_mL_F=0$, choose $y\ne0$ with
$L_Fy=0$; then $y=\kappa R_Fy$, the vector $x=R_Fy$ is nonzero, and
$(I-R_F\kappa)x=0$, so $(I+C)x=0$. Finally $R_F$ and $\kappa$ identify the
kernel and cokernel of $N_F=I-R_F\kappa$ with those of $L_F$: on kernels they
are inverse because $R_F\kappa=I-N_F$ and $\kappa R_F=I_m-L_F$, and on
cokernels the maps are $[x]\mapsto[\kappa x]$ and $[y]\mapsto[R_Fy]$.
The identities $N_FR_F=R_FL_F$ and $\kappa N_F=L_F\kappa$ make them well
defined, while $R_F\kappa=I-N_F$ and $\kappa R_F=I_m-L_F$ make the
composites identities modulo the respective images. Multiplication by the invertible $M_F$ does not change
the dimensions, so both are finite and equal. The empty correction when the
constant coefficient is already invertible is covered by $\det_0=1$.
\end{proof}

\begin{remark}[The exact domain matters]\label{ihs:fr:rem:domain}
The theorem applies to $C$ of nonnegative valuation, not to every series in
$\TT$. It allows a singular ordinary residue and arbitrary positive Hahn
corrections and produces inverses whose valuations may be negative. It does
not extend \eqref{ihs:fr:eq:hahndet} to negative-valuation trace-class
series; \cref{ihs:fr:sec:boundary} explains why that cannot be done by the
same infinite Taylor formula.
\end{remark}

\section{Trace and determinant without eigenvalue sums or products}\label{ihs:fr:sec:lidskii}

The trace-class counterexample yields an obstruction stronger than the
automatic failure of strong summability of infinitely many nonzero
constants: even if ordinary infinite sums are allowed at each Hahn
coefficient, spectral trace and product formulas fail. For $A_\star$ of
\eqref{ihs:fr:eq:counterexample} put $s=t^\eta$, $a_n=[s^2]\lambda_n$, and
$g_n=2^{n+1}/n^4$ for $n\ge1$, $g_0=0$; \cref{ihs:fr:prop:eigencoeff} says
exactly that $a_n=g_n-g_{n-1}$.

\begin{theorem}[Failure of the spectral trace sum]\label{ihs:fr:thm:nolidskii}
$\Tr_\Gamma A_\star=1$. However, $\sum_{n\ge1}\lambda_n$ does not exist even
under the convention of taking an ordinary complex sum separately at every
Hahn exponent: its second coefficient has the exact partial-sum law
\begin{equation}\label{ihs:fr:eq:traceboundary}
 \sum_{n=1}^Na_n=\frac{2^{N+1}}{N^4}\longrightarrow+\infty.
\end{equation}
This rules out a universal eigenvalue-sum formula for the trace in the
coefficientwise trace-class category $\TT$.
\end{theorem}
\begin{proof}
$\Tr D_\star=1$ and the off-diagonal trace-class $Y$ has trace zero. The
second-coefficient partial sum telescopes to $g_N$. The individual $a_n$ do
not tend to zero, so no rearrangement makes their sum convergent; the
spectral sum is undefined rather than equal to a different Hahn scalar.
\end{proof}

For comparison, truncate to the first $N$ basis vectors. The finite-matrix
eigenbranches agree to second order with $\lambda_n$ for $n<N$, but the last
branch loses the edge to $N+1$, so its second coefficient is
$a_N-g_N=-g_{N-1}$, and the finite-section second trace coefficient is
$\sum_{n<N}a_n+(a_N-g_N)=0$, as the finite trace identity requires. The
missing boundary term $g_N$ grows rather than tends to zero: stabilization of
each fixed branch is insufficient to exchange the branch sum with formal
coefficient extraction.

Put $\mathfrak d_0=\det_H(I+D_\star)=\prod_{n\ge1}(1+2^{-n})>0$ and
$q_n=1/\bigl(n^4(1+2^{-n})(1+2^{-(n+1)})\bigr)$; these products and sums are
ordinary real ones.

\begin{theorem}[Failure of a spectral product formula]\label{ihs:fr:thm:noproduct}
\begin{equation}\label{ihs:fr:eq:detsecond}
 [s^2]\Det(I+A_\star)=-\mathfrak d_0\sum_{n\ge1}q_n<0,
\end{equation}
whereas the coefficient of $s^2$ in $\prod_{n=1}^N(1+\lambda_n)$ tends to
$+\infty$. This spectral product has no coefficientwise ordinary limit and
cannot represent $\Det(I+A_\star)$.
\end{theorem}
\begin{proof}
Factor $I+D_\star+sY=(I+D_\star)(I+sK_0)$ with $K_0=(I+D_\star)^{-1}Y$, an
ordinary trace-class operator of trace zero. The determinant logarithm
\eqref{ihs:fr:eq:detlog} gives $[s^2]\Det(I+A_\star)=-\tfrac12\mathfrak
d_0\Tr(K_0^2)$. The diagonal of $K_0^2$ consists of adjacent two-step
contributions, and absolute convergence gives
$\Tr(K_0^2)=2\sum_{n\ge1}q_n$, finite and positive. Let
$P_N=\prod_{n=1}^N(1+2^{-n})$. Since all first eigenvalue coefficients
vanish, $[s^2]\prod_{n=1}^N(1+\lambda_n)=P_N\sum_{n=1}^Na_n/(1+2^{-n})$, and
summation by parts, or pairing the two ends of each edge, gives
\begin{equation}\label{ihs:fr:eq:productboundary}
 \sum_{n=1}^N\frac{a_n}{1+2^{-n}}=\frac{g_N}{1+2^{-N}}-\sum_{n=1}^{N-1}q_n.
\end{equation}
The first term tends to $+\infty$, the sum of the $q_n$ converges, and
$P_N\to\mathfrak d_0>0$.
\end{proof}

The same missing edge that obstructs the trace sum supplies the divergent
boundary term. Finite-section determinants remove that term: their second
coefficients converge to \eqref{ihs:fr:eq:detsecond}. In fact every coefficient
converges. The compressions of $D_\star$ and $Y$ converge in trace norm, so
the ordinary determinant perturbation estimate
\cite[Equation~(4.1)]{Bornemann} gives uniform convergence on each fixed
circle in the complex $s$-plane; Cauchy's formula then gives convergence of
each Taylor coefficient. Products of the labelled branches of the infinite
operator do not converge coefficientwise. There is no conflict with ordinary
Lidskii theory or the ordinary trace-class eigenvalue product: at each fixed
ordinary real parameter, $D_\star+sY$ is an ordinary compact self-adjoint
trace-class operator and the ordinary theory applies. What fails is
interchanging a countable spectral sum with a formal perturbation expansion
around infinitely many spectral gaps shrinking to zero, for want of a
uniform analytic labelling.

\section{Finite-coupling determinants and the accumulation boundary}\label{ihs:fr:sec:boundary}

\subsection{One coherent entire family on ordinary coupling space}\label{ihs:fr:sec:coupling}
Assume now that $A=T+E$ satisfies the hypotheses of \cref{ihs:fr:thm:main}
and also $A\in\TT_+$, so $T$ and every coefficient of $E$ are trace class.
Let $\mathcal O(\C)$ be the ordinary entire functions. For an ordinary complex
variable $z$, expand the ordinary determinant map at $-zT$ in the directions
$-zE_\gamma$; the coefficient of any Hahn exponent is a finite sum of such
derivatives and is entire in $z$, with common support in $(\supp E)^*$. This
defines a coherent family
\begin{equation}\label{ihs:fr:eq:couplingdet}
 \mathcal F_A\in\mathcal O(\C)((t^\Gamma)),\qquad
 \mathcal F_A(z)=\Det(I-zA)\quad(z\in\Ocal),
\end{equation}
where at $z=z_0+\delta$, with ordinary $z_0$ and $v(\delta)>0$, the entire
coefficients are evaluated by their Taylor series at $z_0$ using the
positive support of $\delta$; \cref{ihs:fr:lem:lift} identifies this with
\cref{ihs:fr:def:determinant} applied to $-zA$.

\begin{theorem}[Exact finite-coupling detection]\label{ihs:fr:thm:coupling}
For every finite $z\in\Ocal$,
$\mathcal F_A(z)=0$ if and only if $I-zA$ is not invertible in $\Ahh$. Its
finite zeros are exactly $\{\lambda^{-1}:\lambda\in\sigma_{p,\Gamma}(A)\}$;
in particular $\mathcal F_A(0)=1$ and no infinitesimal coupling is a zero.
\end{theorem}
\begin{proof}
For finite $z$, $-zA$ has nonnegative support with trace-class coefficients,
so \cref{ihs:fr:thm:fredholm} gives the equivalence. If $z=0$ there is
nothing to prove; if $v(z)>0$, $I-zA$ has a Neumann inverse. For $v(z)=0$,
invertibility is equivalent to that of $A-z^{-1}I$, where $z^{-1}$ has a
nonzero standard part, and \cref{ihs:fr:thm:main} makes it singular exactly
at the listed eigenvalues. Conversely every eigenvalue in
\eqref{ihs:fr:eq:mainpoint} has a finite inverse with nonzero standard part.
\end{proof}

The entire \emph{ordinary coefficient functions} do not mean that the
Taylor series can be evaluated at every infinite surcomplex argument.

\begin{proposition}[Negative-valuation evaluation can fail]\label{ihs:fr:prop:negeval}
Write the ordinary entire determinant of $D_\star$ as
$\prod_{n\ge1}(1-2^{-n}z)=\sum_{k\ge0}c_kz^k$. Every $c_k$ is nonzero, and
for $\eta>0$ the family $(c_kt^{-k\eta})_{k\ge0}$ is not strongly Hahn
summable; so the entire Taylor series does not define the determinant at
$t^{-\eta}$ by Hahn substitution.
\end{proposition}
\begin{proof}
Up to sign, $c_k$ is the $k$th elementary symmetric sum of the positive
numbers $2^{-n}$, finite (at most $(\sum_n2^{-n})^k/k!$) and strictly
positive. The substituted terms have exponents $0,-\eta,-2\eta,\ldots$, which
are not well ordered.
\end{proof}

This does not exclude a separately defined regularization or another function
category at infinite arguments; it excludes the termwise extension by the
entire Taylor series. Nonzero finite coupling $z$ corresponds to spectral parameter
$\lambda=z^{-1}$ away from the zero monad, and nonzero $\lambda\in\mK$ would
require infinite coupling, exactly where that substitution need not exist.

\subsection{No coherent analytic characteristic function through zero}\label{ihs:fr:sec:nochar}
Let $\mathbb D_r=\{z\in\C:\abs z<r\}$ with ordinary $r>0$. A \emph{coherent
Hahn analytic function} on its halo is a series
$\mathcal G(z)=\sum_{\gamma\in S}g_\gamma(z)t^\gamma$ with
$g_\gamma\in\mathcal O(\mathbb D_r)$ and one well-ordered support $S$,
evaluated at $z=c+\delta$ by ordinary Taylor expansion of its coefficients at
$c\in\mathbb D_r$ and positive-order Hahn substitution of $\delta$. Its halo
is $\mathbb D_r^\#=\{c+\delta:c\in\mathbb D_r,\ \delta\in\mK\}$.

\begin{lemma}[A nonzero coherent germ cannot vanish on a whole monad]\label{ihs:fr:lem:nonvanishing}
If $\mathcal G\in\mathcal O(\mathbb D_r)((t^\Gamma))$ is nonzero, there is
$\delta_0\in\Gamma_{\ge0}$ with $\mathcal G(x)\ne0$ whenever $0\ne x\in\mK$
and $v(x)>\delta_0$.
\end{lemma}
\begin{proof}
Let $m$ be the least vanishing order at $0$ among the nonzero coefficient
germs. Each such order is finite by the identity theorem on the connected
disc. Every coefficient is divisible by $z^m$ on that same disc, so $\mathcal G=z^m\mathcal H$
with $\mathcal H(0)\ne0$. Let $\alpha$ be the least exponent of a nonzero
coefficient function of $\mathcal H$ and $\beta=v(\mathcal H(0))\ge\alpha$.
All Taylor coefficients of $\mathcal H$ at zero have valuation at least
$\alpha$, so for $x\ne0$ of positive valuation the strong substitution gives
$v(\mathcal H(x)-\mathcal H(0))\ge\alpha+v(x)$; the substitution is legitimate
after extracting $t^\alpha$, by \cref{ihs:lem:support}. If
$v(x)>\beta-\alpha$, the bound exceeds $\beta$, so $\mathcal H(x)$ has the
leading term of $\mathcal H(0)$; since $x^m\ne0$, $\mathcal G(x)\ne0$. Take
$\delta_0=\beta-\alpha$.
\end{proof}

\begin{theorem}[No analytic characteristic determinant at the accumulation point]\label{ihs:fr:thm:nocharacteristic}
For $A$ as in \cref{ihs:fr:thm:main}, no coherent Hahn analytic function on
$\mathbb D_r^\#$, for any ordinary $r>0$, has zero set exactly
$\sigma^{\mathrm{adj}}_\Gamma(A)\cap\mathbb D_r^\#$.
\end{theorem}
\begin{proof}
The desired zero set contains all of $\mK$ by \cref{ihs:fr:cor:reservoir},
and a nonzero coherent function cannot vanish there by
\cref{ihs:fr:lem:nonvanishing}: choose any $\eta_0>0$ in the nonzero group
and evaluate at $t^{\delta_0+\eta_0}$, whose valuation exceeds the stated
threshold. This uses no Archimedean or cofinality condition.
The zero function fails too: a nonreal
ordinary $c$ with $0<\abs c<r$ lies outside the real ordinary spectrum of
$T$, so $A-cI$ is invertible by the second case of \cref{ihs:fr:thm:main}.
\end{proof}

The theorem concerns a scalar analytic characteristic function of the
\emph{spectral parameter} near zero. It does not conflict with the
finite-coupling family $\mathcal F_A$, which equals one at coupling zero, or
with the finite polynomial $\det(\lambda I-A_c)$ on an isolated nonzero halo.
It does not exclude set-theoretic indicators, meromorphic objects with a
singularity at zero, or function classes whose coefficients do not share an
ordinary neighbourhood; its content is the failure inside the named coherent
analytic category.

\section{Genuinely higher-rank examples}\label{ihs:fr:sec:rank}

\subsection{A scale invisible to every power of another scale}\label{ihs:fr:sec:invisible}
Take $\Gamma=\Q^2$ lexicographically, $\eta=(0,1)$ and $\beta=(1,0)$, so
$n\eta<\beta$ for every ordinary integer $n$. The strong identity
$(1-t^\eta)^{-1}=\sum_{n\ge0}t^{n\eta}$ holds, but its finite partial sums
do not converge to it in the full valuation topology: the remainder after
order $N$ has valuation $(N+1)\eta<\beta$ (compare \cref{ihs:sec:ranktwo}).
All proofs above use finite coefficient words, not such convergence. For a
noncommuting example at both scales take
\[
 A=D_\star+t^\eta Y+t^\beta\lvert e_1\rangle\langle e_1\rvert,
\]
with $D_\star,Y$ from \eqref{ihs:fr:eq:counterexample}. The main spectral
theorem and the determinant theorem apply; the obstruction at exponent
$\eta$ is unchanged, and the second eigenvalue coefficients at $2\eta$ are
unchanged because $\beta>n\eta$ for every finite $n$. The failures of global
diagonalization and of the spectral trace sum therefore persist in a
genuinely higher-rank perturbation. Identifying $(a,b)$ with $a\omega+b$
gives $t^\eta=\omega^{-1}$ and $t^\beta=\omega^{-\omega}$: the example has
actual surcomplex coefficients involving both scales.

\subsection{Finite clusters may need exponents outside the input monoid}\label{ihs:fr:sec:monoid}
Keep $\Gamma=\Q^2$, $\eta=(0,1)$ and $\beta=(1,0)$ from the preceding
example. Let the first two ordinary coordinates share the eigenvalue $c=1$, let the
remaining eigenvalues be $2^{-1},2^{-2},\ldots$, and perturb only the first
block by
\begin{equation}\label{ihs:fr:eq:twoscaleblock}
 A_c=\begin{pmatrix}c+t^\eta&t^\beta\\t^\beta&c-t^\eta\end{pmatrix}.
\end{equation}
Its eigenvalues are exactly
\begin{equation}\label{ihs:fr:eq:twoscaleeigen}
 \lambda_\pm=c\pm\sqrt{t^{2\eta}+t^{2\beta}}
 =c\pm t^\eta\sum_{k\ge0}\binom{1/2}{k}t^{2k(\beta-\eta)},
\end{equation}
the binomial series being strongly summable because $\beta-\eta>0$. The
first two nonconstant exponents are $\eta$ and $2\beta-\eta$; the second is
not in the nonnegative integer monoid generated by $\eta$ and $\beta$,
although it lies in the original group. The cluster matrix has support
controlled by the input positive monoid, its eigenvalues lie in the original
divisible field, and their supports need not stay in that monoid; the last
claim would be false and is not part of \cref{ihs:fr:thm:cluster}. This is a
repeated-residue phenomenon of the same kind as the finite example cited in
\cref{ihs:rf:rem:monoid}. Under the surreal identification,
$\lambda_\pm=1\pm\omega^{-1}\bigl(1+\tfrac12\omega^{-2(\omega-1)}
-\tfrac18\omega^{-4(\omega-1)}+\cdots\bigr)$, with set-sized normal-form
supports.

\section{Part IV: proof dependencies, limits and further questions}\label{ihs:fr:sec:limits}

\subsection{Dependencies}\label{ihs:fr:sec:dependencies}
The main spectrum theorem uses ordinary compact self-adjoint spectral theory,
finite-word support, the inherited range factorization and finite Hermitian
diagonalization. The obstruction uses only its simple-eigenvalue instance and
a first-coefficient commutator equation. The spectral trace and product
failures additionally use the explicit second-order recurrence and
elementary finite summation identities. The determinant theorem uses the
ordinary trace-class determinant as an analytic map, finite-word Taylor
lifting, the ordinary index-zero Fredholm alternative and finite-dimensional
linear algebra. The analytic accumulation obstruction uses only a common
coefficient neighbourhood, the least ordinary vanishing order and a
valuation comparison, and is independent of trace-class assumptions. The
matrix identities, finite recurrence and support algebra are natural
formalization targets; the ordinary compact and trace-ideal inputs would
remain substantial analytic prerequisites in a Lean development. No Lean
files are included and no theorem here is represented as machine checked.

\subsection{Precisely limited claims}\label{ihs:fr:sec:nonclaims}
This part does not classify spectra of arbitrary elements of $\Ahh$; its
complete formula requires an injective compact self-adjoint ordinary residue,
a self-adjoint positive-order perturbation, a separable infinite-dimensional
$H$ and a divisible $\Gamma$. It does not supply a positive spectral measure,
a general non-Archimedean Hilbert-space spectral theorem, or a
nuclear-operator theory independent of the chosen coefficient model. Its
determinant is defined on the nonnegative-valuation component of the
coefficientwise trace-class ideal, not on every trace-class-valued Hahn
series. The no-product theorem concerns products of the labelled Hahn
eigenbranches and ordinary coefficientwise limits; it does not claim that all
conceivable regularized products are impossible. The no-characteristic
theorem is confined to the named coherent analytic category. The
counterexample to global diagonalization is confined to bounded ordinary
operator coefficients, not to the row- and column-finite algebra of Part~II
or to every possible unbounded-operator framework. No complete
classification of bounded-coefficient Hahn operators, no equivalence with
intrinsic nuclear or compactoid operator ideals, and no determinant for every
negative-valuation trace-class series is claimed. The finite checks of
\cref{ihs:sec:checksfour} are not evidence of priority and do not replace
the infinite-dimensional arguments; the part has not been independently
refereed.

\subsection{Further questions}\label{ihs:fr:sec:questions}
First, what uniform spectral-gap and coefficient-ideal hypotheses ensure that
all cluster rotations assemble into a bounded-coefficient Hahn unitary? The
first commutator equation requires a bounded solution of
$[D,X]=-B_{\mathrm{off}}$ for the diagonal residue $D$ and the off-diagonal
first perturbation coefficient $B_{\mathrm{off}}$; that necessary first-order
condition is not asserted to be sufficient at every higher order. Second,
which altered coefficient categories support a meaningful regularized
spectral product compatible with the Fredholm determinant when the ordinary
coefficientwise product fails? Third, how should the classification change
when the compact residue has a nontrivial kernel, so that genuine
infinitesimal eigenvalues can coexist with the range defect? These are
directions beyond the proved statements, not missing steps in their proofs.

\emph{Status (batch 34).} The first question is answered in exact-criterion
form, not in the form posed, by \texttt{hgeo:us:thm:main} of
\path{docs/surcomplex/hahn-hilbert-geometry/}. One bounded-coefficient unitary
conjugating every residual cluster projection to its deformed cluster
projection is a simultaneous straightener of the family of cluster
projections, and such a unitary exists exactly when the canonical coefficient
maps of the family are bounded at every exponent and their active support is
well ordered. The first-order equation there (\texttt{hgeo:us:prop:first}) is
the commutator equation above written for the projection family; the missing
higher-order conditions are the boundedness at every exponent together with the
well-ordering. In terms of spectral gaps and coefficient ideals of the
operator the question stays open, and the second and third questions are not
addressed there.

\clearpage
\part*{Shared back matter}
\addcontentsline{toc}{part}{Shared back matter}

\section{The four exact check suites}\label{ihs:sec:checks}

The four parts arrived with separate verification programs, which are kept
separate in \texttt{code/} and \texttt{data/} under the numbers \texttt{02}
(Part I), \texttt{03} (Part II), \texttt{04} (Part III) and \texttt{05}
(Part IV). No program checks any statement of another part.

\begin{warning}[A coincidence, not a dependence]\label{ihs:warn:coincidence}
The Part I and Part II suites both report exactly \textbf{218 exact checks},
and both recorded runs used \textbf{Python 3.13.5} with \textbf{SymPy 1.14.0}.
This is a coincidence of two independently written programs, and the shared
tool versions are a coincidence of two runs on the same day. The two counts
are of entirely different assertions: the Part I total is
$218$ assertions in six groups about finite convolutions, adjoints,
truncated Neumann identities, the isolated-point resolvent, a finite
rectangular analogue, finite sections and a backward-shift residual, while
the Part II total is $218$ checks about perturbation jets, an accumulating
chain, sharp boundary coefficients, bidegree truncations, projection
identities and a closed-form eigenpair. Neither number is evidence about the
other suite, and neither program copied the other. The Part III and Part IV
suites report $484$ and $1{,}268$ checks; their recorded runs also used
Python 3.13.5 (Part IV with SymPy 1.14.0, Part III with the standard library
only). The four counts are counts of program assertions, not of theorems, and
are not comparable with one another.
\end{warning}

\subsection{The Part I suite}\label{ihs:sec:checksone}
The program \texttt{code/02-hahn-hilbert-spectral-verify.py} uses exact
rational and symbolic arithmetic. It checks finite-support convolution and
adjoints, including a lexicographically ordered two-dimensional exponent
group; truncated noncommutative Neumann identities; the isolated-point
resolvent formula \eqref{ihs:hh:eq:isolatedres}; the algebraic left-inverse
and complement projection identities of \cref{ihs:hh:thm:defect} in a finite
rectangular analogue; the finite-section formulas of
\cref{ihs:sec:finitesections}; and the backward-shift residual of
\cref{ihs:hh:ex:backshift}. The recorded run passed 218 exact assertions in
six groups; its output is
\texttt{data/02-hahn-hilbert-spectral-verification.txt}.

The finite rectangular analogue is deliberately identified as such. An
injective nonsurjective endomorphism cannot exist on a finite-dimensional
space, so a finite square-matrix test cannot validate the infinite range
defect. Nor can finite tests establish Baire category, arbitrary Hahn
support well-ordering, an infinite cokernel dimension, or the spectral
formula. Those assertions rest on the proofs, not on the program.

\subsection{The Part II suite}\label{ihs:sec:checkstwo}
The program \texttt{code/03-infinite-spectral-verify.py} passed
\textbf{218 of 218} exact checks; the machine-readable record is
\texttt{data/03-infinite-spectral-verification.json}. The suite verifies a
complex Hermitian three-by-three example with two perturbation orders
through degree six; the six-vertex accumulating-level chain through degree
ten; the sharp finite-section coefficient for path distances one through
four; a two-variable example through total degree four, including the
coefficients in \eqref{ihs:rf:eq:ranktwocoeffs}; projection idempotence and
self-adjointness at every retained bidegree; the symbolic inverse-gap
obstruction of \cref{ihs:rf:prop:bounded}; and the recurrence, quadratic
identity, normalization generating function, and Taylor coefficients of the
closed-form eigenpair of \cref{ihs:rf:prop:closed}.

These are exact algebraic checks, not floating-point evidence for an
infinite theorem. They are not a Lean formalization, do not certify novelty,
and cannot replace the summability arguments. The script raises an exception
at the first failed check. The two-scale test uses finite bidegrees, not an
allegedly finite initial segment of a lexicographic valuation group.

\subsection{The Part III suite}\label{ihs:sec:checksthree}
The program \texttt{code/04-drazin-halos-verify.py} uses exact
\texttt{fractions.Fraction} arithmetic and only the Python standard library.
Its recorded run, \nolinkurl{data/04-drazin-halos-verification.json}, passed
\textbf{484} assertions in six categories: pure nilpotent inverses (30),
mixed Drazin decompositions after a nonorthogonal rational change of basis
with left and right residuals of truncated regular resolvents (175),
nilpotent pole-order transformation under ramified polynomials (210), local
ramified germs and their compositional inverses (18), nonmonomial shifted
nilpotent inverses (24), and finite-order instances of the straightening
inverse construction (27). An exact finite matrix identity is verified for
the tested matrix only, and a truncated formal identity to its recorded
order only. Neither the Volterra spectrum nor the increasing-Jordan-block
counterexample is established by these tests; their evidence is the proof in
\cref{ihs:dz:sec:examples}. The program prints its report and writes a file
only when given \texttt{--output}.

\subsection{The Part IV suite}\label{ihs:sec:checksfour}
The program \texttt{code/05-fredholm-verify.py} uses exact rational and
symbolic arithmetic with SymPy and the deterministic seed \texttt{20260922}.
Its recorded run, \nolinkurl{data/05-fredholm-verification.json}, passed
\textbf{1,268} checks in eighteen categories. They cover exact eigenpair
recursions through order six for $24$ finite symmetric systems, with
residuals, normalization and the second-order Rayleigh coefficient; truncated
coefficient identities for $U^*U=I$ and $UP^\circ=PU$ in the direct rotation
of \cref{ihs:fr:lem:rotation}; the determinant--exponential-trace identity,
finite-rank Sylvester identities and the singular-residue finite-rank
correction; the forced generator entries
\eqref{ihs:fr:eq:unboundedgenerator}, the second eigenvalue coefficients,
vanishing odd coefficients and the growth bound of
\cref{ihs:fr:prop:eigencoeff}; the telescoping laws
\eqref{ihs:fr:eq:traceboundary} and \eqref{ihs:fr:eq:productboundary} and
their finite-section cancellations; and finite instances of $n\eta<\beta$ in
lexicographic order. No finite run verifies arbitrary-rank support
well-ordering, existence of all Hahn coefficients, continuum-dimensional
cokernels, spectral exhaustion over the whole field, or the absence of all
possible global unitaries. The program writes its report to
\texttt{data/verification.json} relative to the working directory unless
\texttt{--output} is given, so it should be run on a copy.

\subsection{Building}\label{ihs:sec:building}
The \LaTeX\ source contains its bibliography and needs no external images,
fonts, or bibliography database. A standard build is
\begin{verbatim}
pdflatex -interaction=nonstopmode -halt-on-error article.tex
pdflatex -interaction=nonstopmode -halt-on-error article.tex
pdflatex -interaction=nonstopmode -halt-on-error article.tex
python code/02-hahn-hilbert-spectral-verify.py
python code/03-infinite-spectral-verify.py
python code/04-drazin-halos-verify.py
python code/05-fredholm-verify.py --output rerun/05-verification.json
\end{verbatim}
The scripts support Python 3.10 or later; those of Parts I, II and IV need
SymPy, and that of Part III needs only the standard library. Some of them
write their reports next to themselves or into the working directory, so run
them on a copy of this directory. The delivered \texttt{build.sh} wrappers
refer to the source packages' original unnumbered file names and do not run
as placed. The source is not Lean code, and no proof-assistant certification
is claimed for any part.

\appendix
\section{Part I: repository and literature audit}\label{ihs:app:auditA}

\subsection{Repository snapshot}\label{ihs:app:snapshotA}
The Part I coverage comparison is pinned to
\begin{center}\small\texttt{\commitA}.\end{center}
It inspected the repository tree, the project \texttt{README.md}, the full
\texttt{docs/README.md} catalogue, the opening 250 lines of
\texttt{docs/surcomplex/spectral-theory/article.tex}, and the opening 230
lines of the differential-equations article. This was a targeted scope
audit, not a claim to have reread every archived source manuscript or
verified every existing Lean file. The spectral opening explicitly states
its finite-dimensional scope, while the catalogue maps the already-covered
analytic, differential, Berkovich, foundational and computational strands.
No repository files were modified by that audit.

\subsection{Primary-source comparison and date boundary}\label{ihs:app:datesA}
The Part I literature search was carried out through September 22, 2026. The
publisher records and accessible abstracts for \cite{ANS,Ishizuka}, the
author-provided text associated with \cite{AN}, and the arXiv records for
\cite{Boasso,Miabey} were used to locate nearby topics. Ishizuka's paper
appears in volume 39(2), labeled 2024, with the publisher recording online
publication on January 3, 2025. Both dates are recorded in the bibliography
to avoid silently conflating the issue year with the online date.

The comparison supports the modest statement that the exact theorem package
of Part I was not identified in the inspected sources. It does not establish
that no earlier work contains any equivalent result. In particular, formal
support inversion, failure of non-Archimedean orthogonality, and ordinary
resolvent expansions have substantial pre-existing literatures.

\section{Part II: provenance and research scope}\label{ihs:app:auditB}

\subsection{Repository snapshot}\label{ihs:app:snapshotB}
The Part II comparison is pinned to
\begin{center}\small\texttt{\commitB}\end{center}
with commit timestamp 2026-09-22T15:58:14Z as returned by GitHub. It used
the documentation catalogue and the spectral-theory README at that commit;
the differential-equations and spectral-theory inventories were also
examined to avoid repeating finite-dimensional work. The pinned
spectral-theory scope explicitly restricts dimensions to ordinary finite
integers. This is a targeted scope comparison, not a line-by-line audit of
every theorem in the repository. Initially returned unpinned documentation
was older; the provenance is based on the pinned snapshot above. No change
was made to the GitHub repository.

The two audits used two different pinned revisions, and neither was repeated
against the other's pin when the two manuscripts were assembled here.

\subsection{Primary sources consulted}\label{ihs:app:sourcesB}
Gonshor for foundational normal forms \cite{Gonshor}; Neumann
\cite{Neumann}; Higman \cite{Higman}; Greenbaum, Li and Overton \cite{glo};
Kenmoe, Smerlak and Zadorin \cite{ksz}; Giscard, Thwaite and Jaksch
\cite{pathsum}; and Ara, Goodearl and O'Meara \cite{ago}, particularly for
the operator-category distinctions and failures of general
infinite-dimensional spectral assertions.

The source list and comparison do not certify the absence of all prior
versions of the new package. The elementary perturbation coefficients, the
formal recursion mechanism, the Neumann support lemma, the Higman word
lemma, the walk interpretation of products, and the broad phenomenon of
positive infinite matrices with no eigenvectors are not claimed as
discoveries.

\subsection{Candidate contribution}\label{ihs:app:candidateB}
The Part II candidate-original contribution is the global arbitrary-rank
Hahn operator/module synthesis: unique normalized infinite diagonalization
with support control, all coherent resolvents, a complete commuting
projection calculus, marked-walk stability, and sharp finite-section error
thresholds, under simple ordinary residues and positive global support. The
explicit accumulating-level closed-form eigenpair of
\cref{ihs:rf:prop:closed} is proved directly and used as an independent
check; no separate priority claim is made for that solvable model.

\subsection{A proof not reprinted}\label{ihs:app:higmannote}
The Part II source manuscript carried an appendix with the classical
minimal-bad-sequence proof of Higman's lemma, specialized to a well-ordered
alphabet, so that the summability argument would be auditable. That proof is
not reprinted here, because the identical argument is already written out in
this collection, as \texttt{spec:lem:words} and \texttt{spec:lem:neumann} of
the finite-dimensional spectral report \cite{RepoSpectral}.
\Cref{ihs:lem:support} above states what both parts use, and strict
positivity of the letters in an ordered group is an additional hypothesis
used only after the word-order result, to pass from embeddings to
comparisons of sums.

\section{Part III: repository and literature audit}\label{ihs:app:auditC}

\subsection{Repository snapshot}\label{ihs:app:snapshotC}
The Part III source compared its results with the repository at
\begin{center}\small\texttt{\commitC},\end{center}
reviewed on September 22, 2026 \cite{RepoPinC}. It consulted the root
README, \texttt{docs/README.md}, \texttt{docs/manifest.tex}, the
incoming-report inventory, and the README and source of this report as it
then stood, with Parts I and II only. It identified Part~I's
constant-normal spectral formula as the closest existing result and credited
it as an overlap. It noted that the report's related-work paragraph already
mentioned Drazin spectra and cited Boasso, so that a claim that the
repository contained no Drazin mention would have been false. A GitHub code
search for ``Drazin'' returned no matches but marked its results incomplete,
and was not used as evidence of absence; a search of the retrieved catalogue
found no occurrence, and the retrieved spectral source contained the
related-work mention. The novelty assessment rests on comparison of
statements and hypotheses. No repository file was changed. The delivered
audit is kept unchanged as \texttt{04-drazin-halos-RESEARCH\_AUDIT.md}.

\subsection{Primary sources}\label{ihs:app:sourcesC}
The bibliographic record of Drazin's paper \cite{Drazin} (the audit does not
claim a full reading of it); the relevant portions of Boasso's two papers
\cite{BoassoOperators,Boasso}, including definitions, the resolvent-pole
characterization and spectral mapping; Koliha's distinction between
finite-index and generalized Drazin inverses \cite{Koliha}; Gonshor and van
den Dries--Ehrlich for normal forms and field foundations
\cite{Gonshor,DriesEhrlich}; and Kaplan--Krapp--Serra for contemporary Hahn
and surreal automorphism context \cite{KaplanKrappSerra}, not credited with
the spectral statements. Targeted searches combining Hahn, Drazin, Laurent
inverse, surreal spectrum and ramified spectral mapping did not locate the
combined theorem package; this does not establish priority.

\subsection{Statements of the source that are stale or had to be corrected}\label{ihs:app:staleC}
\begin{itemize}
\item \emph{Divisibility in Part I.} The source says that Part~I works
 ``under a standing divisibility assumption'' (its audit: ``with
 divisibility as a standing convention''), and presents the absence of
 divisibility as one of the hypotheses Part~III drops. At its pin
 \texttt{048b72c} this was accurate: Part~I then kept divisibility as a
 standing convention although its norm did not need it. Commit
 \texttt{6e34651} removed that assumption from Part~I
 (\cref{ihs:hh:sec:field}) before Part III was placed. Relative to the
 present Part~I, Part~III drops normality and the Hilbert and Banach
 structure, not divisibility (\cref{ihs:dz:rem:recover}).
\item \emph{README line numbers.} The audit locates the normal formula at
 lines 67--70 of this report's README at the pin. That was correct there;
 the README has since been rewritten, and the formula is
 \cref{ihs:hh:thm:spectrum}.
\item \emph{What the report contains.} The audit describes this report as two
 inequivalent theories, which was accurate at the pin. The report now also
 contains Parts III and IV, and the related-work paragraph it cites
 (\cref{ihs:sec:nearby}) now points to Part~III.
\item \emph{Notation.} The source's letters $A$, $D$, $P$, $Q$, $B$, $V$,
 $h$, $r$, $K_\Gamma$ and $\mathfrak m_\Gamma$ are renamed as stated at the
 start of Part~III; no statement changes.
\end{itemize}

\section{Part IV: repository and literature audit}\label{ihs:app:auditD}

\subsection{Repository snapshot}\label{ihs:app:snapshotD}
The Part IV source is pinned to the same revision \texttt{048b72c} as
Part~III and was inspected on September 22, 2026; its build and verification
timestamps fall on September 23 in UTC. It retrieved, through a GitHub
connector, the root \texttt{README.md}, \texttt{docs/README.md}, this
report's README and source, and the READMEs of the dynamics and
analytic-geometry reports, and checked the introduction and relevant scope
passages of this report. It identified Part~I's constant-normal formula and
positive-order range-defect factorization as predecessors, not discoveries.
The comparison was targeted, not a review of every file or proof, and no
repository file was modified. The delivered record is kept unchanged as
\texttt{05-fredholm-PROVENANCE.md}.

\subsection{Primary sources and search scope}\label{ihs:app:sourcesD}
Conway and Gonshor for surreal normal forms \cite{Conway,Gonshor}; Higman
\cite{Higman}; Poonen, Corollary~4, for Hahn-field algebraic closedness
\cite{Poonen}; Kato for compact spectral and isolated-cluster perturbation
inputs \cite{Kato}; Simon for the direct-rotation formula
\cite{SimonProjections}; Bornemann, Section~3, for ordinary trace-class
determinant identities \cite{Bornemann}; and Serre for prior
non-Archimedean Fredholm theory in a different operator category
\cite{Serre}. The Bornemann, Simon and Poonen material was also inspected
through page images. Targeted searches combined Hahn, surreal, compact
operator, trace class, Fredholm determinant, Lidskii and noncommuting
perturbation; exact-phrase searches for ``Hahn series'' with ``Fredholm
determinant'' and for ``Hahn'' with ``Lidskii'' produced no directly
matching treatment. That negative outcome is not a proof of priority, and
unrelated search snippets were not used as evidence.

\subsection{Statements of the source that are stale or had to be corrected}\label{ihs:app:staleD}
\begin{itemize}
\item \emph{Trace class.} The source says this report ``explicitly excludes
 a trace-class theory''. That was true at the pin and remains true of
 Parts I--III. The report now contains Part~IV's coefficientwise trace and
 Fredholm determinant, which is still not a general trace-class theory
 (\cref{ihs:sec:nonclaim}).
\item \emph{Constant spectra.} The source describes the bounded-coefficient
 part as computing the spectrum of constant normal operators. The report now
 also computes the adjointable-algebra spectrum of every constant bounded
 operator (Part~III).
\item \emph{Divisibility.} The source's divisibility hypothesis is kept for
 the spectral classification. The subsequent proof review removes it, and
 the Hilbert dimension and separability restrictions, from the independent
 determinant and Fredholm alternative in \cref{ihs:fr:sec:determinant}.
 This extension is justified there, not attributed to the delivered source.
 Unlike the Part III source, that source makes no claim about Part~I's
 hypotheses on $\Gamma$.
\item \emph{Duplicated inputs.} The source re-proves the range-defect
 factorization, the support lemma and the Neumann inverse. They are printed
 once (\cref{ihs:hh:thm:defect,ihs:lem:support,ihs:lem:neumann}); the
 source's different proof of the first is kept as
 \cref{ihs:fr:rem:secondproof}, and its deduction of the word lemma from
 Higman's theorem as a paragraph after \cref{ihs:lem:support}.
\item \emph{Notation.} The source's $\mathscr B$, $\sigma_{\mathscr B}$,
 $\sigma_{\mathrm{bij}}$ and $\operatorname{res}$ are written $\Ahh$,
 $\sigma^{\mathrm{adj}}_\Gamma$, $\sigma^{\mathrm{alg}}_\Gamma$ and $\st$;
 its letters $D$, $B$, $P_0$, $W$ and $M$ (rotation), $U,V$ (finite-rank
 factors), $M$, $L$, $R$ (Fredholm reduction), $d_0$, $d_A$ and $G$ become
 $D_\star$, $Y$, $P^\circ_c$, $\Omega$ and $I-(P-P^\circ)^2$, $\iota,\kappa$,
 $M_F$, $L_F$, $R_F$, $\mathfrak d_0$, $\mathcal F_A$ and $\mathcal G$; no
 statement changes.
\end{itemize}

\section{How this report was assembled}\label{ihs:app:provenance}
The report is built from \textbf{four} research manuscripts, all dated
September 22, 2026. Two arrived together and were assembled first as Parts
I and II; two arrived in a later batch and were added as Parts III and IV.
\begin{center}\small
\begin{tabularx}{\textwidth}{@{}p{.05\textwidth}p{.31\textwidth}p{.09\textwidth}Y@{}}
\toprule
Part & Files & Pin & Contribution\\
\midrule
I & \texttt{02-hahn-hilbert-spectral-*} & \texttt{608dd23} & Hahn--Hilbert spaces: automatic adjoints, defect rigidity, normal spectral formula, coercive defects, norm attainment\\
II & \texttt{03-infinite-spectral-*} & \texttt{a3124af} & Row- and column-finite Hahn matrices: diagonalization, synthesis, commutants, algebra-relative spectrum, walks\\
III & \texttt{04-drazin-halos-*} & \texttt{048b72c} & Drazin halos in any unital algebra, exact resolvent valuation, ramified spectral mapping, Volterra and block examples\\
IV & \texttt{05-fredholm-*} & \texttt{048b72c} & Compact self-adjoint operators under noncommuting perturbations; obstruction to global diagonalization; Hahn Fredholm determinant\\
\bottomrule
\end{tabularx}
\end{center}
The source manuscripts themselves are not shipped; their code, data and
provenance records are, under the numbers above. Where the merge had to
choose:
\begin{itemize}
\item Results duplicated between sources are printed once: the support
 lemma and Neumann inverse (all four sources), and the range-defect
 factorization (Parts I and IV), with the Part IV proof kept as a second
 route. The Part III normal-case corollary is kept as a second proof of the
 algebra half of \cref{ihs:hh:thm:spectrum}, not merged into it.
\item Parts III and IV are kept as separate parts, not merged into Part~I,
 because Part III's spectrum is relative to an arbitrary algebra and Part IV's
 spectral classification assumes divisibility and a separable
 infinite-dimensional $H$, which Part~I does not. Its determinant section
 now states its independently broader hypotheses.
\item Letters with an existing meaning in Parts I and II were renamed in the
 new parts (\cref{ihs:app:staleC,ihs:app:staleD}); no existing label or
 statement of Parts I and II was changed, and cross-references to the new
 parts were added to them.
\item The comparison of the algebra-relative spectra (\cref{ihs:prop:unify})
 and the cross-references between parts were made in assembly and are
 marked as such; they add no hypothesis to any source theorem and no result
 to any source.
\end{itemize}

\section{Dependency and proof-obligation ledgers}\label{ihs:app:ledgers}

\subsection{Part I}\label{ihs:app:ledgerA}
\begin{table}[ht]
\centering\small
\begin{tabularx}{\textwidth}{@{}p{.27\textwidth}YY@{}}
\toprule
Result & Inputs & Independent restriction\\
\midrule
Positive inner product and norm & Finite Hahn convolution; even leading exponent; binomial root & No divisibility needed\\
Automatic adjoint structure & Ordinary closed graph; Baire; Hahn support & Everywhere-defined adjoints\\
Defect splitting & Injective leading term; Neumann support lemma & No topology or commutation needed\\
Normal spectral formula & Defect splitting; ordinary normal spectral theorem & Constant normal operator\\
New eigenvalue exclusion & Orthogonal eigenspace projection & Normality is essential\\
Continuum cokernel & Leading-coefficient obstruction; power-series divergence test & Ordinary $\ell^2$ coefficients\\
Closed coercive range & Continuous defect projection; kernel identity & Arbitrary value group\\
Least norm-bound criterion & Ordinary positive square root; leading norm ratio & Constant bounded operator\\
\bottomrule
\end{tabularx}
\caption{Where each Part I hypothesis enters; ``complete'' does not replace
``orthogonally complete.''}\label{ihs:tab:ledgerA}
\end{table}

All general statements in \cref{ihs:hh:thm:package} have proofs in the text,
relative to the named classical inputs. No computational observation is used
as a premise of those proofs. The most consequential steps to audit
independently are the Baire uniform-support argument, the complement
projection in \cref{ihs:hh:thm:defect}, and the accumulation-point case in
\cref{ihs:hh:thm:spectrum}; their details are given without an appeal to a
numerical experiment.

\subsection{Part II}\label{ihs:app:ledgerB}
The logical dependence of the Part II claims is as follows. The Hahn algebra
and inner product depend on finite coefficient convolution, not on any
spectral theorem. Ordered-word finiteness licenses the homogeneous
construction. Simple residues make the divided commutator invertible on
off-diagonal coefficients. Self-adjointness and positivity make the
constructed eigenvalues real and permit near-one unitary normalization. The
invertibility already proved for the normalization is what supplies
two-sided unitarity. Synthesis uses vector-valued strong sums. Commutant
classification uses only the pairwise nonzero eigenvalue differences. The
full spectrum theorem adds a separate inverse-coherence argument. Locality
uses the homogeneous formulas and finite weighted walks, and surcomplex
transport uses the classical normal-form embedding.

The Part II results do not assert a solution to a named published open
problem; uniqueness or priority of every underlying idea; existence of a
bounded operator-norm diagonalizer; a Hilbert-space or probability
interpretation; convergence of the Hahn series in the fine surreal topology
or in every intrinsic valuation topology; an operator-valued Hahn sum of all
spectral projections; a spectral theorem for arbitrary repeated residues;
coherence of arbitrary entrywise Hahn matrices; decidability for arbitrary
surreal presentations; or a completed Lean formalization.

\subsection{Part III}\label{ihs:app:ledgerC}
The positive-support lemma supports the straightening automorphism and the
formal substitutions. The elementary Laurent recurrence supports the Drazin
criterion. Straightening plus cyclic-support projection yields inverse
descent, which yields the halo formula and the exact scale. The power image
and the formal inverse germ, with the halo formula and the commuting-factor
lemma, yield the ramified image, and the finite-union subgroup lemma turns
that image into the necessary and sufficient criterion. Among the general
results, only the Banach section and the pole-order profile
(\cref{ihs:dz:sec:banach,ihs:dz:prop:polemap}) use ordinary analytic spectral
theory; the Volterra and block examples also use ordinary operator-norm
estimates and convergence. Only the surreal specialization
uses Conway normal forms. The steps most worth auditing independently are
the surjectivity argument in \cref{ihs:dz:lem:straighten} and the converse
direction of \cref{ihs:dz:lem:laurent}.

The Part III results do not assert a classification of nonconstant Hahn
operators, of all linear endomorphisms, or of bijectivity spectra beyond the
worked examples; algebraicity for non-Banach algebras; a least field-valued
norm, trace-class determinant or spectral measure; a solution of a named
published conjecture; priority; or a Lean formalization.

\subsection{Part IV}\label{ihs:app:ledgerD}
The main spectrum theorem uses ordinary compact self-adjoint spectral
theory, finite-word support, Part~I's range-defect theorem and finite
Hermitian diagonalization over the divisible field. The obstruction uses
only the simple-eigenvalue instance and one commutator equation; the trace
and product failures add the explicit second-order recurrence and finite
summation identities. The determinant theorem uses the ordinary trace-class
determinant, Taylor lifting, the ordinary index-zero Fredholm alternative and
finite-dimensional linear algebra. It allows arbitrary complex Hilbert $H$
and set-sized ordered abelian $\Gamma$, including zero; no scalar root
extraction, compact residue or self-adjointness is used. The accumulation
obstruction retains the spectral hypotheses and uses a
common coefficient neighbourhood and a valuation comparison. The steps most
worth auditing independently are the self-adjointness of the coefficientwise
Riesz projection (\cref{ihs:fr:prop:riesz}), the well-ordered induction in
\cref{ihs:fr:thm:nounitary}, and the kernel--cokernel identification in
\cref{ihs:fr:thm:fredholm}.

The Part IV results do not assert a spectral classification of arbitrary
elements of $\Ahh$ or of compact residues with nontrivial kernel; the main
cluster classification without divisibility; a positive spectral measure or general
non-Archimedean Hilbert-space spectral theorem; equivalence with intrinsic
nuclear or compactoid ideals; a determinant for negative-valuation
trace-class series; the impossibility of every regularized spectral product;
anything outside the named coherent analytic category; a failure of global
diagonalization outside bounded coefficients; a solution of a named published
conjecture; priority; or a Lean formalization.

\subsection{What this report as a whole does not claim}\label{ihs:app:jointnonclaims}
Nothing above unifies the four parts into a single theorem, and no hypothesis
is transferred between them. Parts I--III allow nondivisible $\Gamma$ by
their respective proofs; Part IV's spectral classification assumes
divisibility, while its determinant section independently allows arbitrary
ordered abelian groups. The row- and
column-finiteness of Part~II is not available in Part~I or Part~IV. The
spectra of \cref{ihs:tab:senses} are not proved to agree, and no claim is
made that one determines another, beyond the identities of definitions in
\cref{ihs:prop:unify} and the equalities proved for particular operators
(\cref{ihs:hh:thm:spectrum,ihs:fr:thm:main,ihs:fr:thm:fredholm},
\cref{ihs:dz:sec:volterra,ihs:dz:sec:blocks}).
Parts I--III are not trace-class theories; Part IV supplies a coefficientwise
trace and a Fredholm determinant on the nonnegative-valuation component of
$\SH((t^\Gamma))$, and neither a general trace-class nor a nuclear-operator
theory. Part II supplies a vectorwise Hahn-additive projection calculus, not
a general topologically countably additive spectral-measure theorem; the
counterexample in \cref{ihs:rf:rem:notopsum} states the precise limitation.
No part supplies a positive spectral measure. The local cluster theory of
Part IV does not assemble into a global bounded-coefficient spectral
theorem (\cref{ihs:fr:thm:nounitary}).

\clearpage
\phantomsection
\addcontentsline{toc}{section}{References}
\begin{thebibliography}{99}
\raggedright
\bibitem{RepoMap}
V. Reshetnikov, \emph{Surreal repository: documentation catalogue},
\texttt{docs/README.md}, revision \texttt{\commitA}, inspected September 22, 2026.
\url{https://github.com/VladimirReshetnikov/Surreal/blob/608dd23c0539fce73723843949ee627641c27b30/docs/README.md}.

\bibitem{RepoSpectral}
\emph{Surcomplex Spectral Theory: Singular Values, Infinitesimal Rank, and
Multiscale Stability}, AI-assisted research exposition in the Surreal
repository, \texttt{docs/surcomplex/spectral-theory/article.tex}, revision
\texttt{\commitA}, September 2026. Cited above for its closing scope
paragraph, for its example \texttt{spec:ex:monoid}, and for its support
engine \texttt{spec:lem:words} and \texttt{spec:lem:neumann}.
\url{https://github.com/VladimirReshetnikov/Surreal/blob/608dd23c0539fce73723843949ee627641c27b30/docs/surcomplex/spectral-theory/article.tex}.

\bibitem{RepoPinB}
V. Reshetnikov, \emph{Surreal} repository, research documentation,
inspected September 22, 2026 at commit \texttt{\commitB}. Pinned sources:
the spectral-theory README and the documentation catalogue at that revision.
The repository identifies its research reports as AI-assisted drafts.
\url{https://github.com/VladimirReshetnikov/Surreal/blob/a3124af79f66b8b9c196d76b4cbc5ac3938907c4/docs/surcomplex/spectral-theory/README.md}.

\bibitem{RepoPinC}
V. Reshetnikov, \emph{Surreal} repository at commit \texttt{\commitC},
consulted September 22, 2026 by the Part III and Part IV sources; in
particular this report's README and source as they then stood (Parts I and
II only).
\url{https://github.com/VladimirReshetnikov/Surreal/tree/048b72cf7cbfc8ab246e4f73788c10460cb3f6e0/docs/surcomplex/infinite-dimensional-hahn-spectral-theory}.

\bibitem{Gonshor}
H. Gonshor, \emph{An Introduction to the Theory of Surreal Numbers},
London Mathematical Society Lecture Note Series 110, Cambridge University
Press, 1986.
\url{https://www.cambridge.org/core/books/an-introduction-to-the-theory-of-surreal-numbers/312AE504A3E88E804054BFB390446374}.

\bibitem{Neumann}
B. H. Neumann, On ordered division rings,
\emph{Transactions of the American Mathematical Society} \textbf{66}
(1949), 202--252.
\url{https://doi.org/10.1090/S0002-9947-1949-0032593-5}.

\bibitem{Higman}
G. Higman, Ordering by divisibility in abstract algebras,
\emph{Proceedings of the London Mathematical Society}, series 3,
\textbf{2} (1952), 326--336.
\url{https://doi.org/10.1112/plms/s3-2.1.326}.

\bibitem{Douglas}
R. G. Douglas, \emph{Banach Algebra Techniques in Operator Theory},
second edition, Graduate Texts in Mathematics 179, Springer, 1998.
\url{https://doi.org/10.1007/978-1-4612-1656-8}.

\bibitem{Williams}
D. P. Williams, \emph{Lecture Notes on the Spectral Theorem},
Definition~5.1 and footnote~18, pp.~15--16.
Cited only for the classical strong-operator additivity convention.
\url{https://www.math.dartmouth.edu/~dana/bookspapers/ln-spec-thm.pdf}.

\bibitem{AN}
J. Aguayo and M. Nova, Non-Archimedean Hilbert like spaces,
\emph{Bulletin of the Belgian Mathematical Society -- Simon Stevin}
\textbf{14} (2007), no.~5, 787--797.
\url{https://doi.org/10.36045/bbms/1197908895}.

\bibitem{ANS}
J. Aguayo, M. Nova and K. Shamseddine,
Characterization of compact and self-adjoint operators on free Banach
spaces of countable type over the complex Levi-Civita field,
\emph{Journal of Mathematical Physics} \textbf{54} (2013), 023503.
\url{https://doi.org/10.1063/1.4789541}.

\bibitem{Ishizuka}
K. Ishizuka, A spectral theorem for a non-Archimedean valued field whose
residue field is formally real,
\emph{Extracta Mathematicae} \textbf{39} (2024), no.~2, 235--254;
publisher online date January 3, 2025.
\url{https://doi.org/10.17398/2605-5686.39.2.235}.

\bibitem{Boasso}
E. Boasso, The Drazin spectrum in Banach algebras,
in \emph{An Operator Theory Summer: Timisoara, June 29--July 4, 2010},
International Book Series of Mathematical Texts, Theta Foundation,
Bucharest, 2012, 21--28; arXiv:1309.5025, 2013.
\url{https://arxiv.org/abs/1309.5025}.

\bibitem{Miabey}
T. Miabey, Spectral and Essential Spectral Analysis of Finite-Rank
Perturbations of Unbounded Diagonal Operators on Non-Archimedean Hilbert
Spaces, arXiv:2606.06783v1, June 5, 2026, preprint.
\url{https://arxiv.org/abs/2606.06783}.

\bibitem{glo}
A. Greenbaum, R.-c. Li and M. L. Overton,
First-order perturbation theory for eigenvalues and eigenvectors,
arXiv:1903.00785 (2019).
\url{https://arxiv.org/abs/1903.00785}.

\bibitem{ksz}
M. Kenmoe, M. Smerlak and A. Zadorin,
Dynamical perturbation theory for eigenvalue problems,
arXiv:2002.12872 (2020).
\url{https://arxiv.org/abs/2002.12872}.

\bibitem{pathsum}
P.-L. Giscard, S. J. Thwaite and D. Jaksch,
Evaluating matrix functions by resummations on graphs:
the method of path-sums, arXiv:1112.1588 (2011).
\url{https://arxiv.org/abs/1112.1588}.

\bibitem{ago}
P. Ara, K. Goodearl and K. C. O'Meara,
Positive definite matrices and involutions: the manners of their infinite
cousins, arXiv:2607.25134v1, July 27, 2026.
\url{https://arxiv.org/html/2607.25134v1}.
This reference concerns ordinary row- and column-finite matrices; it is
not a source for the Hahn diagonalization asserted in Part~II.

\bibitem{Drazin}
M. P. Drazin, Pseudo-inverses in associative rings and semigroups,
\emph{The American Mathematical Monthly} \textbf{65} (1958), no.~7,
506--514.

\bibitem{BoassoOperators}
E. Boasso, Drazin spectra of Banach space operators and Banach algebra
elements, arXiv:1307.6942, 2013.
\url{https://arxiv.org/abs/1307.6942}.

\bibitem{Koliha}
J. J. Koliha, A generalized Drazin inverse,
\emph{Glasgow Mathematical Journal} \textbf{38} (1996), no.~3,
367--381. \url{https://doi.org/10.1017/S0017089500031803}.

\bibitem{DriesEhrlich}
L. van den Dries and P. Ehrlich, Fields of surreal numbers and
exponentiation, \emph{Fundamenta Mathematicae} \textbf{167} (2001),
173--188. \url{https://doi.org/10.4064/fm167-2-3}.

\bibitem{KaplanKrappSerra}
E. Kaplan, L. S. Krapp and M. Serra, Decomposing the automorphism group
of the surreal numbers, arXiv:2509.22374v3, April 23, 2026.
\url{https://arxiv.org/abs/2509.22374v3}.

\bibitem{Conway}
J. H. Conway, \emph{On Numbers and Games}, second edition,
A K Peters, 2001. First edition, Academic Press, 1976.

\bibitem{Poonen}
B. Poonen, Maximally complete fields,
\emph{L'Enseignement Math\'ematique} (2) \textbf{39} (1993), 87--106.
\url{https://math.stanford.edu/~conrad/Perfseminar/refs/poonencomplete.pdf}.

\bibitem{Kato}
T. Kato, \emph{Perturbation Theory for Linear Operators},
second edition, Springer, 1976.
\url{https://doi.org/10.1007/978-3-642-66282-9}.

\bibitem{SimonProjections}
B. Simon, Unitaries permuting two orthogonal projections,
\emph{Linear Algebra and its Applications} \textbf{528} (2017),
436--441. Preprint: arXiv:1703.05437.
\url{https://arxiv.org/abs/1703.05437}.

\bibitem{Bornemann}
F. Bornemann, On the numerical evaluation of Fredholm determinants,
\emph{Mathematics of Computation} \textbf{79} (2010), 871--915.
\url{https://doi.org/10.1090/S0025-5718-09-02280-7}.
Preprint: \url{https://arxiv.org/abs/0804.2543}.

\bibitem{Serre}
J.-P. Serre, Endomorphismes compl\`etement continus des espaces de Banach
$p$-adiques, \emph{Publications Math\'ematiques de l'IH\'ES}
\textbf{12} (1962), 69--85.
\url{https://www.numdam.org/item/PMIHES_1962__12__69_0/}.
\end{thebibliography}
\end{document}
